\documentclass[11pt,oneside,openany,fontset=fandol]{ctexbook}

\usepackage[
  paperwidth=185mm,
  paperheight=260mm,
  top=19mm,
  bottom=20mm,
  inner=19mm,
  outer=17mm,
  headheight=14pt,
  headsep=7mm,
  footskip=11mm
]{geometry}
\usepackage{amsmath,amssymb,mathtools,bm}
\usepackage{booktabs,array,tabularx,longtable}
\usepackage{xcolor}
\usepackage{enumitem}
\usepackage{fancyhdr}
\usepackage[most]{tcolorbox}
\usepackage{hyperref}
\usepackage{microtype}

\definecolor{LectureBlue}{HTML}{174F7A}
\definecolor{NoteBlue}{HTML}{EAF3F8}
\definecolor{WarmGray}{HTML}{F5F3EF}
\definecolor{DeepGray}{HTML}{333333}
\hypersetup{colorlinks=true,linkcolor=LectureBlue,urlcolor=LectureBlue}

\setlength{\parindent}{2em}
\setlength{\parskip}{0.25em}
\setlist{nosep,leftmargin=2.2em}
\renewcommand{\baselinestretch}{1.18}

\ctexset{
  chapter={
    name={第,章},
    number=\arabic{chapter},
    format=\Huge\bfseries\color{LectureBlue},
    titleformat=\Huge\bfseries\color{LectureBlue},
    beforeskip=0pt,
    afterskip=20pt
  },
  section={format=\Large\bfseries\color{LectureBlue}},
  subsection={format=\large\bfseries\color{LectureBlue}},
  subsubsection={format=\normalsize\bfseries\color{DeepGray}}
}

\pagestyle{fancy}
\fancyhf{}
\fancyhead[L]{\small\color{DeepGray}高等统计学中文学习笔记}
\fancyhead[R]{\small\color{DeepGray}\nouppercase{\leftmark}}
\fancyfoot[R]{\small\thepage}
\renewcommand{\headrulewidth}{0.45pt}
\renewcommand{\footrulewidth}{0pt}

\newtcolorbox{definitionbox}[1]{
  breakable,
  enhanced,
  colback=NoteBlue,
  colframe=LectureBlue,
  boxrule=0.7pt,
  arc=1.5mm,
  title={#1},
  fonttitle=\bfseries
}
\newtcolorbox{notebox}[1][学习注记]{
  breakable,
  enhanced,
  colback=WarmGray,
  colframe=DeepGray!55,
  boxrule=0.5pt,
  arc=1.5mm,
  title={#1},
  fonttitle=\bfseries
}
\newtcolorbox{examplebox}[1]{
  breakable,
  enhanced,
  colback=white,
  colframe=LectureBlue!70,
  boxrule=0.6pt,
  borderline west={2.2pt}{0pt}{LectureBlue},
  title={#1},
  fonttitle=\bfseries
}

\newcommand{\ind}{\mathbf 1}
\newcommand{\asto}{\xrightarrow{\mathrm{a.s.}}}
\newcommand{\pto}{\xrightarrow{P}}
\newcommand{\dto}{\xrightarrow{d}}
\newcommand{\R}{\mathbb R}
\newcommand{\E}{\operatorname{E}}
\newcommand{\Var}{\operatorname{Var}}
\newcommand{\Cov}{\operatorname{Cov}}
\newcommand{\argmax}{\operatorname*{arg\,max}}
\newcommand{\argmin}{\operatorname*{arg\,min}}


% Keep running heads short enough for the original 185 x 260 mm page.
\renewcommand{\chaptermark}[1]{\markboth{\ifnum\value{chapter}>0 第\thechapter 章\fi}{}}
\setlength{\emergencystretch}{1em}
\raggedbottom

\newcommand{\NoteSubtitle}{第一至五章与习题解答}
\begin{document}
\hypersetup{pageanchor=false,pdftitle={高等统计学中文学习笔记},pdfauthor={Studyer-Tang},pdfsubject={概率极限、统计模型、估计与检验}}
\begin{titlepage}
  \centering
  \vspace*{30mm}
  {\Huge\bfseries\color{LectureBlue}高等统计学中文学习笔记\par}
  \vspace{8mm}
  {\LARGE \NoteSubtitle\par}
  \vspace{22mm}
  {\large 定义、证明与习题\par}
  \vspace{5mm}
  {\Large\bfseries Studyer-Tang\par}
  \vspace{3mm}
  {\small \url{https://github.com/Studyer-Tang}\par}
  \vfill
  \begin{tcolorbox}[width=.84\textwidth,colback=WarmGray,colframe=LectureBlue!50,arc=1.5mm]
    \small
    依据 Fang Yao 的 Advanced Mathematical Statistics 课程材料及相关学习推导整理。感谢原课程作者提供讲义与课程框架。本笔记为个人学习记录，非官方讲义或官方译本；原课件与作业未随笔记发布，相关权利归原作者所有。具体范围与校订记录见仓库 SOURCES.md 与 ERRATA.md。
  \end{tcolorbox}
  \vspace{10mm}
  {\small 2026 年 10 月 24 日修订\par}
\end{titlepage}
\frontmatter
\hypersetup{pageanchor=true}
\begingroup
\small
\renewcommand{\baselinestretch}{1.0}\selectfont
\tableofcontents
\endgroup
\mainmatter

\chapter*{阅读说明}
\addcontentsline{toc}{chapter}{阅读说明}
本册将课件内容与学习中反复核对的推导合并编排。正文按数学主题组织，不逐页复制课件。第一、二章分别建立渐近工具与模型结构，第三章讨论风险和信息，第四章转向无偏估计与 U 统计量，第五章研究最优检验及置信集合。

定义和定理用蓝色框，例题用白底侧线框，条件与辨析用灰底框。原有第一、二章的编号继续保留；新增章节按本文结构编号，不与原 Slides 页码逐项对应。概率中的“几乎必然”一般按每个模型分布分别解释；若证明需要共同版本或共同零集，则在相应位置单列说明。

附录 A 整理第一次作业实际讨论的第 7--12 题，附录 B 整理第二次作业全部 11 题。重复讨论的题目合并成一份解答，较短证明保留足够的计算步骤，较长证明把使用条件写在相应等式之前。课件已整理为笔记并不表示已经独立完成全部证明；复习时宜先遮住答案，核对条件、推理和退化情形。

第五章完整收录至置信集合。若按原课程 Slides 的 UMP 范围复习，对应第 1--27 页；UMPU 从第 28 页开始，属于本册保留的后续材料。该范围仅记录已有课程讨论，不能替代后续考试通知。

\chapter{从样本到经验分布}

数理统计以概率规律以及随机样本如何代表这些规律为起点。本章先复习描述分布、密度、条件期望和积分所需的测度论语言，再引入经验分布函数和经验测度，最后利用大数定律、中心极限定理、连续映射定理、Slutsky 定理和 Delta 方法建立基本渐近工具。

本章各部分的关系是
\[
  \text{总体分布 }P
  \longrightarrow
  \text{随机样本 }X_1,\ldots,X_n
  \longrightarrow
  \text{经验分布 }P_n,F_n
  \longrightarrow
  \text{大样本极限理论}.
\]

\section{概率论回顾}

\subsection{测度论的基本语言}

\begin{definitionbox}{定义 1.1.1：$\sigma$-域}
设 $\mathcal F$ 是样本空间 $\Omega$ 的若干子集组成的非空集合族。若
\[
 A\in\mathcal F\Rightarrow A^c\in\mathcal F,
 \qquad
 A_i\in\mathcal F\Rightarrow \bigcup_{i=1}^{\infty}A_i\in\mathcal F,
\]
则称 $\mathcal F$ 为一个 $\sigma$-域。由 De Morgan 律，它也对可列交封闭。实数轴上的 Borel $\sigma$-域 $\mathcal B(\R)$ 是包含所有开集的最小 $\sigma$-域。
\end{definitionbox}

概率模型写作三元组
\[
  (\Omega,\mathcal F,P),
\]
其中 $\Omega$ 给出全部基本结果，$\mathcal F$ 规定哪些集合是可测事件，$P$ 则为事件分配概率。

\begin{definitionbox}{定义 1.1.2：测度与概率}
测度 $\mu$ 是定义在 $(\Omega,\mathcal F)$ 上的非负、可列可加集合函数，且 $\mu(\varnothing)=0$。若 $A_i$ 两两不交，则
\[
 \mu\!\left(\bigcup_{i=1}^{\infty}A_i\right)
 =\sum_{i=1}^{\infty}\mu(A_i).
\]
若再有 $\mu(\Omega)=1$，则 $\mu$ 是概率测度。只在一个 $P$-零测集上可能失败的命题称为几乎必然成立，记为 a.s.。
\end{definitionbox}

连续空间中最常见的参考测度是 Lebesgue 测度 $\lambda$，满足
\[
 \lambda((a,b))=b-a.
\]
可数空间中最常见的是计数测度
\[
 \#(A)=|A|.
\]

若存在 $A_1,A_2,\ldots\in\mathcal F$，使得
\[
 \Omega=\bigcup_{n=1}^{\infty}A_n,
 \qquad \mu(A_n)<\infty,
\]
则称 $\mu$ 为 $\sigma$-有限测度。Lebesgue 测度在 $\R$ 上是 $\sigma$-有限的，因为 $\R=\bigcup_n[-n,n]$；而 $\R$ 上的计数测度不是 $\sigma$-有限的。

\begin{definitionbox}{定义 1.1.3：离散概率测度}
若概率测度 $P$ 集中在有限或可数集合 $\{x_1,x_2,\ldots\}$ 上，且
\[
 P(E)=\sum_i p_i\ind_{\{x_i\in E\}},
 \qquad p_i\geq0,
 \qquad \sum_i p_i=1,
\]
则称 $P$ 为离散概率测度。
\end{definitionbox}

\begin{definitionbox}{定义 1.1.4：绝对连续}
若对每个 $E\in\mathcal F$ 都有
\[
 \nu(E)=0\Rightarrow\mu(E)=0,
\]
则称 $\mu$ 关于 $\nu$ 绝对连续，记为 $\mu\ll\nu$，也称 $\mu$ 被 $\nu$ 控制。
\end{definitionbox}

\begin{notebox}[绝对连续的直觉]
$\nu$ 可以看作参考尺子。$\mu\ll\nu$ 并不表示 $\mu(E)\leq\nu(E)$，它只比较零测集：$\nu$ 完全看不见的集合，$\mu$ 不能在上面放置正质量。

例如，若 $P(X=0)=0.4$、$P(X=1)=0.6$，则
\[
 \lambda(\{0\})=0,\qquad P_X(\{0\})=0.4,
\]
所以 $P_X\not\ll\lambda$。但在其可数支撑上，$P_X\ll\#$。
\end{notebox}

\begin{definitionbox}{定义 1.1.5：随机变量与概率律}
随机变量是可测映射 $X:\Omega\to\R$，即对所有 $a\in\R$，
\[
 X^{-1}(( -\infty,a])=\{\omega:X(\omega)\leq a\}\in\mathcal F.
\]
它的概率律是 $\mathcal B(\R)$ 上的概率测度
\[
 \boxed{P_X(B)=P\bigl(X^{-1}(B)\bigr)=P(X\in B)}.
\]
\end{definitionbox}

这里 $P$ 只能直接测量 $\Omega$ 中的事件，而 $B$ 位于数值空间 $\R$。原像
\[
 X^{-1}(B)=\{\omega:X(\omega)\in B\}
\]
把 $B$ 拉回到原样本空间，再由 $P$ 计算概率。注意 $P_X$ 是概率测度，$P_X(B)$ 才是一个具体概率值。

\subsection{Radon--Nikodym 导数与概率密度}

\begin{definitionbox}{定理 1.1.1：Radon--Nikodym 定理}
设 $\mu,\nu$ 是 $(\Omega,\mathcal F)$ 上的 $\sigma$-有限测度。若 $\mu\ll\nu$，则存在非负可测函数 $f$，使得
\[
 \mu(E)=\int_E f\,\mathrm d\nu,
 \qquad E\in\mathcal F.
\]
函数 $f$ 称为 $\mu$ 关于 $\nu$ 的 Radon--Nikodym 导数，记为
\[
 f=\frac{\mathrm d\mu}{\mathrm d\nu}.
\]
它在 $\nu$-几乎处处意义下唯一。
\end{definitionbox}

RN 导数表示“每单位 $\nu$-测度含有多少 $\mu$-质量”。核心恒等式是
\[
 \boxed{
 \mu(E)=\int_E\frac{\mathrm d\mu}{\mathrm d\nu}\,\mathrm d\nu.}
\]

在离散空间中，以计数测度为参考，
\[
 \frac{\mathrm dP}{\mathrm d\#}(x)=P(\{x\}),
\]
就是 PMF。在连续空间中，若 $P_X\ll\lambda$，则
\[
 f_X=\frac{\mathrm dP_X}{\mathrm d\lambda}
\]
就是通常的 PDF，并且
\[
 P(X\in A)=\int_A f_X(x)\,\mathrm dx,
 \qquad
 F_X(x)=\int_{-\infty}^{x}f_X(y)\,\mathrm dy.
\]

\begin{notebox}[密度依赖参考测度]
密度必须连同参考测度理解。对同一个概率分布，若不同参考测度均允许相应的 RN 表示，所得密度可能不同；并不是任意分布对计数测度或 Lebesgue 测度都有密度。密度值也可以超过 $1$；只有对集合积分以后才得到概率。
\end{notebox}

若 $\mathcal P$ 是概率测度族，并存在同一个 $\sigma$-有限测度 $\nu$ 使每个 $P\in\mathcal P$ 都满足 $P\ll\nu$，则称 $\mathcal P$ 被 $\nu$ 控制，并可由密度族
\[
 \left\{\frac{\mathrm dP}{\mathrm d\nu}:P\in\mathcal P\right\}
\]
表示。

RN 导数具有换测度、可加性和链式法则：
\[
 \int g\,\mathrm d\mu
 =\int g\frac{\mathrm d\mu}{\mathrm d\nu}\,\mathrm d\nu,
\]
\[
 \frac{\mathrm d(\mu_1+\mu_2)}{\mathrm d\nu}
 =\frac{\mathrm d\mu_1}{\mathrm d\nu}
 +\frac{\mathrm d\mu_2}{\mathrm d\nu},
\]
以及当 $\tau\ll\mu\ll\nu$ 时，
\[
 \frac{\mathrm d\tau}{\mathrm d\nu}
 =\frac{\mathrm d\tau}{\mathrm d\mu}
  \frac{\mathrm d\mu}{\mathrm d\nu}.
\]

\begin{notebox}[连续 CDF 不一定有密度]
Cantor 分布集中在不可数但 Lebesgue 测度为零的 Cantor 集 $C$ 上：
\[
 P_X(C)=1,\qquad \lambda(C)=0.
\]
同时它没有单点质量，因此 CDF 处处连续。这说明 CDF 连续只能排除点质量，却不能排除概率集中在零长度的不可数集合上：
\[
 \text{CDF 连续}\not\Rightarrow P_X\ll\lambda.
\]
\end{notebox}

\subsection{变量变换、条件期望与积分}

设 $f:(\Omega,\mathcal F,\nu)\to(\Lambda,\mathcal G)$ 可测。它诱导推前测度
\[
 (\nu\circ f^{-1})(B)=\nu(f^{-1}(B)).
\]

\begin{definitionbox}{定理 1.1.2：变量变换}
对任意非负可测函数或可积函数 $g$，
\[
 \int_\Omega g\circ f\,\mathrm d\nu
 =\int_\Lambda g\,\mathrm d(\nu\circ f^{-1}).
\]
\end{definitionbox}

当 $\nu=P$、$f=X$ 时，推前测度就是 $P_X=P\circ X^{-1}$，从而
\[
 \E\{h(X)\}
 =\int_\Omega h(X(\omega))\,\mathrm dP(\omega)
 =\int_\R h(x)\,\mathrm dP_X(x).
\]

\begin{definitionbox}{定义 1.1.7：条件期望}
设 $X$ 可积，$\mathcal G\subseteq\mathcal F$。条件期望 $Z=\E(X\mid\mathcal G)$ 是满足以下条件且 a.s. 唯一的随机变量：
\begin{enumerate}
 \item $Z$ 是 $\mathcal G$-可测的；
 \item 对每个 $A\in\mathcal G$，
 \[
  \int_A Z\,\mathrm dP=\int_A X\,\mathrm dP.
 \]
\end{enumerate}
\end{definitionbox}

第二个条件等价于
\[
 \E\{(X-Z)\ind_A\}=0,
 \qquad A\in\mathcal G,
\]
即残差在每个已知信息事件上的平均值都为零。给定随机变量 $Y$ 的条件期望定义为
\[
 \E(X\mid Y)=\E\{X\mid\sigma(Y)\}.
\]
塔式性质与全期望公式分别为
\[
 \mathcal H\subseteq\mathcal G
 \Rightarrow
 \E(X\mid\mathcal H)
 =\E\{\E(X\mid\mathcal G)\mid\mathcal H\},
\]
\[
 \E\{\E(X\mid\mathcal G)\}=\E X.
\]
完全信息下 $\E(X\mid\mathcal F)=X$ a.s.；平凡信息下 $\E(X\mid\mathcal G)=\E X$。

若 $(X,Y)$ 具有联合密度，则当 $f_Y(y)>0$ 时，
\[
 f_{X\mid Y}(x\mid y)
 =\frac{f_{X,Y}(x,y)}{f_Y(y)}
 =\frac{f_{X,Y}(x,y)}{\int f_{X,Y}(u,y)\,\mathrm du}.
\]

\begin{definitionbox}{定理 1.1.3：Fubini--Tonelli 定理}
设 $(\Omega_i,\mathcal F_i,\nu_i)$（$i=1,2$）是两个 $\sigma$-有限测度空间，
$f$ 是关于 $\mathcal F_1\otimes\mathcal F_2$ 可测的函数。
若 $f\geq0$（Tonelli，积分允许取 $+\infty$），或
$\int|f|\,\mathrm d(\nu_1\times\nu_2)<\infty$（Fubini，积分有限），则可以交换积分次序：
\[
\begin{aligned}
 \int f\,\mathrm d(\nu_1\times\nu_2)
 &=\int_{\Omega_2}\!\left[\int_{\Omega_1}f(\omega_1,\omega_2)\,\mathrm d\nu_1\right]\mathrm d\nu_2\\
 &=\int_{\Omega_1}\!\left[\int_{\Omega_2}f(\omega_1,\omega_2)\,\mathrm d\nu_2\right]\mathrm d\nu_1.
\end{aligned}
\]
\end{definitionbox}

这里明确采用两个测度均为 $\sigma$-有限的版本。概率测度与 $\R$ 上的 Lebesgue 测度满足这一条件，
但前面提到的 $\R$ 上计数测度不满足；涉及非 $\sigma$-有限测度时不能直接援用本框中的定理。

基本收敛定理包括：对 $X_n\geq0$，Fatou 引理给出
\[
 \E(\liminf_nX_n)\leq\liminf_n\E X_n;
\]
若 $X_n\asto X$ 且 $|X_n|\leq Y$、$\E Y<\infty$，控制收敛定理给出
\[
 \E X_n\to\E X;
\]
若 $0\leq X_n\uparrow X$ a.s.，单调收敛定理同样给出 $\E X_n\to\E X$，两端允许为无穷大。

\section{总体分布、样本与经验分布函数}

给定概率测度 $P$，通常写 $X\sim P$。对期望存在的可测函数 $f$，记
\[
 \boxed{Pf:=\int f\,\mathrm dP=\E\{f(X)\}.}
 \tag{1.1}
\]
总体是一个一般观测 $X$ 的概率分布 $P$；统计模型是可能总体分布的集合 $\mathcal P$。参数模型通常写为
\[
 \mathcal P=\{P_\theta:\theta\in\Theta\},
 \qquad \Theta\subseteq\R^d,
\]
非参数模型则允许未知参数本身是函数。

\begin{definitionbox}{定义 1.2.1：随机样本}
若
\[
 X_1,\ldots,X_n\overset{\mathrm{i.i.d.}}{\sim}P,
\]
则称其为来自 $P$ 的随机样本。其联合分布是乘积测度 $P^n$。观测前 $X_i$ 和统计量 $\bar X$ 是随机的；观测后得到的数据 $x_i$ 和数值 $\bar x$ 是固定的。
\end{definitionbox}

总体 CDF $F$ 未知。落入 $(-\infty,x]$ 的观测比例定义经验 CDF：
\[
 \boxed{F_n(x)=\frac1n\sum_{i=1}^n\ind_{\{X_i\leq x\}}.}
\]
每个实现样本对应的 $F_n$ 都是阶梯型 CDF。若 $X_{(1)}\leq\cdots\leq X_{(n)}$ 是次序统计量，则
\[
 F_n^{-1}(p)=\inf\{x:F_n(x)\geq p\}
 =X_{(\lceil np\rceil)},
 \qquad 0<p\leq1.
\]

固定 $x$，令 $Y_i(x)=\ind_{\{X_i\leq x\}}$。则
\[
 Y_i(x)\sim\operatorname{Bernoulli}(F(x)),
 \qquad
 nF_n(x)\sim\operatorname{Binomial}(n,F(x)).
 \tag{1.2}
\]

\begin{definitionbox}{定理 1.2.1：经验 CDF 的均值与方差}
对每个固定 $x\in\R$，
\[
 \E F_n(x)=F(x),
 \qquad
 \Var(F_n(x))=\frac{F(x)\{1-F(x)\}}n.
\]
\end{definitionbox}

该结论是点态的，只控制固定 $x$，并不能直接控制
\[
 \sup_x|F_n(x)-F(x)|.
\]
后者需要同时控制整个阈值函数类。

对 $x,y\in\R$，令 $x\wedge y=\min\{x,y\}$，则
\[
 \boxed{
 \Cov(F_n(x),F_n(y))
 =\frac{F(x\wedge y)-F(x)F(y)}n.}
 \tag{1.3}
\]

\section{经验测度与代入原理}

对随机样本 $X_1,\ldots,X_n$，定义经验测度
\[
 \boxed{P_n=\frac1n\sum_{i=1}^n\delta_{X_i}},
\]
其中 $\delta_x$ 是点质量测度：
\[
 \delta_x(A)=\ind_A(x),
 \qquad
 \int f\,\mathrm d\delta_x=f(x).
\]
因此
\[
 P_n(A)=\frac1n\sum_{i=1}^n\ind_A(X_i),
 \qquad
 P_nf=\frac1n\sum_{i=1}^nf(X_i).
\]
经验 CDF 是特殊情形
\[
 F_n(t)=P_n((-\infty,t])=P_n\ind_{(-\infty,t]}.
\]

\begin{definitionbox}{命题 1.3.1：经验积分的均值与方差}
若 $f$ 可积，则
\[
 \E(P_nf)=Pf.
\]
若 $f\in L^2(P)$，则
\[
 \Var(P_nf)=\frac{\Var(f(X))}{n}.
\]
\end{definitionbox}

\subsection{统计泛函与代入原理}

\begin{definitionbox}{定义 1.3.1：统计泛函}
设 $\mathcal Q$ 包含统计模型 $\mathcal P$ 以及所考虑的经验测度。统计泛函是映射
\[
 T:\mathcal Q\to\mathcal D.
\]
$T(P)$ 表示总体分布的某个特征。用 $P_n$ 替换未知 $P$，得到代入估计量
\[
 \boxed{\widehat\theta_n=T(P_n).}
\]
\end{definitionbox}

对固定可积函数 $\psi$，
\[
 T_\psi(P)=P\psi=\int\psi\,\mathrm dP
\]
是关于分布 $P$ 的线性泛函，因为
\[
 T_\psi(aP_1+bP_2)
 =aT_\psi(P_1)+bT_\psi(P_2).
\]
代入 $P_n$ 后，
\[
 T_\psi(P_n)=P_n\psi=\frac1n\sum_{i=1}^n\psi(X_i).
\]

\begin{notebox}[“线性”是对分布而言]
$\psi(x)$ 本身不必对 $x$ 线性。例如 $\psi(x)=x^2$ 虽是非线性函数，但 $P\mapsto P(x^2)$ 仍对 $P$ 线性。方差
\[
 \Var_P(X)=P(x^2)-(Px)^2
\]
则不是线性泛函，因为它包含积分结果的平方。
\end{notebox}

均值与方差的直接代入估计为
\[
 \bar X=P_nx,
 \qquad
 S_n^2=P_n(x^2)-(P_nx)^2
 =\frac1n\sum_{i=1}^n(X_i-\bar X)^2.
\]
分位数泛函及其代入估计为
\[
 T_p(P)=F^{-1}(p),
 \qquad
 T_p(P_n)=F_n^{-1}(p).
\]

\subsection{M-估计与 Z-估计}

设 $m_\theta(x)$ 衡量候选参数 $\theta$ 在观测 $x$ 上的表现。总体准则与经验准则分别是
\[
 M(\theta)=Pm_\theta,
 \qquad
 M_n(\theta)=P_nm_\theta.
\]
M-估计定义为
\[
 \boxed{
 \widehat\theta_n\in\argmax_{\theta\in\Theta}P_nm_\theta.}
\]
若使用损失 $\rho_\theta=-m_\theta$，等价地写为最小化 $P_n\rho_\theta$。

\begin{examplebox}{平方损失给出样本均值}
取 $\rho_\theta(x)=(x-\theta)^2$。由于
\[
 \E(X-\theta)^2=\Var(X)+(\mu-\theta)^2,
\]
总体最优点是 $\theta=\mu$。经验问题
\[
 \argmin_\theta\frac1n\sum_{i=1}^n(X_i-\theta)^2
\]
的一阶条件给出 $\widehat\theta_n=\bar X$。若取绝对值损失 $|x-\theta|$，相应 M-估计量则是样本中位数。
\end{examplebox}

在受支配模型中取
\[
 m_\theta(x)=\log p_\theta(x),
\]
便得到最大似然估计。因此 MLE 是 M-估计的特殊情形。

Z-估计从总体估计方程出发：设 $\psi_\theta:\mathcal X\to\R^s$，总体参数满足
\[
 P\psi_{\theta(P)}=0.
\]
用 $P_n$ 替换 $P$，得到
\[
 \boxed{P_n\psi_{\widehat\theta_n}=0.}
\]
任何解都称为 Z-估计量。

\begin{examplebox}{样本均值也是 Z-估计}
取 $\psi_\theta(x)=x-\theta$。总体方程
\[
 \E(X-\theta)=0
\]
的解为 $\theta=\mu$；经验方程
\[
 \frac1n\sum_{i=1}^n(X_i-\theta)=0
\]
的解为 $\widehat\theta_n=\bar X$。
\end{examplebox}

若 M-准则光滑且最优点位于参数空间内部，其一阶条件通常给出 Z-方程。最大似然中的得分函数
\[
 \dot\ell_\theta(x)=\nabla_\theta\log p_\theta(x)
\]
满足
\[
 0=\nabla_\theta\ell_n(\widehat\theta_n)
 =nP_n\dot\ell_{\widehat\theta_n}.
\]
边界最大值和不可微似然不一定满足这个方程，需要单独处理。

\section{收敛方式与随机阶}

\subsection{四种收敛}

设 $\bm Y_n,\bm Y\in\R^k$。主要收敛概念如下：
\begin{align*}
 \bm Y_n\asto\bm Y
 &\iff P(\lim_n\bm Y_n=\bm Y)=1,\\
 \bm Y_n\xrightarrow{L^r}\bm Y
 &\iff \E\|\bm Y_n-\bm Y\|_r^r\to0,\\
 \bm Y_n\pto\bm Y
 &\iff P(\|\bm Y_n-\bm Y\|>\varepsilon)\to0
       \quad\forall\varepsilon>0,\\
 \bm Y_n\dto\bm Y
 &\iff P_{\bm Y_n}\Rightarrow P_{\bm Y}.
\end{align*}
依分布收敛等价于：对每个有界连续函数 $h$，
\[
 \E h(\bm Y_n)\to\E h(\bm Y).
\]

\begin{definitionbox}{定理 1.4.1：依概率收敛的子序列刻画}
$\bm Y_n\pto\bm Y$ 当且仅当每个子序列 $(\bm Y_{n_j})$ 都存在一个进一步子序列 $(\bm Y_{n_{j_\ell}})$，使
\[
 \bm Y_{n_{j_\ell}}\asto\bm Y.
\]
\end{definitionbox}

\textbf{证明。} 若 $\bm Y_n\pto\bm Y$，对任意子序列可递归选取进一步子序列，使
\[
 P(\|\bm Y_{n_{j_\ell}}-\bm Y\|>2^{-\ell})<2^{-\ell}.
\]
这些概率之和有限，由第一 Borel--Cantelli 引理，以概率一只有有限多个相应事件发生，故进一步子序列几乎必然收敛。

反之，若不依概率收敛，可选 $\varepsilon,\eta>0$ 及子序列，使每项均有
$P(\|\bm Y_{n_j}-\bm Y\|>\varepsilon)\ge\eta$。按假设，它有几乎必然收敛的进一步子序列；相应事件的示性函数趋于零，被常数 $1$ 控制，由控制收敛定理其概率趋零，矛盾。

\begin{notebox}[弱收敛的闭集判别及证明]
欧氏空间中，$Y_n\dto Y$ 等价于对每个闭集 $F$ 有
\[
 \limsup_nP(Y_n\in F)\le P(Y\in F).
\]
这称为 Portmanteau 定理的闭集刻画。

先假定弱收敛。对非空闭集取连续函数
$h_\delta(x)=\max\{1-d(x,F)/\delta,0\}$，则 $\ind_F\le h_\delta\le1$，$h_\delta\downarrow\ind_F$ 当 $\delta\downarrow0$。弱收敛及控制收敛给出
\[
 \limsup_nP(Y_n\in F)\le\E h_\delta(Y)\downarrow P(Y\in F).
\]
空集情形显然。

反之，闭集不等式取补集便得到开集概率的下极限不等式。若 $h$ 连续且 $0\le h\le M$，由分层积分与 Fatou 引理，
\[
 \begin{aligned}
 \liminf_n\E h(Y_n)
 &=\liminf_n\int_0^M P\{h(Y_n)>t\}\,\mathrm dt\\
 &\ge\int_0^M\liminf_nP\{h(Y_n)>t\}\,\mathrm dt
 \ge\E h(Y).
 \end{aligned}
\]
再对 $M-h$ 应用同一结论得上极限不等式。一般有界连续函数先平移到非负，故其期望收敛，得到弱收敛。
\end{notebox}


标准蕴含关系为
\[
 \boxed{
 \bm Y_n\asto\bm Y
 \Rightarrow
 \bm Y_n\pto\bm Y
 \Rightarrow
 \bm Y_n\dto\bm Y.}
\]
若有 $L^r$ 收敛，也能由 Markov 不等式推出依概率收敛。

\subsection{矩不等式}

\begin{definitionbox}{定理 1.4.2：Markov 不等式}
若 $Z\geq0$ 且 $\E Z<\infty$，则对每个 $t>0$，
\[
 P(Z\geq t)\leq\frac{\E Z}{t}.
\]
\end{definitionbox}

对 $Z=(X-\E X)^2$ 应用 Markov 不等式，得到 Chebyshev 不等式
\[
 \boxed{
 P(|X-\E X|\geq t)
 \leq\frac{\Var(X)}{t^2}.}
\]
更一般地，
\[
 P(\|\bm Y_n-\bm Y\|_r>\varepsilon)
 \leq
 \frac{\E\|\bm Y_n-\bm Y\|_r^r}{\varepsilon^r}.
 \tag{1.4}
\]

\subsection{随机阶}

对正确定数列 $a_n$，定义
\[
 Y_n=O_p(a_n)
 \iff \frac{Y_n}{a_n}\text{ 依概率有界},
\]
\[
 Y_n=o_p(a_n)
 \iff \frac{Y_n}{a_n}\pto0.
\]
更具体地，$Y_n=O_p(a_n)$ 表示对任意 $\varepsilon>0$，存在 $C_\varepsilon<\infty$，使
\[
 \sup_{n\geq1}P(|Y_n|>C_\varepsilon a_n)<\varepsilon.
\]

基本运算规则包括
\[
 o_p(a_n)O_p(b_n)=o_p(a_nb_n),
\]
\[
 O_p(a_n)+O_p(b_n)=O_p(a_n+b_n),
\]
\[
 o_p(a_n)+o_p(b_n)=o_p(a_n+b_n).
\]

\begin{examplebox}{样本均值的随机阶}
若 $X_i$ i.i.d.，均值为 $\mu$、方差为 $\sigma^2$，则中心极限定理给出
\[
 \sqrt n(\bar X_n-\mu)\dto N(0,\sigma^2).
\]
因此
\[
 \bar X_n-\mu=O_p(n^{-1/2})=o_p(1),
\]
但一般不是 $o_p(n^{-1/2})$，因为标准化误差并未收敛到零，而是收敛到非退化正态分布。
\end{examplebox}

\section{极限定理与变换}

\subsection{大数定律与中心极限定理}

若 $X_i$ i.i.d.，均值为 $\mu$、方差为 $\sigma^2<\infty$，Chebyshev 不等式给出
\[
 P(|\bar X_n-\mu|\geq\varepsilon)
 \leq\frac{\sigma^2}{n\varepsilon^2}\to0.
 \tag{1.5}
\]

\begin{definitionbox}{定理 1.5.1：强大数定律}
若 $X_i\overset{\mathrm{i.i.d.}}{\sim}P$ 且 $P|f|<\infty$，则
\[
 \boxed{P_nf\asto Pf.}
\]
特别地，对每个固定 $x$，$F_n(x)\asto F(x)$。
\end{definitionbox}

\begin{definitionbox}{定理 1.5.2：中心极限定理}
若 $Pf^2<\infty$，$\sigma_f^2=\Var(f(X))\in(0,\infty)$，则
\[
 \boxed{
 \sqrt n(P_n-P)f
 =\frac1{\sqrt n}\sum_{i=1}^n\{f(X_i)-Pf\}
 \dto N(0,\sigma_f^2).}
\]
\end{definitionbox}

固定 $x$ 并取 $f_x(z)=\ind_{\{z\leq x\}}$，可得
\[
 \sqrt n\{F_n(x)-F(x)\}
 \dto N\!\left(0,F(x)\{1-F(x)\}\right).
\]
大数定律说明误差消失，中心极限定理则描述误差在 $n^{-1/2}$ 尺度上的波动。

\subsection{Lindeberg 中心极限定理}

\begin{definitionbox}{定理 1.5.3：Lindeberg 中心极限定理}
对每个 $n$，设 $X_{n1},\ldots,X_{nk_n}$ 相互独立。令
\[
 Y_{nj}=X_{nj}-\E X_{nj},
 \qquad
 s_n^2=\sum_{j=1}^{k_n}\Var(X_{nj})\in(0,\infty).
\]
若对每个 $\varepsilon>0$，
\[
 \frac1{s_n^2}\sum_{j=1}^{k_n}
 \E\!\left[Y_{nj}^2
 \ind_{\{|Y_{nj}|>\varepsilon s_n\}}\right]
 \to0,
 \tag{1.6}
\]
则
\[
 \boxed{
 \frac1{s_n}\sum_{j=1}^{k_n}Y_{nj}\dto N(0,1).}
 \tag{1.7}
\]
\end{definitionbox}

Lindeberg 条件表示：相对于总方差 $s_n^2$，异常大项贡献的方差比例趋于零。它排除了由一个或少数几个变量支配全部波动的情形。

\begin{examplebox}{异质 Bernoulli 顾客例子}
设第 $n$ 天有 $n$ 位独立顾客，
\[
 X_{nj}\sim\operatorname{Bernoulli}\!\left(\frac{j}{2n}\right),
 \qquad j=1,\ldots,n.
\]
总购买人数 $T_n=\sum_jX_{nj}$ 满足
\[
 \E T_n=\frac{n+1}{4},
 \qquad
 s_n^2=\frac{(n+1)(4n-1)}{24n}\sim\frac n6.
\]
由于 $|X_{nj}-\E X_{nj}|\leq1$ 而 $s_n\to\infty$，对固定 $\varepsilon>0$，事件
\[
 |Y_{nj}|>\varepsilon s_n
\]
最终不可能发生，故 Lindeberg 条件成立。因此
\[
 \frac{T_n-(n+1)/4}{\sqrt{(n+1)(4n-1)/(24n)}}
 \dto N(0,1).
\]
各顾客购买概率不同，但总购买人数仍具有正态极限。
\end{examplebox}

\subsection{连续映射定理与 Slutsky 定理}

\begin{definitionbox}{定理 1.5.4：连续映射定理}
设 $g$ 为 Borel 可测映射。若 $g$ 在除 $P_{\bm Y}$-零测集外的每一点连续，则
\[
 \bm Y_n\asto\bm Y\Rightarrow g(\bm Y_n)\asto g(\bm Y),
\]
\[
 \bm Y_n\pto\bm Y\Rightarrow g(\bm Y_n)\pto g(\bm Y),
\]
\[
 \bm Y_n\dto\bm Y\Rightarrow g(\bm Y_n)\dto g(\bm Y).
\]
\end{definitionbox}

\textbf{证明。} 设 $g$ 可测，$D_g$ 为其不连续点集，$P(\bm Y\in D_g)=0$。几乎必然收敛时，在 $\bm Y_n(\omega)\to\bm Y(\omega)$ 且 $\bm Y(\omega)\notin D_g$ 的概率一事件上应用连续性定义即可。

依概率收敛时，任取子序列。由子序列刻画，可选进一步子序列几乎必然收敛到 $\bm Y$，其像也几乎必然收敛。再用一次子序列刻画，得 $g(\bm Y_n)\pto g(\bm Y)$。

依分布收敛时，取像空间中的闭集 $F$。若 $x\in\overline{g^{-1}(F)}$ 且 $x\notin D_g$，可取 $x_j\in g^{-1}(F)$ 收敛到 $x$；由连续性与 $F$ 的闭性，$g(x)\in F$。故
\[
 \overline{g^{-1}(F)}\subseteq g^{-1}(F)\cup D_g.
\]
由 Portmanteau 定理的闭集不等式，
\[
\begin{aligned}
 \limsup_n P\{g(\bm Y_n)\in F\}
 &\le\limsup_n P\{\bm Y_n\in\overline{g^{-1}(F)}\}\\
 &\le P\{\bm Y\in\overline{g^{-1}(F)}\}
 \le P\{g(\bm Y)\in F\}.
\end{aligned}
\]
再次使用 Portmanteau 定理即得结论，不要求 $g$ 处处连续。

例如，若 $\bar X_n\pto\mu$ 且 $g$ 在 $\mu$ 处连续，则
\[
 g(\bar X_n)\pto g(\mu).
\]

\begin{definitionbox}{命题 1.5.1：渐近等价}
若
\[
 \bm U_n\dto\bm U,
 \qquad
 \|\bm V_n-\bm U_n\|\pto0,
\]
则 $\bm V_n\dto\bm U$。
\end{definitionbox}

\textbf{渐近等价的证明。} 对 $h\in C_c(\R^k)$，令 $M=\sup|h|$。紧支撑连续函数在整个空间一致连续，故给定 $\eta>0$，可选 $\delta>0$ 使距离不超过 $\delta$ 时函数值之差不超过 $\eta$。于是
\[
 |\E h(\bm V_n)-\E h(\bm U_n)|
 \le\eta+2M P(\|\bm V_n-\bm U_n\|>\delta).
\]
先令 $n\to\infty$，再令 $\eta\downarrow0$，左侧趋于零。由 $\bm U_n\dto\bm U$，得到 $\E h(\bm V_n)\to\E h(\bm U)$。在欧氏空间，对所有 $C_c$ 函数的积分收敛到一个概率律即可推出弱收敛：用紧支撑连续函数 $0\le\chi_R\le1$ 截断，先控制尾部概率，再把有界连续函数写成 $h\chi_R+h(1-\chi_R)$。因此 $\bm V_n\dto\bm U$。一般有界连续函数未必一致连续，不能直接套用上述 $\delta$。

\begin{definitionbox}{定理 1.5.5：Slutsky 定理}
若 $Y_n\dto Y$、$Z_n\pto c$，其中 $c$ 为常数，则
\[
 (Y_n,Z_n)\dto(Y,c).
\]
因此
\[
 Y_n+Z_n\dto Y+c,
 \qquad
 Z_nY_n\dto cY,
\]
且当 $c\neq0$ 时，
\[
 \frac{Y_n}{Z_n}\dto\frac{Y}{c}.
\]
\end{definitionbox}

\textbf{Slutsky 定理的证明。} 连续映射 $y\mapsto(y,c)$ 给出 $(Y_n,c)\dto(Y,c)$，且
\[
 \|(Y_n,Z_n)-(Y_n,c)\|=\|Z_n-c\|\pto0.
\]
由渐近等价命题得到联合收敛，再对加法、乘法、除法应用连续映射定理。除法在分母非零处连续，故需 $c\ne0$；若某些 $n$ 的分母为零，可在该事件任意补定义，其概率趋于零。

\begin{examplebox}{学生化样本均值}
设样本独立同分布，且有限方差 $\sigma^2>0$。令
\[
 S_n^2=\frac1n\sum_{i=1}^n(X_i-\bar X_n)^2.
\]
中心极限定理给出
\[
 \frac{\sqrt n(\bar X_n-\mu)}\sigma\dto N(0,1),
\]
而大数定律与连续映射定理给出 $S_n/\sigma\pto1$。由 Slutsky 定理，
\[
 \boxed{
 \frac{\sqrt n(\bar X_n-\mu)}{S_n}
 \dto N(0,1).}
\]
相合的尺度估计量可以替代未知总体尺度，而不改变一阶极限分布。
\end{examplebox}

\subsection{Delta 方法}

连续映射定理可以处理相合性，却不能描述非线性变换后的一阶波动。Delta 方法将随机变量的 Taylor 展开与已有渐近分布结合起来。

\begin{definitionbox}{命题 1.5.2：随机 Taylor 展开}
若 $\bm X_n\pto\bm c$，且当 $\bm h\to0$ 时
\[
 g(\bm c+\bm h)
 =g(\bm c)+\sum_{j=1}^m\frac1{j!}
 D^jg(\bm c)[\bm h,\ldots,\bm h]+R_m(\bm h),
\]
其中 $R_m(\bm h)/\|\bm h\|^m\to0$，则
\[
\begin{aligned}
 g(\bm X_n)=g(\bm c)
 &+\sum_{j=1}^m\frac1{j!}
 D^jg(\bm c)[\bm X_n-\bm c,\ldots,\bm X_n-\bm c]\\
 &+\|\bm X_n-\bm c\|^mQ_n,
 \qquad Q_n=o_p(1).
\end{aligned}
\]
\end{definitionbox}

\textbf{随机 Taylor 展开的证明。} 令 $q(h)=R_m(h)/\|h\|^m$（$h\ne0$），并令 $q(0)=0$。由 Peano 余项假设，$q(h)\to0$ 当 $h\to0$。取 $Q_n=q(\bm X_n-\bm c)$。给定 $\varepsilon>0$，可选 $\delta>0$ 使 $\|h\|<\delta$ 时 $|q(h)|\le\varepsilon$，于是
\[
 P(|Q_n|>\varepsilon)\le P(\|\bm X_n-\bm c\|\ge\delta)\to0.
\]
故 $Q_n=o_p(1)$，代回确定性展开即可；这不要求每条样本路径都收敛。

\begin{definitionbox}{定理 1.5.6：一阶 Delta 方法}
若 $a_n\to\infty$，
\[
 a_n(\bm X_n-\bm c)\dto\bm Y,
\]
且 $g:\R^k\to\R$ 在 $\bm c$ 处可微，则
\[
 \boxed{
 a_n\{g(\bm X_n)-g(\bm c)\}
 \dto \nabla g(\bm c)^{\mathsf T}\bm Y.}
\]
若 $\bm Y\sim N_k(0,\bm\Sigma)$，则极限为
\[
 N\!\left(0,
 \nabla g(\bm c)^{\mathsf T}
 \bm\Sigma
 \nabla g(\bm c)\right).
\]
\end{definitionbox}

\textbf{证明。} 弱收敛序列依概率有界，所以 $H_n=a_n(\bm X_n-\bm c)=O_p(1)$。因 $a_n\to\infty$，$\bm X_n\pto\bm c$。随机 Taylor 展开给出
\[
 a_n\{g(\bm X_n)-g(\bm c)\}
 =\nabla g(\bm c)^{\mathsf T}H_n+\|H_n\|Q_n.
\]
第一项由连续映射定理依分布收敛，第二项为 $O_p(1)o_p(1)=o_p(1)$。这里乘积规则来自
\[
 P(|A_nB_n|>\varepsilon)\le P(|A_n|>K)+P(|B_n|>\varepsilon/K):
\]
先选 $K$ 控制第一项，再令 $n\to\infty$ 控制第二项。由 Slutsky 定理得证。

\begin{definitionbox}{高阶 Delta 方法}
若 $a_n(\bm X_n-\bm c)\dto\bm Y$，$a_n\to\infty$，$g$ 在 $\bm c$ 有 $m$ 阶 Peano 展开，且 $D^jg(\bm c)=0$ 对 $1\le j<m$ 成立，则
\[
 a_n^m\{g(\bm X_n)-g(\bm c)\}
 \dto\frac1{m!}D^mg(\bm c)[\bm Y,\ldots,\bm Y].
\]
一维时右端为 $g^{(m)}(c)Y^m/m!$。
\end{definitionbox}

\textbf{证明。} 令 $H_n=a_n(\bm X_n-\bm c)$，由低阶导数消失，
\[
 a_n^m\{g(\bm X_n)-g(\bm c)\}
 =\frac1{m!}D^mg(\bm c)[H_n,\ldots,H_n]+\|H_n\|^mQ_n.
\]
$m$ 重线性映射在对角线上是连续齐次多项式，第一项由连续映射定理收敛；$H_n=O_p(1)$ 使余项为 $o_p(1)$。再用 Slutsky 定理。即使 $D^mg(\bm c)\ne0$，$\bm Y$ 也可能落在该多项式的零集上，使极限退化。

\begin{examplebox}{样本均值的平方}
由
\[
 \sqrt n(\bar X_n-\mu)\dto N(0,\sigma^2),
\]
对 $g(t)=t^2$，当 $\mu\neq0$ 时，$g'(\mu)=2\mu$，故
\[
 \sqrt n(\bar X_n^2-\mu^2)
 \dto N(0,4\mu^2\sigma^2).
\]
当 $\mu=0$ 时一阶导数消失，需要二阶 Delta 方法：
\[
 n\bar X_n^2\dto\sigma^2Z^2,
 \qquad Z\sim N(0,1).
\]
\end{examplebox}

\begin{examplebox}{样本方差的渐近分布}
若 $P(x^4)<\infty$，令
\[
 S_n^2=P_n(x^2)-(P_nx)^2.
\]
对 $\bm W_i=(X_i,X_i^2)^{\mathsf T}$ 使用二维 CLT，再对 $h(a,b)=b-a^2$ 使用 Delta 方法。梯度为
\[
 \nabla h(\mu,Px^2)=(-2\mu,1)^{\mathsf T}.
\]
记 $\mu=\E X$、$m_j=\E X^j$、$\mu_4=\E(X-\mu)^4$。二维 CLT 的协方差矩阵为
\[
 \bm\Sigma=\begin{pmatrix}
 m_2-\mu^2&m_3-\mu m_2\\m_3-\mu m_2&m_4-m_2^2
 \end{pmatrix}.
\]
二次型 $(-2\mu,1)\bm\Sigma(-2\mu,1)^{\mathsf T}$ 是线性组合的方差。其中心化变量为
\[
 (X^2-m_2)-2\mu(X-\mu)=(X-\mu)^2-(m_2-\mu^2).
\]
因此中心化线性项化简为
\[
 L(X)=(X-\mu)^2-\sigma^2,
\]
故渐近方差为
\[
 \Var\{L(X)\}=P\{(x-\mu)^4\}-\sigma^4.
\]
最终
\[
 \boxed{
 \sqrt n(S_n^2-\sigma^2)
 \dto N\!\left(0,P\{(x-\mu)^4\}-\sigma^4\right).}
\]
使用分母 $n-1$ 的样本方差记为 $\widetilde S_n^2=\frac{n}{n-1}S_n^2$（$n\geq2$）。
由大数定律与连续映射定理，$S_n^2\pto\sigma^2$，因此 $S_n^2=O_p(1)$。在这里所需的
$\sqrt n$ 中心化尺度下，
\[
 \sqrt n(\widetilde S_n^2-S_n^2)
 =\frac{\sqrt n}{n-1}S_n^2
 \pto0.
\]
于是
\[
 \sqrt n(\widetilde S_n^2-\sigma^2)
 =\sqrt n(S_n^2-\sigma^2)+o_p(1),
\]
由 Slutsky 定理得到相同的极限分布。关键是两者的差在 $\sqrt n$ 尺度下可以忽略；
仅有乘法因子趋于 $1$，一般不足以保证中心化后的极限相同。
\end{examplebox}

\section*{本章总结}
\addcontentsline{toc}{section}{本章总结}

概率分布与期望描述总体量，经验测度则给出对应的样本量。Markov 和 Chebyshev 不等式把矩信息转化为概率上界；大数定律说明估计误差消失；中心极限定理描述 $n^{-1/2}$ 尺度上的随机波动；连续映射、Slutsky 定理和 Delta 方法则把这些极限传递到非线性代入估计量。

第三章的一致性分析、第四章的 U 统计量极限分布都需要这些工具。

\begin{notebox}[复习路线]
建议按以下顺序复习：
\[
 P_nf\to Pf
 \quad\Longrightarrow\quad
 \sqrt n(P_n-P)f\dto N(0,\sigma_f^2)
 \quad\Longrightarrow\quad
 g(P_nf)\text{ 的 Delta 展开}.
\]
首先确认估计量是否相合，再寻找其波动尺度，最后分析非线性变换和学生化。
\end{notebox}

\section*{练习与核对}
\addcontentsline{toc}{section}{练习与核对}
\begin{enumerate}
\item 若 $X_i$ 独立同分布于 $\operatorname{Bernoulli}(p)$，求 $\bar X_n(1-\bar X_n)$ 在 $p\ne1/2$ 与 $p=1/2$ 时适当缩放的极限分布。
\item $A_n\dto A$、$B_n\dto B$ 能否推出 $A_n+B_n\dto A+B$？说明 Slutsky 定理为何要求其中一个量收敛到常数。
\item 在四阶中心矩有限时，直接证明 $S_n^2-\sigma^2=(P_n-P)(x-\mu)^2-(\bar X_n-\mu)^2$，据此重做样本方差的极限分布。
\end{enumerate}
\textbf{答案线索。} (1) $p\ne1/2$ 时，$\sqrt n$ 极限为 $N(0,(1-2p)^2p(1-p))$；$p=1/2$ 时，$n\{\bar X_n(1-\bar X_n)-1/4\}\dto-Z^2/4$，$Z\sim N(0,1)$。(2) 不能：取标准正态 $Z$，$A_n=Z$、$B_n=(-1)^nZ$，各边缘律不变，和却在 $2Z$ 与 $0$ 之间交替。(3) 展开中心化平方，余项的 $\sqrt n$ 倍为 $O_p(n^{-1/2})$。若 $\mu_4=\sigma^4$，正态极限退化为零，不能据此做非退化正态近似。

\chapter{统计模型、指数族、充分性与完备性}

本章依据 Fang Yao 老师《第02章-讲义》整理，原文件共 22 页，对应讲义正文第 27--48 页。定义、定理及例子的编号沿用讲义；“学习注记”和“补充证明”记录个人理解与讨论中的展开，不是讲义逐字译文。Slides 的编号与讲义不同，例如 Slides 引理 2.7 对应讲义引理 2.3.1，Slides 命题 2.9 对应讲义命题 2.4.1。

\begin{notebox}[学习起点：本章到底要做什么？]
第一章主要问：样本量增大时，经验分布和估计量怎样逼近总体？第二章先固定样本量，问：在指定模型下，原始数据中的哪些部分真正参与对参数的比较？

本章依次研究三个问题：怎样用似然比较候选参数；怎样压缩样本而保持充分性；怎样用完备性得到函数唯一性、用 Basu 定理得到独立性。“已整理”不等于已经能够独立证明，实际追问和待复习内容另见本章学习历程。
\end{notebox}

\section{统计模型}

\subsection{从具体样本到模型}
想估计硬币正面概率 $p$，独立抛掷 $n$ 次，把正面记作 $1$。例如
\[
 x=(1,0,1,1,0),\qquad y=(0,1,1,0,1)
\]
都有三次成功。独立性给出
\[
 P_p(\mathbf X=x)=p(1-p)pp(1-p)=p^3(1-p)^2=P_p(\mathbf X=y).
\]
一般地，记 $s=\sum_i x_i$，则
\[
 P_p(\mathbf X=x)
 =\prod_{i=1}^n p^{x_i}(1-p)^{1-x_i}
 =p^s(1-p)^{n-s}.
\]
固定 $n$ 后，这个公式只需要成功次数，不需要成功顺序。后面要把“似乎可以丢弃顺序”变成严格的充分性结论。

\begin{notebox}[具体排列与成功次数不是同一个事件]
某一个排列的概率是 $p^s(1-p)^{n-s}$；成功次数等于 $s$ 的概率要对所有相容排列求和，才有组合数：
\[
 P_p(S=s)=\binom ns p^s(1-p)^{n-s}.
\]
知道 $S$ 通常不能恢复完整样本，但这不必然意味着损失了关于 $p$ 的信息。
\end{notebox}

参数模型是一族共同样本空间上的概率分布：
\[
 \mathcal P=\{P_\theta:\theta\in\Theta\}.
\]
这里 $\Theta$ 是允许的参数集合，$\theta$ 是未知常数，$P_\theta$ 是对应的数据生成规律。i.i.d. 样本的联合模型为
\[
 \mathcal P^{(n)}=\{P_\theta^{\otimes n}:\theta\in\Theta\}.
\]
$\mathbf X=(X_1,\ldots,X_n)$ 表示随机样本，$\mathbf x=(x_1,\ldots,x_n)$ 表示已观测到的实现值。统计量是计算规则不含未知参数的数据函数；若 $\mu$ 未知，$\bar X-\mu$ 不是直接可计算的统计量。

\begin{definitionbox}{定义 2.1.1：可识别性}
模型可识别是指
\[
 P_\theta=P_{\theta'}\ \Longrightarrow\ \theta=\theta'.
\]
等号指两个概率测度对每个事件都给出相同概率，不是某份样本的两个似然值碰巧相同。
\end{definitionbox}

已知方差的正态位置模型可识别：分布相同则期望相同，故均值参数相同。反之，模型 $N(\alpha\beta,1)$ 中，$(\alpha,\beta)=(1,2)$ 和 $(2,1)$ 产生同一分布，只能识别乘积，不能在整个模型中分别识别两者。增加样本量不能消除这种参数化问题。

\subsection{密度与似然：同一个公式的两个视角}
\begin{definitionbox}{定义 2.1.2：受支配分布族}
若存在共同的 $\sigma$-有限测度 $\nu$，使
\[
 P\ll\nu\quad\forall P\in\mathcal P,
\]
则称模型受 $\nu$ 支配。记 $p_\theta=\mathrm dP_\theta/\mathrm d\nu$。
\end{definitionbox}

计数测度对应离散 PMF，Lebesgue 测度对应通常的连续 PDF。密度满足
\[
 P_\theta(A)=\int_Ap_\theta\,\mathrm d\nu.
\]
独立同分布样本的联合密度与似然为
\[
 p_\theta^{(n)}(\mathbf x)=\prod_i p_\theta(x_i),
 \qquad L_{\mathbf x}^{(n)}(\theta)=p_\theta^{(n)}(\mathbf x).
\]
固定 $\theta$，它描述数据的概率规律；固定 $\mathbf x$，它比较候选参数。连续模型使用密度值，不是通常为零的单点概率。

\begin{notebox}[似然不是参数的概率分布]
当 $n=5,s=3$ 时，$L(p)=p^3(1-p)^2$，但
\[
 \int_0^1L(p)\,\mathrm dp
 =\frac14-\frac25+\frac16=\frac1{60}\ne1.
\]
比较 $L(0.6)/L(0.2)=6.75$，只是比较同一数据在两个候选参数下的概率，不是两个参数自身的概率比。后验分布还需要先验和归一化。
\end{notebox}

\subsection{最大似然估计与边界}
\begin{definitionbox}{定义 2.1.3：最大似然估计}
若统计量 $\widehat\theta(\mathbf X)$ 的观测值满足
\[
 \widehat\theta(\mathbf x)\in\argmax_{\theta\in\Theta}L_{\mathbf x}^{(n)}(\theta),
\]
则称为 MLE。$\argmax$ 是最大化参数的集合，不是最大的函数值。
\end{definitionbox}

对数严格递增，故最大化似然等价于最大化扩展对数似然，约定 $\log0=-\infty$：
\[
 \ell_{\mathbf x}^{(n)}(\theta)=\sum_i\log p_\theta(x_i).
\]
这与第一章相接：
\[
 n^{-1}\ell_{\mathbf X}^{(n)}(\theta)=P_n\log p_\theta.
\]
本章先研究有限样本结构，后续一致性理论才讨论这个经验准则的极限。非空紧参数空间上的上半连续似然可以保证最大值取到；严格凹性可以保证已存在的最大化解唯一，但不自动保证存在性。

\begin{examplebox}{例 2.1.2：Bernoulli 的 MLE}
固定二元样本，令 $s=\sum_i x_i$。在 $0<s<n$ 时，
\[
 \ell(p)=s\log p+(n-s)\log(1-p),\qquad
 \ell'(p)=\frac sp-\frac{n-s}{1-p}.
\]
因为 $p(1-p)>0$，
\[
 \ell'(p)=0
 \iff s(1-p)=(n-s)p
 \iff s=np
 \iff p=s/n.
\]
又
\[
 \ell''(p)=-s/p^2-(n-s)/(1-p)^2<0,
\]
所以这是唯一内点最大值。若 $s=0$，$L(p)=(1-p)^n$ 递减；若 $s=n$，$L(p)=p^n$ 递增。在 $[0,1]$ 上分别得到边界 MLE $0,1$，在 $(0,1)$ 上却都没有 MLE，因为上确界只能在被排除的边界逼近。
\end{examplebox}

\subsection{参数无关因子与参数相关支撑}
若固定样本上 $a(\mathbf x)>0$，且
\[
 p_\theta^{(n)}(\mathbf x)=a(\mathbf x)b_\theta(\mathbf x),
\]
其中 $a$ 不含参数，则相同因子在参数之间的似然比中消去，也不改变最大化位置。它仍是概率模型的一部分，计算样本事件概率时不能删除。

Poisson 样本满足
\[
 p_\lambda^{(n)}(\mathbf x)
 =e^{-n\lambda}\lambda^{\sum_i x_i}\prod_i\frac1{x_i!}.
\]
例如 $(0,2)$ 与 $(1,1)$ 的似然分别为
\[
 \tfrac12e^{-2\lambda}\lambda^2,\qquad e^{-2\lambda}\lambda^2.
\]
二者不相等，却始终成固定比例，对任意两个参数的相对比较相同。

给定总和为 $2$，三个可能样本为 $(0,2)$、$(1,1)$、$(2,0)$。它们的条件概率依次为 $1/4$、$1/2$、$1/4$，不等概率但都不含参数。

对 $U(0,\theta)$，联合密度则为
\[
 p_\theta^{(n)}(\mathbf x)
 =\theta^{-n}\ind_{\{x_{(n)}<\theta\}}\ind_{\{x_{(1)}>0\}}.
\]
示性函数含参数，不能忽略。若最大观测是 $3$，则 $\theta=2$ 与数据不相容；仅保留 $\theta^{-n}$ 会错误地给它正似然。

\begin{notebox}[两种似然比较，以及与后文的联系]
固定同一个样本，比较 $L_{\mathbf x}(\theta)/L_{\mathbf x}(\theta')$，用于参数推断；比较两组样本的整条似然是否满足
\[
 L_{\mathbf x}(\theta)=c(\mathbf x,\mathbf y)L_{\mathbf y}(\theta)
 \quad\forall\theta,\qquad c(\mathbf x,\mathbf y)>0
\]
且 $c$ 不含参数，用于判断哪些样本可以合并。

参数依赖支撑时要同时比较零值区域，不能只在公共正值区域约分。按严格端点密度 $\ind_{\{x<\theta\}}$，均匀模型的 MLE 可能只有未取到的上确界；改用含端点密度会改变逐点最大化。计算时应保持同一密度版本。
\end{notebox}

\section{指数族}

\subsection{指数形式与自然参数}
前面发现 Bernoulli 和 Poisson 的参数比较都只需要总和。指数族为这种现象提供统一形式：参数和数据通过一个有限维内积相结合，独立样本相乘后变成自然统计量的求和。

\begin{definitionbox}{定义 2.2.1：自然指数族}
相对于共同 $\sigma$-有限测度 $\nu$，密度为
\[
 p_\eta(x)=\exp\{\eta^\top T(x)-\zeta(\eta)\}h(x),
 \qquad \eta\in\Xi\subseteq\R^k.
\]
这里 $T$ 是自然统计量，$\eta$ 是自然参数，$h\ge0$ 不含参数，而
\[
 \zeta(\eta)=\log\int e^{\eta^\top T(x)}h(x)\,\nu(\mathrm dx)
\]
负责归一化。最大自然参数空间由该积分有限且为正的参数组成；实际子模型可只使用其中一部分。
\end{definitionbox}

积分核对：
\[
 \int p_\eta\,\mathrm d\nu
 =e^{-\zeta(\eta)}\int e^{\eta^\top T(x)}h(x)\,\mathrm d\nu=1.
\]
若用可逆矩阵 $A$ 替换为 $\widetilde T=AT$、$\widetilde\eta=A^{-\top}\eta$，内积不变。这说明表示不唯一，统计意义取决于所生成的信息，而不取决于选用的坐标名称。

\begin{examplebox}{例 2.2.1--2.2.3：从密度中识别自然参数}
Bernoulli 模型中，对 $x\in\{0,1\}$，
\[
 p^x(1-p)^{1-x}
 =\exp\left\{x\log\frac p{1-p}+\log(1-p)\right\}.
\]
取 $\eta=\log[p/(1-p)]$，反解得 $p=e^\eta/(1+e^\eta)$，所以
\[
 T(x)=x,\quad \zeta(\eta)=\log(1+e^\eta),\quad \Xi=\R.
\]
Poisson 模型中，
\[
 e^{-\lambda}\lambda^x/x!
 =e^{x\eta-e^\eta}/x!,\qquad\eta=\log\lambda\in\R.
\]
所以 $T(x)=x$、$h(x)=1/x!$、$\zeta(\eta)=e^\eta$。

正态两参数模型中，
\[
 -\frac{(x-\mu)^2}{2\sigma^2}
 =\frac{\mu}{\sigma^2}x-\frac{x^2}{2\sigma^2}
 -\frac{\mu^2}{2\sigma^2}.
\]
可取
\[
 T(x)=(x,x^2),\quad
 \eta_1=\mu/\sigma^2,\quad \eta_2=-1/(2\sigma^2),\quad h(x)=1,
\]
\[
 \zeta(\eta)=-\frac{\eta_1^2}{4\eta_2}
 +\frac12\log\frac{-\pi}{\eta_2},
 \qquad \Xi=\R\times(-\infty,0).
\]
\end{examplebox}

对于 i.i.d. 样本，
\[
 \begin{aligned}
 p_\eta^{(n)}(\mathbf x)
 &=\prod_i e^{\eta^\top T(x_i)-\zeta(\eta)}h(x_i)\\
 &=\exp\left\{\eta^\top\sum_iT(x_i)-n\zeta(\eta)\right\}
   \prod_i h(x_i).
 \end{aligned}
\]
因此参数相关部分只需要 $T_n=\sum_iT(X_i)$。正态模型对应 $(\sum_iX_i,\sum_iX_i^2)$，第二项是“平方后求和”，不是 $(\bar X)^2$。

\subsection{满秩表示的两个条件}
讲义对满秩的约定同时要求：
\begin{enumerate}
\item 统计量一侧没有仿射冗余：
\[
 a^\top T(X)=b\ \text{几乎必然}\quad\Longrightarrow\quad a=0;
\]
\item 自然参数一侧含非空开集：
\[
 \operatorname{int}_{\R^k}\Xi\ne\varnothing.
\]
\end{enumerate}
第一条的 a.s. 按共同载体测度或其等价概率版本理解；“rank”是指数族表示的性质，不是在计算某个观测向量的矩阵秩。

\begin{notebox}[追问：是不是统计量各分量线性无关？]
还要把常数 $1$ 算进去。$T=(X,2X)$ 满足 $2T_1-T_2=0$，有冗余；$T=(X,X+1)$ 满足 $T_2-T_1=1$，也有冗余。需排除的是右侧为任意常数的非平凡仿射关系，而不只是齐次线性关系。

对具有正态分布的 $X$，$T=(X,X^2)$ 没有仿射冗余：若 $a_1x+a_2x^2=b$ 几乎处处成立，连续多项式必恒等于零，从而 $a_1=a_2=b=0$。其取值虽在抛物线上，却不被一条直线包含。没有仿射冗余既不意味着分量独立，也不排除非线性函数关系。
\end{notebox}

“本质值域”忽略零概率修改。对于欧氏空间中的随机向量，点 $t$ 属于本质值域意味着每个邻域都有正概率：
\[
 P(\|T-t\|<\varepsilon)>0\quad\forall\varepsilon>0.
\]
不要求 $P(T=t)>0$。例如将 $(X,2X)$ 在正态零概率事件 $\{X=0\}$ 上改为 $(0,100)$，仍不能消除几乎必然的仿射冗余。

第二条则要求自然参数集合内至少容得下一个 $k$ 维开球。正态两参数模型的 $\Xi=\R\times(-\infty,0)$ 可容纳以 $(0,-1)$ 为中心、半径 $1/2$ 的圆盘，因此满足第二条。

\subsection{弯曲指数族与参数约束}
通常在光滑、正则的参数化下，弯曲族写为
\[
 \eta=\eta(\theta)\in\R^k,\qquad \theta\in\Theta\subseteq\R^d,\quad d<k.
\]
原始参数 $\theta$ 只有 $d$ 个自由方向，决定了 $k$ 个自然参数坐标，所以自然参数只能沿低维曲面变化。仅凭任意映射 $d<k$ 不能保证其像为光滑曲面，这里使用通常的正则参数化语境。

例如 $X\sim N(\theta,\theta^2)$，$\theta>0$，仍有 $T=(X,X^2)$，但
\[
 \eta_1=1/\theta,\qquad\eta_2=-1/(2\theta^2)
 \quad\Longrightarrow\quad \eta_2=-\eta_1^2/2.
\]
参数在二维平面中沿一条抛物线变化，不能包含二维开球。

\begin{notebox}[两个抛物线，为什么结论不同？]
统计量 $(x,x^2)$ 落在抛物线上，不违反第一条，因为第一条仅排除仿射超平面，二维中即直线。自然参数落在抛物线上，却违反第二条，因为第二条要求真正的二维开区域。

因此，统计量无仿射冗余，并不推出整个表示满秩。缺少开集也只意味着后面的指数族完备性定理不能直接使用，并非自动证明不完备；不完备需要另构造反例。
\end{notebox}

\section{充分性}

\subsection{充分性的条件分布定义}
\begin{definitionbox}{定义 2.3.1：充分统计量}
若样本 $\mathbf X$ 给定 $T(\mathbf X)$ 的条件分布可以对所有参数选取同一个不含参数的版本，则称 $T$ 充分。
\end{definitionbox}

充分性不要求 $T$ 的边缘分布不含参数，也不表示 $\mathbf X$ 与 $T$ 独立。完整样本本身总是充分，故“充分”也不自动表示低维或简洁。

Bernoulli 样本中，若 $\sum_i x_i=t$，则
\[
 P_p(\mathbf X=\mathbf x\mid T=t)
 =\frac{p^t(1-p)^{n-t}}{\binom nt p^t(1-p)^{n-t}}
 =\binom nt^{-1}.
\]
否则条件概率为零。给定成功次数后，所有相容排列的条件规律都不含 $p$，故总和充分。

\subsection{共同支配概率的构造与用途}
\begin{definitionbox}{引理 2.3.1：共同支配概率}
非空模型 $\mathcal P$ 若被某个 $\sigma$-有限测度支配，则可选有限或可数个 $P_j\in\mathcal P$ 及严格正权重 $w_j$，使
\[
 \sum_jw_j=1,\qquad Q=\sum_jw_jP_j,\qquad P\ll Q\quad\forall P\in\mathcal P.
\]
讲义用无穷序列表示；有限族也可用重复分布写成序列，或直接用有限和。
\end{definitionbox}

\begin{notebox}[追问：是不是统一在一个概率分布下研究？]
可以这样理解，但统一的是参考分布，并没有把不同的 $P_\theta$ 变成同一个分布：
\[
 E_\theta f(X)=E_Q\left[f(X)\frac{\mathrm dP_\theta}{\mathrm dQ}(X)\right].
\]
不同分布的差异保留在各自的权重中。$Q$ 固定，不随当前研究的 $\theta$ 改变；$Q$ 不必属于原模型，也不表示真实数据来自混合机制。引理只适用于满足共同 $\sigma$-有限支配条件的模型，不是所有可能的离散、连续分布都可任意混在一起。
\end{notebox}

定义上，$P\ll Q$ 意味着 $Q(A)=0\Rightarrow P(A)=0$，不是 $P(A)\le Q(A)$。严格正权重保证
\[
 Q(A)=0\Rightarrow P_j(A)=0\quad\forall j.
\]
单个模型成员未必能支配整个族。对 $\{U(0,\theta):\theta>0\}$，可取
\[
 Q=\sum_{j\ge1}2^{-j}U(0,j).
\]
任取 $\theta$，选整数 $j\ge\theta$，则 $U(0,\theta)\ll U(0,j)\ll Q$。选择 $j$ 可以依赖 $\theta$，但 $Q$ 已经固定。

\subsubsection{补充证明：用可数个正密度区域覆盖整个模型}
讲义以“本质上确界可由可数子族实现”为工具。下面给出针对本引理的集合版本，避免直接调用那个函数定理。

首先将支配测度 $\nu$ 换成等价有限测度。取可数可测分割 $\mathcal X=\bigcup_mE_m$，每个 $\nu(E_m)<\infty$，令
\[
 \mu(A)=\sum_{m\ge1}2^{-m}\frac{\nu(A\cap E_m)}{1+\nu(E_m)}.
\]
则 $\mu(\mathcal X)\le1$，且 $\mu(A)=0\iff\nu(A)=0$。

对每个 $P$ 取密度 $p_P=\mathrm dP/\mathrm d\nu$，记 $S_P=\{p_P>0\}$，定义
\[
 M=\sup\{\mu(S_{P_1}\cup\cdots\cup S_{P_r}):r\ge1,\ P_i\in\mathcal P\}.
\]
$M$ 有限。对每个 $n$，选一个有限并 $B_n$ 使 $\mu(B_n)>M-1/n$，再令 $B=\bigcup_nB_n$。由下连续性和有限并的定义，
\[
 \mu(B)=M.
\]
确实，$\mu(B)\ge\mu(B_n)>M-1/n$；另一方面每个 $B_1\cup\cdots\cup B_N$ 仍是有限个正密度区域的并，其测度不超过 $M$。

任取其他分布 $P$，同理 $\mu(B\cup S_P)\le M$，所以
\[
 \mu(S_P\setminus B)=0,\qquad \nu(S_P\setminus B)=0.
\]
组成 $B$ 的模型成员至多可数，将它们记为 $P_j$，取正权重构成 $Q$。它的密度是
\[
 q=\sum_jw_jp_j,\qquad\{q>0\}=B.
\]
若 $Q(A)=0$，则 $\int_Aq\,\mathrm d\nu=0$，由 $q>0$ 于 $B$ 得 $\nu(A\cap B)=0$。任意 $p_P$ 又在 $B$ 外几乎处处为零，故
\[
 P(A)=\int_Ap_P\,\mathrm d\nu=0.
\]
因此每个 $P\ll Q$，引理得证。

\subsection{Fisher--Neyman 因子分解定理}
\begin{definitionbox}{定理 2.3.1：因子分解}
在讲义的标准 Borel 样本空间和值域、共同 $\sigma$-有限支配等条件下，$T$ 充分，当且仅当联合密度可写为
\[
 p_\theta(x)=a_\theta(T(x))h(x)
\]
（按相应几乎处处意义理解），其中 $h$ 不含参数。这里为简化记号，把完整样本也写作 $X$，其联合分布记为 $P_\theta$。
\end{definitionbox}

\subsubsection{方向一：因子分解推出充分性}
离散情形下，在同一个纤维 $\{x:T(x)=t\}$ 内，
\[
 P_\theta(X=x\mid T=t)
 =\frac{a_\theta(t)h(x)}
 {a_\theta(t)\sum_{y:T(y)=t}h(y)}
 =\frac{h(x)}{\sum_{y:T(y)=t}h(y)}.
\]
分母为正时参数因子消去。连续模型不能直接除以零概率事件，故使用共同概率 $Q$。

由引理，取 $Q=\sum_jw_jP_{\theta_j}$。其密度为
\[
 q(x)=h(x)\bar a(T(x)),\qquad
 \bar a(t)=\sum_jw_ja_{\theta_j}(t).
\]
在有效正密度集合上，
\[
 \frac{\mathrm dP_\theta}{\mathrm dQ}(x)
 =\frac{a_\theta(T(x))}{\bar a(T(x))}
 =:r_\theta(T(x)).
\]
零集上任意补定义。换分布时的权重只依赖 $T$，这正是离散消因子的测度论版本。

对任意事件 $A$，先在固定的 $Q$ 下定义
\[
 K_A(T)=E_Q(\ind_A\mid T).
\]
它是与参数无关的同一个函数。要证明其也是 $P_\theta$ 下的条件概率，除可测性外，只需对每个 $C\in\sigma(T)$ 验证积分：
\[
 \begin{aligned}
 \int_CK_A(T)\,\mathrm dP_\theta
 &=\int_CK_A(T)r_\theta(T)\,\mathrm dQ\\
 &=\int_C\ind_A r_\theta(T)\,\mathrm dQ\\
 &=P_\theta(A\cap C).
 \end{aligned}
\]
第二步使用 $\ind_Cr_\theta(T)$ 为 $\sigma(T)$-可测非负权重；无界时先截断，再用单调收敛。由条件期望唯一性，$K_A(T)=E_\theta(\ind_A\mid T)$，$P_\theta$-a.s.。在标准 Borel 条件下可通过 $Q$ 下的正则条件分布核统一组织这些等式，从而 $T$ 充分。

\subsubsection{方向二：充分性推出因子分解}
充分性给出不含参数的共同版本
\[
 E_\theta(\ind_A\mid T)=K_A(T).
\]
混合各个分布，对于 $C\in\sigma(T)$，
\[
 \int_CK_A(T)\,\mathrm dQ
 =\sum_jw_jP_{\theta_j}(A\cap C)=Q(A\cap C).
\]
所以 $K_A(T)=E_Q(\ind_A\mid T)$，$Q$-a.s.。

令 $P_\theta^T,Q^T$ 为 $T$ 的边缘分布。由 $P_\theta\ll Q$ 得 $P_\theta^T\ll Q^T$，定义
\[
 r_\theta(t)=\frac{\mathrm dP_\theta^T}{\mathrm dQ^T}(t).
\]
这是边缘密度比，不是条件期望。对每个事件 $A$，
\[
 \begin{aligned}
 P_\theta(A)
 &=E_\theta K_A(T)
 =\int K_A(t)\,\mathrm dP_\theta^T(t)\\
 &=\int K_A(t)r_\theta(t)\,\mathrm dQ^T(t)\\
 &=E_Q[K_A(T)r_\theta(T)]\\
 &=E_Q[\ind_A r_\theta(T)]
 =\int_A r_\theta(T)\,\mathrm dQ.
 \end{aligned}
\]
故 $\mathrm dP_\theta/\mathrm dQ=r_\theta(T)$，$Q$-a.s.。取 $h=\mathrm dQ/\mathrm d\nu$，RN 链式法则给出
\[
 p_\theta(x)=r_\theta(T(x))h(x)\quad\nu\text{-a.e.}
\]
因子分解得证。

\begin{notebox}[证明中各测度的作用]
引理提供一个固定参考概率 $Q$；因子分解使换测度权重只依赖 $T$；条件期望定义把“纤维内相对权重不变”变成严谨的共同条件分布。

样本空间积分 $E_Q[\phi(T)]$ 与边缘空间积分 $\int\phi(t)\,\mathrm dQ^T(t)$ 相等，但空间和测度不同。不能把 $Q$ 与 $Q^T$、条件期望与边缘密度比混为一谈。
\end{notebox}

\subsection{典型充分统计量与模型依赖性}
因子分解给出：Bernoulli、Poisson 样本的总和充分；$U(0,\theta)$ 的最大值充分。正态两参数模型中，充分统计量可以写为
\[
 \left(\sum_iX_i,\sum_iX_i^2\right).
\]
令 $V=\sum_i(X_i-\bar X)^2$，有
\[
 \bar X=\frac1n\sum_iX_i,\quad
 V=\sum_iX_i^2-n\bar X^2,
\]
\[
 \sum_iX_i=n\bar X,\qquad
 \sum_iX_i^2=V+n\bar X^2.
\]
所以 $(\bar X,V)$ 是同一信息的另一种编码。

已知正态方差时仅 $\bar X$ 已充分；两参数未知时一般还需 $V$。在全部连续密度构成的 i.i.d. 模型中，排序后的完整样本充分：给定排序值，所有排列等概率且与密度无关。讲义同时给出该大模型的最小充分性结论；这里不把它当作已从头证明的非参数定理。

若 $T$ 充分且 $T=\phi(S)$，由
\[
 p_\theta(x)=a_\theta(\phi(S(x)))h(x)
\]
可知 $S$ 也充分。增加信息不破坏充分性，但对充分统计量任意取函数可能丢失信息。

\section{最小充分性}

\subsection{最小充分性的定义}
充分性排除信息损失，却不排除无关细节。已知正态方差的模型中，完整样本、两块样本均值和总体样本均值都充分，但保留的细节不同。

\begin{definitionbox}{定义 2.4.1：最小充分统计量}
$T$ 对 $\mathcal P$ 充分，且对任意其他充分统计量 $S$，存在同一个不含参数的可测函数 $\phi$，使
\[
 T=\phi(S)\qquad P\text{-a.s.},\quad\forall P\in\mathcal P,
\]
则 $T$ 最小充分。
\end{definitionbox}

\begin{notebox}[追问：“最粗的精确压缩”是什么意思？]
$T=\phi(S)$ 的方向是“知道 $S$ 就能恢复 $T$”，因此 $S$ 至少包含 $T$ 的信息。最小充分统计量是每一种充分摘要都必须保留的共同核心。

统计量把样本按相同取值分组。摘要越粗，合并的样本越多；最小充分性要求在保持充分性的前提下尽可能合并。它不是简单比较向量维数。具有可测逆的一一变换仅改变编码，不改变信息。
\end{notebox}

\subsection{从子模型推广最小充分性}
\begin{definitionbox}{命题 2.4.1：最小充分性的推广原则}
设 $\mathcal P_0\subseteq\mathcal P$，且对每个事件 $A$，
\[
 [P(A)=0\ \forall P\in\mathcal P_0]
 \ \Longrightarrow\
 [P(A)=0\ \forall P\in\mathcal P].
\]
若 $T$ 对 $\mathcal P$ 充分，并对 $\mathcal P_0$ 最小充分，则 $T$ 对 $\mathcal P$ 最小充分。
\end{definitionbox}

这里“大、小”指候选分布的多少，不是样本量。例如正态均值遍历 $\R$ 是大模型，只保留均值为 $0,1$ 的两个候选分布是小模型，观察的样本向量并未改变。

\textbf{证明分四步。}
\begin{enumerate}
\item 任取对大模型充分的 $S$。同一个条件分布版本当然也适用于小模型，所以 $S$ 对小模型充分。
\item 由 $T$ 对小模型最小充分，存在同一个 $\phi$，使 $T=\phi(S)$ 在每个 $P\in\mathcal P_0$ 下 a.s. 成立。
\item 定义等式失败的事件
\[
 N=\{T\ne\phi(S)\}.
\]
则 $P(N)=0$ 对小模型每个分布成立。由假设，它对大模型每个分布也成立。
\item 因此任意大模型充分统计量 $S$ 都能恢复 $T$。结合已经假设的“大模型下 $T$ 充分”，得到最小充分性。
\end{enumerate}

\begin{notebox}[为什么不能省略共同零事件条件？]
取小模型 $\{U(0,1)\}$，大模型 $\{U(0,1),U(0,2)\}$。事件 $(1,2)$ 在小模型下概率为零，却在大模型某成员下有概率 $1/2$。一个等式在小模型几乎必然成立，不能自动推广到大模型。

证明中的异常集必须是 $\{T\ne\phi(S)\}$。
\end{notebox}

共同支配引理恰好提供满足零事件条件的可数小模型。若
\[
 Q=\sum_jw_jP_j,\qquad P\ll Q\quad\forall P\in\mathcal P,
\]
取 $\mathcal P_0=\{P_j\}$，则
\[
 P_j(A)=0\ \forall j\Rightarrow Q(A)=0
 \Rightarrow P(A)=0\ \forall P\in\mathcal P.
\]
这就是引理、推广原则和下面可数模型构造之间的衔接。

\subsection{可数模型的显式构造}
\begin{definitionbox}{定理 2.4.1：可数混合构造}
设 $\mathcal P_0=\{P_j:j\in J\}$ 为有限或可数族，取
\[
 Q=\sum_jw_jP_j,\quad r_j=\frac{\mathrm dP_j}{\mathrm dQ},
 \quad w_j>0,\quad\sum_jw_j=1.
\]
则 $R=(r_j)_{j\in J}$ 最小充分。若族内某个 $P_0$ 已支配所有 $P_j$，也可用 $(\mathrm dP_j/\mathrm dP_0)_j$。
\end{definitionbox}

\textbf{证明。} 取共同测度 $\nu$，记 $p_j=\mathrm dP_j/\mathrm d\nu$、$q=\mathrm dQ/\mathrm d\nu$。因为 $p_j=r_jq$，因子分解给出 $R$ 充分。

任取充分统计量 $S$，因子分解写为 $p_j(x)=a_j(S(x))h(x)$。由 $J$ 可数，可以合并各个等式的异常零集。于是
\[
 q(x)=h(x)\sum_jw_ja_j(S(x)).
\]
在有效正密度集合上，
\[
 r_j(x)=\frac{a_j(S(x))}{\sum_\ell w_\ell a_\ell(S(x))}.
\]
每个坐标都是 $S$ 的函数，因此整个可数向量 $R=\phi(S)$，$Q$-a.s.。由 $P_j\ll Q$，这个表示对每个 $P_j$ 也成立，故最小充分。用 $P_0$ 作参考时，同样由 $p_j/p_0=a_j(S)/a_0(S)$ 得到结论。

\subsection{似然比判别与典型例子}
\begin{definitionbox}{定理 2.4.2：似然比判别准则}
在讲义的受支配、标准 Borel 框架下，设 $T$ 已充分。使用讲义指定的密度版本，且在同一个共同零集之外，若
\[
 \left[
 p_P(x)=c(x,y)p_P(y)\ \forall P,\quad
 c(x,y)>0\text{ 不含 }P
 \right]\Rightarrow T(x)=T(y),
\]
则 $T$ 最小充分。
\end{definitionbox}

讲义的版本约定不可省略：使用证明中最小充分统计量 $S$ 给出的密度 $p_P(x)=a_P(S(x))h(x)$，并在共同满测 Borel 集上进行比较。不可数个参数分别拥有的零集，其并不一定还是零集。

\textbf{讲义证明的逻辑。} 引用 Bahadur 的存在性结果得到最小充分统计量 $S$。在共同满测集上取 $0<h<\infty$。若 $S(x)=S(y)$，则
\[
 p_P(x)=p_P(y)\frac{h(x)}{h(y)}\quad\forall P.
\]
假设推出 $T(x)=T(y)$。Blackwell--Mackey 定理及 Doob--Dynkin 引理保证 $T=\psi(S)$ 可测地成立。任意充分统计量 $U$ 都能恢复 $S$，进而恢复 $T$；加上 $T$ 充分，故最小充分。这里的存在性和可测表示结果是讲义引用的基础定理，本笔记不从描述集合论重新证明。

\begin{examplebox}{例 2.4.1：Bernoulli 与 Poisson}
Bernoulli 中，令 $\Delta=\sum_i x_i-\sum_i y_i$，
\[
 \frac{p_p^{(n)}(x)}{p_p^{(n)}(y)}
 =\left(\frac p{1-p}\right)^\Delta.
\]
$p/(1-p)$ 遍历正半轴，比值与 $p$ 无关当且仅当 $\Delta=0$。

Poisson 中，
\[
 \frac{p_\lambda^{(n)}(x)}{p_\lambda^{(n)}(y)}
 =\lambda^\Delta\frac{\prod_i y_i!}{\prod_i x_i!},
\]
同样当且仅当 $\Delta=0$。结合充分性，两个模型中的总和均最小充分。
\end{examplebox}

正态两参数模型中，令 $\Delta_1=\sum_i x_i-\sum_i y_i$、$\Delta_2=\sum_i x_i^2-\sum_i y_i^2$，则
\[
 \log\frac{p_{\mu,\sigma^2}^{(n)}(x)}{p_{\mu,\sigma^2}^{(n)}(y)}
 =\frac{\mu}{\sigma^2}\Delta_1-\frac1{2\sigma^2}\Delta_2.
\]
固定方差改变均值，得 $\Delta_1=0$；再改变方差，得 $\Delta_2=0$。反之两差为零时比值为 $1$。因此 $(\sum_iX_i,\sum_iX_i^2)$ 最小充分。

\subsubsection{例 2.4.3：指数族的两条证明路线}
设自然统计量维数为 $k$，存在 $\theta_0,\ldots,\theta_k$，使
\[
 v_i=\eta(\theta_i)-\eta(\theta_0),\quad i=1,\ldots,k
\]
线性无关。

第一条路线是选有限子模型。其对数密度比为
\[
 \log\frac{p_{\theta_i}(x)}{p_{\theta_0}(x)}
 =v_i^\top T(x)-[\zeta(\eta(\theta_i))-\zeta(\eta(\theta_0))].
\]
把 $v_i^\top$ 排成可逆矩阵 $V$，就可从这些对数比恢复 $VT$，继而恢复 $T$。比值向量在有限子模型中最小充分，所以 $T$ 也最小充分。共同正支撑保证共同零事件条件，再用命题 2.4.1 推广到整个模型。

第二条路线是直接比较整条似然。若 $p_\theta(x)/p_\theta(y)$ 与 $\theta$ 无关，令 $\Delta T=T(x)-T(y)$，对数等式减去 $\theta_0$ 时的等式得
\[
 v_i^\top\Delta T=0,\quad i=1,\ldots,k.
\]
$V$ 可逆，故 $\Delta T=0$，由判别准则得最小充分。

这里“存在 $k+1$ 个仿射独立自然参数”比“包含开集”更弱。弯曲族仍可能满足本条最小充分性条件，不能把它与后面的完备性开集条件混淆。

对 $U(0,\theta)$ 的正样本，似然为正的参数范围是 $(x_{(n)},\infty)$。两条完整似然成比例当且仅当最大值相同，因此最大值最小充分。

\section{完备性}

\subsection{完备性与函数唯一性}
充分性和最小充分性讨论保留哪些数据。完备性改问：两个关于 $T$ 的可积函数，如果在所有参数下期望相同，是否必须几乎必然相同？

\begin{definitionbox}{定义 2.5.1：完备与有界完备}
对任意不含未知参数的可测函数 $g$，若
\[
 E_\theta|g(T)|<\infty,\quad E_\theta g(T)=0\quad\forall\theta
\]
必推出
\[
 g(T)=0\quad P_\theta\text{-a.s.},\quad\forall\theta,
\]
则 $T$ 完备。若只对有界可测 $g$ 要求该结论，则称有界完备。
\end{definitionbox}

完备推出有界完备。不能取 $g(T)=T-\theta$ 来证明不完备，因为这个函数含有未知参数。结论是随机变量几乎必然为零，不要求 $g$ 在不可能取到的点上也为零。

定义期望变换 $(Kg)(\theta)=E_\theta g(T)$，则完备性表示其核只有几乎必然为零的函数。若 $a(T),b(T)$ 均可积且对所有参数期望相同，对差 $g=a-b$ 应用完备性就有 $a(T)=b(T)$。这保证同一完备统计量的无偏函数的唯一性，但不保证无偏估计存在，也不表示所有基于完整数据的无偏估计都相同。

\subsection{有界完备且充分推出最小充分}
\begin{definitionbox}{定理 2.5.1}
在讲义的标准 Borel 条件下，$T$ 有界完备且充分，则 $T$ 最小充分。
\end{definitionbox}

\textbf{证明。} 任取充分统计量 $U$，选 $T$ 值域上的有界一一 Borel 编码 $b$，其像上的逆可测。实值统计量可取 $b(t)=\arctan t$。定义
\[
 q(U)=E_\theta[b(T)\mid U],\qquad
 r(T)=E_\theta[q(U)\mid T].
\]
由 $U,T$ 充分，两个条件期望依次可取同一个不含参数的版本。它们有界，且全期望给出
\[
 E_\theta[r(T)-b(T)]=0\quad\forall\theta.
\]
有界完备性推出 $r(T)=b(T)$ a.s.。分别对 $U,T$ 条件化，
\[
 \begin{aligned}
 E_\theta[b(T)q(U)]
 &=E_\theta[q(U)E_\theta(b(T)\mid U)]
 =E_\theta[q(U)^2],\\
 E_\theta[b(T)q(U)]
 &=E_\theta[b(T)E_\theta(q(U)\mid T)]
 =E_\theta[b(T)^2].
 \end{aligned}
\]
故
\[
 E_\theta[(b(T)-q(U))^2]
 =E_\theta b(T)^2-2E_\theta[b(T)q(U)]+E_\theta q(U)^2=0.
\]
非负变量期望为零，得 $b(T)=q(U)$ a.s.。在 $b$ 的像上取可测逆，在像外任意补定义，便有 $T=b^{-1}(q(U))$ a.s.。由于 $U$ 任意，$T$ 最小充分。

\begin{notebox}[为什么不能直接用塔式性质？]
事先不知道 $\sigma(T)$ 与 $\sigma(U)$ 谁包含谁，不能直接断言
$E_\theta[E_\theta(b(T)\mid U)\mid T]=b(T)$。证明中是先由完备性得到这个等式，再由平方期望消除残差。
\end{notebox}

\subsection{两种直接证明：多项式与积分上限}
\begin{examplebox}{例 2.5.1--2.5.2：二项分布的完备性}
令 $T\sim\operatorname{Binomial}(n,p)$，$p\in(0,1)$。若所有参数下 $E_pg(T)=0$，则
\[
 \sum_{t=0}^n g(t)\binom ntp^t(1-p)^{n-t}=0.
\]
除以 $(1-p)^n>0$，令 $u=p/(1-p)$，得到
\[
 \sum_{t=0}^ng(t)\binom ntu^t=0\quad\forall u>0.
\]
一个有限次多项式在整个区间上为零，所有系数必须为零。由 $\binom nt>0$，得 $g(t)=0$ 对 $t=0,\ldots,n$ 成立。故 $T$ 完备，Bernoulli 样本总和因而完备充分。
\end{examplebox}

若把模型限制为单个已知 $p_0\in(0,1)$，$n\ge1$，则 $g(T)=T-np_0$ 是合法的固定非零零期望函数，完备性反而失效。完备性也相对于整族参数定义。

\begin{examplebox}{例 2.5.3：均匀上端点模型的完备性}
若 $X_i$ 独立同分布于 $U(0,\theta)$，$T=X_{(n)}$，则对 $0<t<\theta$，
\[
 P_\theta(T\le t)=(t/\theta)^n,\qquad
 f_\theta^T(t)=nt^{n-1}\theta^{-n}.
\]
任取在所有参数下可积的 $g$，零期望给出
\[
 \int_0^\theta g(t)t^{n-1}\,\mathrm dt=0\quad\forall\theta>0.
\]
对每个 $B>0$，可积性保证 $g(t)t^{n-1}$ 在 $(0,B)$ 绝对可积。因此
$F(b)=\int_0^b g(t)t^{n-1}\,\mathrm dt$ 在 $[0,B]$ 绝对连续，且
\[
 F'(b)=g(b)b^{n-1}\quad\text{Lebesgue-a.e.}
\]
$F$ 恒为零且 $b>0$，故 $g(b)=0$ a.e.。对正整数 $B$ 取可数并得到整个正半轴上的结论。$T$ 有 Lebesgue 密度，所以每个参数下 $g(T)=0$ a.s.，完备性成立。
\end{examplebox}

\subsection{指数族完备性及开集的作用}
\begin{definitionbox}{定理 2.5.2：指数族完备性}
若 $T$ 的分布相对于 $\nu$ 的密度为
\[
 q_\eta(t)=e^{\eta^\top t-\zeta(\eta)}h_*(t),\quad\eta\in\Xi\subseteq\R^k,
\]
且 $\Xi$ 包含非空开集，则 $T$ 完备。这里 $\nu$ 不必是 $k$ 维 Lebesgue 测度。
\end{definitionbox}

\textbf{证明展开。} 任取所有参数下可积且零期望的 $g$，写 $g=g^+-g^-$。选内点 $\eta_0$ 与 $\varepsilon>0$ 使 $B(\eta_0,\varepsilon)\subseteq\Xi$。消去密度中的归一化因子，得到
\[
 \int g^+(t)e^{\eta^\top t}h_*(t)\,\nu(\mathrm dt)
 =\int g^-(t)e^{\eta^\top t}h_*(t)\,\nu(\mathrm dt).
\]
在 $\eta_0$ 处记共同有限值为 $C$。若 $C=0$，非负积分为零，直接得到 $g=0$ 于载体测度 $h_*\nu$ 几乎处处。

若 $C>0$，定义两个概率密度
\[
 f_\pm(t)=C^{-1}g^\pm(t)e^{\eta_0^\top t}h_*(t).
\]
对于 $\|s\|<\varepsilon$，使用参数 $\eta_0+s$ 的积分恒等式，
\[
 \begin{aligned}
 \int e^{s^\top t}f_+(t)\,\nu(\mathrm dt)
 &=C^{-1}\int g^+(t)e^{(\eta_0+s)^\top t}h_*(t)\,\nu(\mathrm dt)\\
 &=C^{-1}\int g^-(t)e^{(\eta_0+s)^\top t}h_*(t)\,\nu(\mathrm dt)\\
 &=\int e^{s^\top t}f_-(t)\,\nu(\mathrm dt).
 \end{aligned}
\]
可积性保证两边在该邻域有限。引用原点邻域内有限的多元矩母函数唯一性，得到两个概率测度相同，进而 $f_+=f_-$，$\nu$-a.e.。约去正因子可得 $g(t)h_*(t)=0$，$\nu$-a.e.。所有模型分布都关于载体测度绝对连续，故 $g(T)=0$ 在每个参数下 a.s. 成立。

开集允许所有足够小方向的指数倾斜，提供一个原点邻域上的矩母函数信息。“所有阶矩相同必同分布”没有一般保证，不能替代这里的唯一性定理。

\begin{notebox}[弯曲族的不完备必须另证]
单观测 $X\sim N(\theta,\theta^2)$，$\theta>0$ 时，
\[
 X/\theta\sim N(1,1),\qquad P_\theta(X<0)=\Phi(-1).
\]
令 $T=(X,X^2)$，取
\[
 g(T)=\ind_{\{X<0\}}-\Phi(-1).
\]
该函数有界、对所有参数期望为零，但两个可能取值均非零。因此此例不有界完备，也不完备。证明依据是这个具体反例函数，而不只是自然参数空间没有开集。
\end{notebox}

\section{辅助统计量与 Basu 定理}

\subsection{辅助性要求整个分布不含参数}
\begin{definitionbox}{定义 2.6.1：辅助统计量}
若 $A=A(\mathbf X)$ 的分布不依赖 $\theta$，即对每个可测集合 $B$，
$P_\theta(A\in B)$ 都是关于参数的常数，则 $A$ 辅助。
\end{definitionbox}

均值恒定不够。例如 $X\sim N(0,\sigma^2)$ 的均值始终为零，但分布随方差变化，因此 $X$ 对该模型不辅助。

正态样本 $n\ge2$ 时，标准化残差方向
\[
 A=\frac{(X_1-\bar X,\ldots,X_n-\bar X)}
 {\{\sum_i(X_i-\bar X)^2\}^{1/2}}
\]
对完整位置尺度模型辅助。写 $X_i=\mu+\sigma Z_i$，其中 $Z_i$ 独立标准正态，中心化消去 $\mu$，归一化消去正数 $\sigma$，所以 $A$ 成为标准正态向量的固定函数。分母为零是概率零事件，可任意补定义。

\subsection{Basu 定理的条件与完整证明}
\begin{definitionbox}{定理 2.6.1：Basu 定理}
若 $T$ 有界完备且充分，$A$ 辅助，则它们在模型中的每个分布下独立。
\end{definitionbox}

\textbf{证明。} 任取 $A$ 值域内的可测集合 $B$，辅助性给出不含参数的常数
\[
 c_B=P_\theta(A\in B).
\]
充分性给出不含参数的共同版本
\[
 q_B(T)=E_\theta[\ind_{\{A\in B\}}\mid T],\qquad 0\le q_B\le1.
\]
全期望公式给出 $E_\theta[q_B(T)-c_B]=0$ 对所有参数成立。差是 $T$ 的固定有界函数，由有界完备性，
\[
 q_B(T)=c_B\quad P_\theta\text{-a.s.}
\]
再取 $T$ 值域内的可测集合 $C$，
\[
 \begin{aligned}
 P_\theta(T\in C,A\in B)
 &=E_\theta[\ind_{\{T\in C\}}E_\theta(\ind_{\{A\in B\}}\mid T)]\\
 &=E_\theta[\ind_{\{T\in C\}}q_B(T)]\\
 &=P_\theta(T\in C)P_\theta(A\in B).
 \end{aligned}
\]
这就是独立性；若从联合测度角度表述，可由矩形集上的相等延拓到乘积 $\sigma$-域。

\begin{notebox}[三个条件各用在哪里？]
辅助性使 $c_B$ 不含参数；充分性使 $q_B$ 可以统一选择；有界完备性使“所有参数下均值为零”升级为“几乎必然为零”。无条件概率不含参数，不能直接推出条件概率是同一常数。
\end{notebox}

\subsection{正态与指数模型的独立性应用}
\begin{examplebox}{例 2.6.2：正态均值与样本方差独立}
设 $n\ge2$，$X_i$ 独立同分布于 $N(\mu,\sigma^2)$。先固定任意 $\sigma^2>0$，仅令 $\mu$ 变化。联合密度的参数项为
\[
 \exp\left\{\frac{\mu}{\sigma^2}\sum_i x_i-\frac{n\mu^2}{2\sigma^2}\right\}.
\]
所以 $\sum_iX_i$ 充分；自然参数 $\mu/\sigma^2$ 遍历 $\R$，由指数族定理完备。$\bar X$ 与总和一一对应，也完备充分。

写 $X_i=\mu+\sigma Z_i$，则
\[
 S^2=\frac{\sigma^2}{n-1}\sum_i(Z_i-\bar Z)^2.
\]
在固定方差的子模型中，其分布不含 $\mu$，所以辅助。Basu 定理给出 $\bar X$ 与 $S^2$ 独立；固定的 $\sigma^2$ 任意，因此对整个两参数正态模型中的每个参数点都成立。
\end{examplebox}

不能直接说 $S^2$ 对完整两参数模型辅助，因为其分布依赖 $\sigma^2$。上述证明也没有预先使用均值、方差独立性，避免循环论证。

\begin{examplebox}{例 2.6.3：指数尺度族}
设 $X_i$ 独立同分布，密度为 $\theta^{-1}e^{-x/\theta}\ind_{\{x>0\}}$。总和充分，自然参数 $-1/\theta$ 遍历 $(-\infty,0)$，故总和及其一一变换 $\bar X$ 完备充分。

令 $Z_i=X_i/\theta$，则 $Z_i$ 独立同分布于尺度为 $1$ 的指数分布，且
\[
 \frac{X_1}{\bar X}=\frac{Z_1}{\bar Z}.
\]
比例的分布不含参数，因此辅助。Basu 定理给出
\[
 \bar X\ \perp\!\!\!\perp\ X_1/\bar X.
\]
$n=1$ 时比例恒为 $1$；$n\ge2$ 时是非平凡结论。它不表示 $X_1$ 与 $\bar X$ 独立。
\end{examplebox}

\subsection{反例：最小充分并不意味着完备}
\begin{examplebox}{例 2.6.4：均匀位置族}
设 $X_i$ 独立同分布于 $U(\theta,\theta+1)$，$\theta\in\R$，$n\ge2$。联合似然为
\[
 p_\theta^{(n)}(\mathbf x)
 =\ind_{\{x_{(n)}-1<\theta<x_{(1)}\}}.
\]
在极差小于 $1$ 的共同满测样本集上，两条非零似然成比例当且仅当参数区间两个端点相同。因而 $T=(X_{(1)},X_{(n)})$ 最小充分。

写 $X_i=\theta+U_i$，$U_i$ 独立同分布于 $U(0,1)$。极差
\[
 R=X_{(n)}-X_{(1)}=U_{(n)}-U_{(1)}
\]
是 $T$ 的辅助函数。极小、极大的联合密度为
\[
 f(u,v)=n(n-1)(v-u)^{n-2},\qquad0<u<v<1.
\]
令 $r=v-u$，Jacobian 为 $1$，在 $0<u<1-r$ 上积分得
\[
 f_R(r)=n(n-1)r^{n-2}(1-r),\qquad0<r<1.
\]
因此
\[
 ER=n(n-1)\int_0^1r^{n-1}(1-r)\,\mathrm dr
 =n(n-1)\left(\frac1n-\frac1{n+1}\right)=\frac{n-1}{n+1}.
\]
于是 $g(T)=R-(n-1)/(n+1)$ 有界、所有参数下零期望，而 $R$ 有连续密度，故 $g(T)\ne0$ 的概率为 $1$。$T$ 不有界完备，也不完备。
\end{examplebox}

这里是区间长度固定、位置未知的模型，不能与上端点未知的 $U(0,\theta)$ 混淆。也可反证：若 $T$ 有界完备，Basu 定理给出 $R\perp\!\!\!\perp T$。由于 $R$ 是 $T$ 的函数，事件 $\{R\le r\}$ 属于 $\sigma(T)$；独立性便给出
\[
 P_\theta(R\le r)=P_\theta(R\le r)^2.
\]
所以这个概率只能为 $0$ 或 $1$。但 $R$ 的密度在 $(0,1)$ 上为正，任意 $r\in(0,1)$ 都满足 $0<P_\theta(R\le r)<1$，矛盾。这也解释了“一个实随机变量若与自身独立，便只能几乎必然为常数”。

\section{总结与学习检查}

\begin{notebox}[充分、最小充分、完备与辅助]
充分性：给定摘要后，剩下样本的条件分布是否与参数无关？

最小充分性：每个充分摘要是否都必须保留这份共同核心？

完备性：一个固定可积函数能否在所有参数下均值为零，却不几乎必然为零？

辅助性：统计量自身的整个分布是否与参数无关？
\end{notebox}

本章的可靠逻辑关系是
\[
 \text{完备}\Rightarrow\text{有界完备},\qquad
 \text{有界完备且充分}\Rightarrow\text{最小充分}\Rightarrow\text{充分}.
\]
最小充分不推出完备。常数统计量对任意模型都完备，却在分布随参数改变的模型中不充分，说明完备不是“信息最多”。

\subsection*{考场证明的起手与收尾}
\begin{enumerate}
\item 证明充分：写全联合密度和支撑，明确 $g_\theta(T(x))$ 与 $h(x)$，引用因子分解。
\item 证明最小充分：先有充分性，再证明完整似然成比例推出 $T(x)=T(y)$；或选择合适小模型并核对共同零事件。
\item 证明完备：任取所有参数下可积的固定函数 $g$，从零期望推出 $g(T)=0$ a.s.。用指数族定理时写明自然参数及其开集。
\item 否定完备：明确构造与未知参数无关的 $g$，核对可积、所有参数下零期望、至少一个参数下不几乎必然为零。
\item 应用 Basu：明确当前模型，分别核对充分、完备、辅助；若固定干扰参数使用子模型，最后说明所固定值任意。
\end{enumerate}

\subsection*{练习与核对}
\begin{enumerate}
\item 不看答案，解释为什么统计量和自然参数都落在曲线上时，两条满秩条件的判断仍可能不同。
\item 从条件期望定义出发，重写因子分解两个方向，标明每一步使用的测度。
\item 说明共同支配引理、可数混合构造与推广原则怎样共同支持最小充分性证明。
\item 独立证明二项与均匀最大值的完备性，并构造均匀位置族的不完备反例。
\item 不预先使用卡方分布或独立性结论，应用 Basu 证明正态样本均值与方差独立。
\end{enumerate}

\begin{notebox}[复习题答案线索]
第 1 题：对统计量检验 $a^\top T=b$，对自然参数检验同维开球；曲线可以仿射张成平面。第 2 题：先构造共同概率 $Q$，将密度比写为 $T$ 的函数，再验证条件期望定义，区分原样本测度与边缘测度。第 3 题：由可数混合得到共同零事件，由密度比向量得到最小充分，最后用推广原则回到全模型。第 4 题：二项模型用多项式系数，均匀上端点用绝对连续积分的导数，位置族用极差减去其常数期望。第 5 题：固定方差，仅令均值变化，此时均值完备充分、样本方差辅助；对每个固定方差应用 Basu。
\end{notebox}

\subsection*{来源与引用说明}
主要来源为 Fang Yao《Advanced Mathematical Statistics》第二章讲义，正文第 27--48 页，并结合第二章 Slides 与相关推导整理。原讲义文件未随笔记重新发布。一般最小充分性证明沿用讲义引用：Bahadur (1954), Theorem 6.2；Srivastava (1998), \emph{A Course on Borel Sets}, Theorem 4.5.7；Kallenberg (2002), \emph{Foundations of Modern Probability}, 2nd ed., Lemma 1.13。书目与编号按课程材料记录，未另行核对相应书籍全文。

\chapter{决策、估计与信息下界}

评价一个统计方法，须先指定模型、估计目标、损失函数及允许比较的规则类。无偏性约束估计量的均值，一致性描述样本量增大时的集中，而风险同时计入误差的大小与发生概率。本章先建立这些比较标准，再由得分恒等式推导 Cramér--Rao 下界，并讨论等号、参数变换和非正则模型。

\section{决策问题与风险比较}

\subsection{从估计规则到风险函数}
\begin{definitionbox}{统计决策问题}
给定统计模型 $\mathcal P=\{P_\theta:\theta\in\Theta\}$、行动空间 $\mathcal A$、非负损失 $L(\theta,a)$，以及一类可测规则 $\mathcal D$。规则 $\delta:\mathcal X\to\mathcal A$ 的风险为
\[
 R_\theta(\delta)=\E_\theta L(\theta,\delta(X)).
\]
固定参数时，损失 $L(\theta,\delta(X))$ 是随机变量，风险是对样本分布取平均后的数值；改变参数，便得到整条风险函数。
\end{definitionbox}

估计 $g(\theta)\in\R^k$ 时，常用平方损失 $\|a-g(\theta)\|^2$、绝对距离损失 $\|a-g(\theta)\|$，或加权二次损失
\[
 L(\theta,a)=(a-g(\theta))^\top W_\theta(a-g(\theta)),
 \qquad W_\theta\succeq0.
\]
平方损失下的风险称均方误差。权重可以改变不同方向的代价，但这种二次损失仍对误差的正负对称。高估与低估代价不同时，应另设非对称损失。

\begin{examplebox}{正态均值：两条相交的风险曲线}
设 $X_1,\ldots,X_n\overset{\mathrm{i.i.d.}}\sim N(\mu,\sigma^2)$，$\sigma^2$ 已知，估计 $\mu$。平方损失下，
\[
 R_\mu(\bar X)=\frac{\sigma^2}{n},\qquad
 R_\mu(m)=(m-\mu)^2,
\]
其中 $m$ 是事先固定的常数。若 $|\mu-m|<\sigma/\sqrt n$，常数规则的风险较低；区间外则样本均值较低。常数规则在 $\mu=m$ 处风险为零，并不能使它成为整个模型中的好规则。
\end{examplebox}

风险的定义先固定参数，只对数据平均。给参数配置先验 $\pi$ 后再平均，才得到 Bayes 风险
\[
 r_\pi(\delta)=\int_\Theta R_\theta(\delta)\,\pi(\mathrm d\theta).
\]
这两个期望的对象不同。

\subsection{支配、可容许与极小极大}
\begin{definitionbox}{三种风险比较}
若 $R_\theta(\delta_1)\le R_\theta(\delta_2)$ 对所有 $\theta$ 成立，且至少一点严格小于，称 $\delta_1$ 支配 $\delta_2$。在指定规则类中不受任何规则支配的规则称可容许。

若某个 $\delta^*$ 满足
\[
 R_\theta(\delta^*)\le R_\theta(\delta)
 \quad\text{对所有 }\theta\in\Theta,\ \delta\in\mathcal D,
\]
则它逐点最优。极小极大规则则要求达到
\[
 \inf_{\delta\in\mathcal D}\sup_{\theta\in\Theta}R_\theta(\delta).
\]
\end{definitionbox}

逐点最优必然可容许；若逐点最优规则存在，任何可容许规则都须与它具有相同风险函数，否则便被支配。可容许性本身不提供风险的大小：它仅排除一种统一改进。规则类扩大后，原来的可容许性也可能失效。

例如 $X_i\sim N(\theta,1)$ 且 $n\ge2$，$X_1$ 与 $\bar X$ 均无偏，但风险分别为 $1$ 与 $1/n$，故 $X_1$ 不可容许。这个比较尚未证明 $\bar X$ 优于所有估计量。又如 $X\sim N(\theta,1)$，在只含 $X$ 与常数 $3$ 的两元素规则类中，二者风险相交，因此都可容许，却无逐点最优规则。

极小极大比较先固定规则，再找它最坏的参数点，最后在规则间优化。交换这两个优化顺序，会变成知道参数后另选行动的问题。一个实用下界为
\[
 \inf_\delta\sup_\theta R_\theta(\delta)
 \ge \inf_\delta r_\pi(\delta),
\]
因为任意规则的最大风险不小于其先验平均风险。若某个规则的最大风险等于右端，就证明了它的极小极大性。

\subsection{检验也是决策}
检验 $H_0:\theta\in\Theta_0$ 对 $H_1:\theta\in\Theta_1$，约定拒绝为行动 $1$。随机化检验 $\varphi(X)\in[0,1]$ 表示给定数据后的拒绝概率，可用独立 $V\sim U(0,1)$ 实现为 $\ind\{V\le\varphi(X)\}$。记功效
\[
 \beta_\varphi(\theta)=\E_\theta\varphi(X).
\]
若错误行动的损失为 $1$、正确行动为 $0$，则
\[
 R_\theta(\varphi)=
 \begin{cases}
 \beta_\varphi(\theta),&\theta\in\Theta_0,\\
 1-\beta_\varphi(\theta),&\theta\in\Theta_1.
 \end{cases}
\]
水平不超过 $\alpha$ 要求 $\sup_{\Theta_0}\beta_\varphi\le\alpha$。随后提高备择下的功效，才是带水平约束的风险优化。总不拒绝虽然没有第一类错误，却没有发现备择的能力。

\paragraph{复习题}
\begin{enumerate}
 \item 在上面的正态均值例中，把损失改成 $|a-\mu|$，求 $\bar X$ 与常数 $m$ 的风险及交点。
 \item 逐点最优、可容许、极小极大是否是同一个概念？说明其量词或优化对象的区别。
\end{enumerate}
\begin{notebox}[答案线索]
第一题：$R_\mu(\bar X)=\sigma\sqrt{2/(\pi n)}$，$R_\mu(m)=|m-\mu|$。第二题：逐点最优比较每个竞争者在每个参数点的风险；可容许只排除支配；极小极大只比较最大风险。有限规则类中的结论不能直接推广到所有可测规则。
\end{notebox}

\section{偏差、方差与收缩}

\subsection{均方误差分解}
设 $T$ 二阶可积，目标为 $g(\theta)\in\R$，偏差记作
\[
 b_\theta(T)=\E_\theta T-g(\theta).
\]
将误差拆为 $T-\E_\theta T+b_\theta(T)$，平方后取期望，交叉项因中心化消失，故
\[
 \operatorname{MSE}_\theta(T)
 =\Var_\theta(T)+b_\theta(T)^2.
\]
无偏性只使第二项为零，不保证第一项小。允许适当偏差，有时可以显著降低方差。

\begin{examplebox}{固定系数的正态收缩}
设 $X\sim N(\theta,1)$，$0<a<1$，比较 $aX$ 与 $X$。有
\[
 b_\theta(aX)=(a-1)\theta,\qquad
 R_\theta(aX)=a^2+(1-a)^2\theta^2.
\]
因此
\[
 R_\theta(aX)<R_\theta(X)=1
 \iff |\theta|<\sqrt{\frac{1+a}{1-a}}.
\]
收缩降低了方差，却引入随 $|\theta|$ 增大的偏差；整个实参数空间上不存在这里的统一支配关系。
\end{examplebox}

\subsection{二项比例的加权估计}
设 $Y\sim\operatorname{Bin}(n,p)$，$a,b\ge0$，考虑
\[
 \delta_{a,b}=\frac{Y+a}{n+a+b}.
\]
当 $a+b>0$ 时，它是样本比例与固定中心的加权平均：
\[
 \delta_{a,b}=w\frac Yn+(1-w)q,
 \qquad w=\frac n{n+a+b},\quad q=\frac a{a+b}.
\]
直接计算得
\[
 b_p(\delta_{a,b})=\frac{a-(a+b)p}{n+a+b},\qquad
 R_p(\delta_{a,b})=
 \frac{np(1-p)+\{a-(a+b)p\}^2}{(n+a+b)^2}.
\]
取 $n=16,a=b=1$，令 $u=p(1-p)$，则两个风险分别为 $(1+12u)/324$ 与 $u/16$。前者较小当且仅当
\[
 u>\frac4{33}
 \iff \left|p-\frac12\right|<\frac12\sqrt{\frac{17}{33}}.
\]
靠近中心时，减少的方差补偿了偏差；靠近端点时则相反。

\subsection{对称收缩的极小极大计算}
在 $\delta_c=(Y+c)/(n+2c)$、$c\ge0$ 的规则类内，风险为
\[
 R_p(c)=\frac{np(1-p)+c^2(1-2p)^2}{(n+2c)^2},
 \qquad
 \frac{\partial R_p(c)}{\partial p}
 =\frac{(n-4c^2)(1-2p)}{(n+2c)^2}.
\]
当 $c<\sqrt n/2$，最大风险在 $p=1/2$；当 $c>\sqrt n/2$，在 $p=0,1$。因而
\[
 \sup_{0\le p\le1}R_p(c)=
 \begin{cases}
 \dfrac{n}{4(n+2c)^2},&0\le c\le\sqrt n/2,\\[5pt]
 \dfrac{c^2}{(n+2c)^2},&c\ge\sqrt n/2.
 \end{cases}
\]
第一段随 $c$ 递减，第二段递增。最优选择是
\[
 c^*=\frac{\sqrt n}{2},\qquad
 R_p(c^*)=\frac1{4(\sqrt n+1)^2},
\]
此时风险不依赖 $p$。

\begin{notebox}[从类内优化到全部规则：还需一个下界]
上述微分计算只证明对称规则类内的极小极大性。本例可以再补一步：取真先验 $p\sim\operatorname{Beta}(c^*,c^*)$。平方损失下的后验均值正是 $\delta_{c^*}$，因为条件均方误差由后验均值最小化。它的 Bayes 风险等于常数风险 $1/\{4(\sqrt n+1)^2\}$。Bayes 下界与该规则的最大风险相等，故它实际上也在全部规则中极小极大。若不补这个论证，便不能扩大原有结论的范围。
\end{notebox}

固定 $a,b$ 时，加权估计的偏差为 $O(n^{-1})$；但 $c^*=\sqrt n/2$ 随样本量改变，偏差一般仅为 $O(n^{-1/2})$。两种渐近比较的前提不同。

\paragraph{复习题}
\begin{enumerate}
 \item 设独立 $X_i$ 具有共同均值 $\theta$ 与方差 $\sigma^2>0$。证明所有确定权重的线性无偏估计 $\sum_i c_iX_i$ 中，$\bar X$ 的方差最小，并求等号条件。
 \item 当 $n=16$ 时，分别计算 $Y/n$ 与 $\delta_{c^*}$ 的最大风险。偏差是否必然增加最大均方误差？
\end{enumerate}
\begin{notebox}[答案线索]
第一题：无偏要求 $\sum_i c_i=1$，Cauchy--Schwarz 给出 $\sum_i c_i^2\ge1/n$；等号当且仅当所有 $c_i=1/n$。结论只在线性无偏类中成立。第二题：最大风险分别为 $1/64$ 与 $1/100$，故有偏规则也可能降低最大风险。
\end{notebox}

\section{无偏性、一致性与渐近尺度}

\subsection{三种性质应分别检验}
\begin{definitionbox}{无偏、渐近无偏与一致}
$T_n$ 对 $g(\theta)$ 无偏，指每个 $n$ 均有 $\E_\theta T_n=g(\theta)$。若期望存在且 $\E_\theta T_n\to g(\theta)$，称渐近无偏。若对每个固定 $\theta$、每个 $\varepsilon>0$，
\[
 P_\theta(|T_n-g(\theta)|>\varepsilon)\to0,
\]
则称一致。给定 $a_n\to\infty$，若 $a_n(T_n-g(\theta))=O_{P_\theta}(1)$，称具有 $a_n$-一致性。
\end{definitionbox}

依概率有界可推出 $T_n-g(\theta)=O_{P_\theta}(a_n^{-1})\pto0$。这些都是固定参数下的陈述，未必在整个参数空间上一致，也未声称该速度最优。

由 Markov 不等式，
\[
 P_\theta(|T_n-g(\theta)|>\varepsilon)
 \le\frac{\E_\theta(T_n-g(\theta))^2}{\varepsilon^2}.
\]
故偏差与方差都趋零足以推出一致；反向需要额外的尾部控制。

\begin{examplebox}{无偏性与一致性互不替代}
设 $X_i\sim N(\mu,1)$。
\begin{enumerate}
 \item $T_n=(1-1/n)\bar X$ 有偏，偏差 $-\mu/n$，方差 $(1-1/n)^2/n$，却一致。
 \item $T_n=X_1$ 每个样本量均无偏，但
 \[
 P_\mu(|T_n-\mu|>\varepsilon)=P(|Z|>\varepsilon)>0,
 \quad Z\sim N(0,1),
 \]
 所以不一致。
 \item 允许外加随机化，取与数据独立的 $U\sim U(0,1)$，令
 \[
 T_n=\bar X+n\ind\{U\le1/n\}.
 \]
 则 $\E_\mu T_n=\mu+1$，不渐近无偏，但
 \[
 P_\mu(|T_n-\mu|>\varepsilon)
 \le \frac1{n\varepsilon^2}+\frac1n\to0.
 \]
 此时均方误差为 $1/n+n$，甚至发散。若规则只允许原始样本的确定函数，这个随机化例子须先把 $U$ 纳入实验。
\end{enumerate}
\end{examplebox}

证明不一致应回到尾概率，而不能仅以“方差未趋零”代替。相反，上例说明依概率收敛本身并不保证期望或二阶矩收敛。

\subsection{一阶相同不等于有限样本相同}
设 $X_i\sim N(\mu,\sigma^2)$，$\sigma^2$ 已知，估计 $\mu^2$。考虑
\[
 T_n=\bar X^2,\qquad U_n=\bar X^2-\frac{\sigma^2}{n}.
\]
前者偏差为 $\sigma^2/n$，后者无偏。利用恒等式
\[
 \bar X^2-\mu^2=2\mu(\bar X-\mu)+(\bar X-\mu)^2
\]
及 $\sqrt n(\bar X-\mu)=O_P(1)$，有
\[
 \sqrt n(T_n-\mu^2)
 =2\mu\sqrt n(\bar X-\mu)+o_P(1).
\]
当 $\mu\ne0$，两个估计量均满足
\[
 \sqrt n(T_n-\mu^2),\ \sqrt n(U_n-\mu^2)
 \dto N(0,4\mu^2\sigma^2),
\]
因为 $\sqrt n(T_n-U_n)=\sigma^2/\sqrt n\to0$。

在 $\mu=0$ 处，一阶导数消失。写 $Z_n=\sqrt n\bar X/\sigma\sim N(0,1)$，便有精确分布
\[
 nT_n\overset d=\sigma^2\chi_1^2,
 \qquad nU_n\overset d=\sigma^2(\chi_1^2-1).
\]
此处应使用 $n$ 尺度，且偏差修正改变极限分布。尺度和有效的一阶展开可能随参数点改变。

\paragraph{复习题}
\begin{enumerate}
 \item 若 $\sqrt n(T_n-g(\theta))\dto N(0,V)$，能否推出一致性？能否直接推出 $n\operatorname{MSE}_\theta(T_n)\to V$？
 \item 计算上述 $U_n$ 的精确方差，比较 $\mu=0$ 与 $\mu\ne0$ 时的误差尺度。
\end{enumerate}
\begin{notebox}[答案线索]
第一题：弱收敛蕴含依概率有界，因此推出一致；矩收敛不能直接推出，需统一可积性等附加条件。第二题：令 $Z=\bar X-\mu$，使用 $\E Z^3=0$、$\E Z^4=3\sigma^4/n^2$，得到
\[
 \Var_\mu(U_n)=\frac{4\mu^2\sigma^2}{n}+\frac{2\sigma^4}{n^2}.
\]
\end{notebox}

\section{得分与 Fisher 信息}

\subsection{定义与正则条件}
令 $\theta\in\Theta\subset\R^d$，单次观测密度为 $p_\theta$，得分取列向量：
\[
 \ell_\theta(x)=\log p_\theta(x),\qquad
 s_\theta(x)=\nabla_\theta\ell_\theta(x),\qquad
 I(\theta)=\E_\theta[s_\theta(X)s_\theta(X)^\top].
\]
信息矩阵总半正定，因为 $a^\top I(\theta)a=\E_\theta(a^\top s_\theta)^2\ge0$，但不一定可逆。导数针对参数，求导时观测值固定；把 $X$ 代入后，得分才成为随机向量。

以下是一组足以使用信息恒等式的局部条件：参数点位于开放参数空间内部；各密度在同一支撑上为正；除一个共同测度零集外，密度关于参数二次连续可微；密度的一、二阶导数在真值邻域内具有可积包络；得分二阶矩与对数密度 Hessian 的期望存在。这些条件使微分与积分可以交换。

\begin{definitionbox}{得分与信息恒等式}
在上述条件下，
\[
 \E_\theta s_\theta(X)=0,\qquad
 I(\theta)=\Cov_\theta(s_\theta(X))
 =-\E_\theta\nabla_\theta^2\ell_\theta(X).
\]
\end{definitionbox}

\paragraph{证明}
由 $\nabla p_\theta=p_\theta s_\theta$ 和密度归一化，
\[
 \E_\theta s_\theta
 =\int\nabla p_\theta\,\mathrm d\nu
 =\nabla\int p_\theta\,\mathrm d\nu=0.
\]
逐坐标使用乘积法则，
\[
 \nabla^2p_\theta
 =p_\theta\{\nabla^2\ell_\theta+s_\theta s_\theta^\top\}.
\]
再对两边积分，左边为归一化常数的二阶导数 $0$，故得第二条恒等式。单个观测的负 Hessian 不必半正定；这里等于信息矩阵的是它的期望。

对于 i.i.d. 样本，记
\[
 \ell_n(\theta)=\sum_{i=1}^n\ell_\theta(X_i),\quad
 S_n(\theta)=\sum_{i=1}^ns_\theta(X_i),\quad
 H_n(\theta)=\nabla_\theta^2\ell_n(\theta).
\]
独立性与得分均值为零使外积中的交叉项消失，故
\[
 I_n(\theta)=\E_\theta S_nS_n^\top=nI(\theta).
\]
独立但不同分布时信息仍逐项相加；存在依赖时不能直接写成 $nI$。

\subsection{指数族的归一化函数}
自然指数族写成
\[
 p_\eta(x)=h(x)\exp\{\eta^\top T(x)-\zeta(\eta)\},
 \qquad Z(\eta)=e^{\zeta(\eta)}.
\]
在自然参数空间内部，直接微分归一化积分得到
\[
 \nabla\zeta(\eta)=\frac{\nabla Z(\eta)}{Z(\eta)}=\E_\eta T,
\]
\[
 \nabla^2\zeta(\eta)
 =\frac{\nabla^2Z}{Z}-\frac{\nabla Z\nabla Z^\top}{Z^2}
 =\Cov_\eta(T).
\]
另一方面，$s_\eta=T-\nabla\zeta$，所以 $I(\eta)=\nabla^2\zeta(\eta)$。若 $T$ 无非平凡的几乎必然仿射关系，则任意非零 $a$ 均有 $\Var(a^\top T)>0$，信息矩阵正定。

微分合法性来自参数空间的内部性。一维时，选取 $\eta_0\pm2a$ 仍在自然参数空间内。对 $m=1,2$ 及 $|\eta-\eta_0|\le a$，
\[
 |t|^me^{\eta t}
 \le C_{m,a}\{e^{(\eta_0+2a)t}+e^{(\eta_0-2a)t}\}.
\]
乘 $h$ 后右边可积，给出导数所需的局部包络。多维时可用参数邻域内有限个立方体顶点的指数项之和控制。

\subsection{参数变换与 Jacobian 的方向}
设新参数 $\lambda\in\R^r$，旧参数为 $\theta=\psi(\lambda)\in\R^d$，Jacobian 定义为
\[
 J(\lambda)=D\psi(\lambda)
 =\left(\frac{\partial\theta_j}{\partial\lambda_r}\right)_{j,r}
 \in\R^{d\times r}.
\]
链式法则给出
\[
 s_\lambda(X)=J(\lambda)^\top s_\theta(X),\qquad
 I_\lambda(\lambda)=J(\lambda)^\top
 I_\theta(\psi(\lambda))J(\lambda).
\]
导数方向是“旧参数对新参数”，因为新得分对 $\lambda$ 求导。这里没有变换样本 $X$；Jacobian 来自参数的链式求导，不是密度换元的行列式。

\begin{examplebox}{指数分布：尺度与率}
尺度参数模型 $p_\theta(x)=\theta^{-1}e^{-x/\theta}\ind\{x>0\}$ 有
\[
 s_\theta(x)=\frac{x-\theta}{\theta^2},\qquad
 I_\theta(\theta)=\frac1{\theta^2}.
\]
改用率 $\lambda=1/\theta$，$J=\mathrm d\theta/\mathrm d\lambda=-1/\lambda^2$，因此
\[
 I_\lambda(\lambda)=\frac1{\lambda^4}\,I_\theta(1/\lambda)
 =\frac1{\lambda^2}.
\]
直接对 $\log p_\lambda(x)=\log\lambda-\lambda x$ 求导，也有 $s_\lambda=1/\lambda-X$，其方差为 $1/\lambda^2$。只把 $\theta=1/\lambda$ 代入旧信息所得的 $\lambda^2$，仍是旧坐标导数的二阶矩，缺少链式因子。
\end{examplebox}

\subsection{常见模型核对}
\begin{center}
\begin{tabularx}{\textwidth}{@{}lXX@{}}
\toprule
模型与参数 & 单次得分 & 单次信息\\
\midrule
$N(\mu,\sigma^2)$，$\sigma^2$ 已知 & $(X-\mu)/\sigma^2$ & $1/\sigma^2$\\
$\operatorname{Bernoulli}(p)$ & $(X-p)/\{p(1-p)\}$ & $1/\{p(1-p)\}$\\
$\operatorname{Poisson}(\lambda)$ & $X/\lambda-1$ & $1/\lambda$\\
指数分布，尺度 $\theta$ & $(X-\theta)/\theta^2$ & $1/\theta^2$\\
\bottomrule
\end{tabularx}
\end{center}
Bernoulli 的自然参数为 $\eta=\log\{p/(1-p)\}$，其信息为 $p(1-p)$，不是原参数 $p$ 下的信息。信息的数值必须连同参数坐标一起解释。

\paragraph{复习题}
\begin{enumerate}
 \item 若 Poisson 均值改写为 $\lambda=e^\eta$，用链式法则计算 $I_\eta$，再用对数密度直接验证。
 \item 为什么从 $\int p_\theta=1$ 合法求导，尚不足以保证可以从 $\int T p_\theta$ 合法求导？
\end{enumerate}
\begin{notebox}[答案线索]
第一题：$I_\eta=(e^\eta)^2/e^\eta=e^\eta$；新得分是 $X-e^\eta$。第二题：多出的 $T$ 可能无界，需要另行控制 $T\,\partial_\theta p_\theta$ 的可积性，例如提供局部可积包络。
\end{notebox}

\section{Cramér--Rao 下界与等号条件}

\subsection{标量下界的完整证明}
\begin{definitionbox}{标量 Cramér--Rao 不等式}
设标量参数模型满足前述得分恒等式，$0<I_n(\theta)<\infty$。统计量 $T$ 无偏估计可微目标 $g(\theta)$，且 $\E_\theta T^2<\infty$。另假设在 $\E_\theta T=\int T p_\theta^{(n)}\,\mathrm d\nu^n$ 中可以交换微分与积分。则
\[
 \Cov_\theta(T,S_n)=g'(\theta),\qquad
 \Var_\theta(T)\ge\frac{g'(\theta)^2}{I_n(\theta)}.
\]
i.i.d. 情形 $I_n=nI$。
\end{definitionbox}

\paragraph{证明}
$T=T(x)$ 不依赖未知参数。对无偏性求导，
\[
 \begin{aligned}
 g'(\theta)
 &=\int T(x)\,\partial_\theta p_\theta^{(n)}(x)\,\mathrm d\nu^n(x)\\
 &=\E_\theta[T S_n(\theta)]
 =\Cov_\theta(T,S_n),
 \end{aligned}
\]
末步用了 $\E_\theta S_n=0$。Cauchy--Schwarz 不等式给出
\[
 g'(\theta)^2
 =\{\E_\theta[(T-g(\theta))S_n]\}^2
 \le\Var_\theta(T)I_n(\theta).
\]
除以正的信息即得下界。等号成立当且仅当
\[
 T-g(\theta)=\frac{g'(\theta)}{I_n(\theta)}S_n(\theta)
 \quad P_\theta\text{-a.s.}
\]
当 $g'(\theta)=0$，下界为零，取等要求 $T=g(\theta)$ 几乎必然，而非任意估计量均取等。

\subsection{投影分解与下界不可达的可能性}
令
\[
 r=T-g(\theta)-\frac{g'(\theta)}{I_n(\theta)}S_n(\theta).
\]
则 $\E_\theta r=0$、$\E_\theta(rS_n)=0$。平方展开得到精确分解
\[
 \Var_\theta(T)=\frac{g'(\theta)^2}{I_n(\theta)}+\E_\theta r^2.
\]
这是 $L^2(P_\theta)$ 中把误差投影到得分方向的勾股关系。残差正交表示不相关，不一般表示独立。投影项可能依赖未知参数，因此不能自动把残差删去、构造一个合法的更好估计量；信息下界可能不可达。

\begin{examplebox}{正态均值与均值平方}
已知 $\sigma^2$ 时，$S_n(\mu)=n(\bar X-\mu)/\sigma^2$，$I_n=n/\sigma^2$。估计 $g(\mu)=\mu$，样本均值满足等号条件，方差 $\sigma^2/n$ 达到下界。

估计 $g(\mu)=\mu^2$ 时，下界为 $4\mu^2\sigma^2/n$，而无偏估计 $U_n=\bar X^2-\sigma^2/n$ 的方差为
\[
 \Var_\mu(U_n)=\frac{4\mu^2\sigma^2}{n}+\frac{2\sigma^4}{n^2}.
\]
它未达到有限样本下界。在 $\mu=0$，下界甚至为零。这并不证明存在方差更小的无偏估计，第四章的最优无偏理论还需单独判断。
\end{examplebox}

\subsection{有偏估计的修正}
若 $\E_\theta T=g(\theta)+b(\theta)$，$b$ 可微且同样允许求导交换，则
\[
 \Cov_\theta(T,S_n)=g'(\theta)+b'(\theta).
\]
因此
\[
 \Var_\theta(T)\ge\frac{\{g'(\theta)+b'(\theta)\}^2}{I_n(\theta)},
\]
\[
 \operatorname{MSE}_\theta(T)
 \ge b(\theta)^2+
 \frac{\{g'(\theta)+b'(\theta)\}^2}{I_n(\theta)}.
\]
例如 $X\sim N(\theta,1)$、$T=aX$，$g'=1$、$b'=a-1$，修正后的方差下界是 $a^2$，恰与实际方差相等。不能把无偏下界 $1$ 用于这个有偏估计。

\paragraph{复习题}
\begin{enumerate}
 \item 指数尺度模型中，证明 $\bar X$ 达到估计 $\theta$ 的信息下界。
 \item 正态均值模型下，常数估计 $T=m$ 方差为零，为何不违反信息下界？用有偏版本核对。
\end{enumerate}
\begin{notebox}[答案线索]
第一题：$I_n=n/\theta^2$，$S_n=n(\bar X-\theta)/\theta^2$，等号条件直接成立。第二题：常数规则并非在整个模型中无偏，$b(\theta)=m-\theta$，故 $g'+b'=0$，有偏下界给出 $\operatorname{MSE}\ge(m-\theta)^2$，正好取等。
\end{notebox}

\section{向量参数与矩阵信息下界}

设 $\theta\in\R^d$，目标 $g(\theta)\in\R^k$，$T\in\R^k$ 无偏，Jacobian 约定为
\[
 Dg(\theta)=\left(\frac{\partial g_j}{\partial\theta_r}\right)_{j,r}
 \in\R^{k\times d}.
\]
在二阶矩存在、均值逐分量可交换微分积分、得分正则及 $I_n(\theta)$ 正定的条件下，
\[
 \Cov_\theta(T,S_n)=Dg(\theta),
\]
\begin{definitionbox}{矩阵 Cramér--Rao 不等式}
\[
 \Cov_\theta(T)\succeq
 Dg(\theta)I_n(\theta)^{-1}Dg(\theta)^\top.
\]
$A\succeq B$ 指 $A-B$ 半正定，即 $u^\top(A-B)u\ge0$ 对所有 $u$ 成立，并非逐元素大小。估计参数本身时，下界为 $I_n^{-1}$。
\end{definitionbox}

\paragraph{证明}
记 $E=T-g(\theta)$、$A=Dg(\theta)$、$B=I_n(\theta)$。则 $\Cov(E,S_n)=A$、$\Cov(S_n)=B$。令
\[
 U=E-AB^{-1}S_n,
\]
得到 $\E U=0$、$\Cov(U,S_n)=A-AB^{-1}B=0$。因此
\[
 \Cov(T)=AB^{-1}A^\top+\E(UU^\top).
\]
末项半正定，得到结论。矩阵等号成立当且仅当 $\E\|U\|^2=0$，即
\[
 T-g(\theta)=Dg(\theta)I_n(\theta)^{-1}S_n(\theta)
 \quad P_\theta\text{-a.s.}
\]
这个证明也给出维度核对：$(k\times d)(d\times d)(d\times k)=k\times k$。

\begin{examplebox}{正态均值与未知方差}
以 $\theta=(\mu,v)^\top$ 参数化 $N(\mu,v)$，$v>0$。单次得分为
\[
 s_\theta(X)=
 \begin{pmatrix}
 (X-\mu)/v\\[2pt]
 -1/(2v)+(X-\mu)^2/(2v^2)
 \end{pmatrix},
 \qquad
 I(\theta)=
 \begin{pmatrix}1/v&0\\0&1/(2v^2)\end{pmatrix}.
\]
交叉信息为零来自中心正态的三阶矩为零；右下项来自四阶矩 $3v^2$。估计 $\mu$ 时 $Dg=(1,0)$，下界仍为 $v/n$，由 $\bar X$ 达到。

估计 $v$ 时下界为 $2v^2/n$。无偏样本方差 $S^2$ 的方差为 $2v^2/(n-1)$（$n\ge2$），故未达到此界。方差参数的下界应由整个信息矩阵的逆取得，不能一般地直接用信息矩阵对角元的倒数。
\end{examplebox}

若把 $\theta=(\psi,\lambda)$ 分为目标与干扰参数，对标量目标 $\psi$，分块信息在 $I_{\lambda\lambda}$ 可逆时给出
\[
 (I^{-1})_{\psi\psi}
 =\{I_{\psi\psi}-I_{\psi\lambda}
 I_{\lambda\lambda}^{-1}I_{\lambda\psi}\}^{-1}.
\]
这解释了为何未知干扰参数一般会减少有效信息。正态例中的交叉块为零是一个特殊情形。

\paragraph{复习题}
\begin{enumerate}
 \item 对目标 $g(\theta)=c^\top\theta$，写出标量方差下界。为什么不能逐元素解释矩阵不等式？
 \item 在正态例中，改用参数 $(\mu,\sigma)$，其中 $v=\sigma^2$，求新信息矩阵。
\end{enumerate}
\begin{notebox}[答案线索]
第一题：下界为 $c^\top I(\theta)^{-1}c/n$；矩阵偏序同时约束全部线性方向的方差。第二题：旧参数对新参数的 Jacobian 是 $\operatorname{diag}(1,2\sigma)$，故新信息矩阵为 $\operatorname{diag}(1/\sigma^2,2/\sigma^2)$。
\end{notebox}

\section{达到下界与指数族结构}

\subsection{自然统计量达到哪个目标的下界}
固定样本量，将\emph{整份样本}密度写作
\[
 p_\eta^{(n)}(x)=h_n(x)\exp\{\eta T(x)-\zeta(\eta)\}.
\]
假设自然参数在归一化积分有限的开区间内变化，且 $\zeta''(\eta)>0$。由归一化函数的导数，
\[
 \E_\eta T=\zeta'(\eta),\qquad
 \Var_\eta(T)=I_n(\eta)=\zeta''(\eta).
\]
对目标 $g(\eta)=\zeta'(\eta)$，信息下界恰为
\[
 \frac{g'(\eta)^2}{I_n(\eta)}=\zeta''(\eta)=\Var_\eta(T).
\]
所以 $T$ 达到其均值目标的下界。这里的 $\zeta$ 已是联合模型的归一化函数，不能再给信息乘一个 $n$。目标是 $\zeta'(\eta)$，并不自动是自然参数 $\eta$ 本身。

\begin{examplebox}{Bernoulli 的自然参数}
令 $\eta=\log\{p/(1-p)\}$，$Y=\sum_iX_i$，则
\[
 p_\eta^{(n)}(x)=\exp\{\eta Y-n\log(1+e^\eta)\}.
\]
$Y$ 的均值为 $np$，方差与样本信息均为 $np(1-p)$；$Y$ 达到估计 $np$ 的下界。相应地，$Y/n$ 达到估计 $p$ 的下界，但这没有提供自然参数 $\eta$ 的无偏有效估计。
\end{examplebox}

\subsection{处处取等反推一维指数族}
\begin{definitionbox}{CR 等号的逆向结论}
固定样本量，设标量参数在开区间 $\Theta$ 内变化，模型满足标量 CR 条件；密度在共同支撑上严格为正，除同一个固定测度零集外关于参数连续可微。另设 $I_n(\theta)>0$ 与 $g'(\theta)\ne0$ 连续。

若同一个无偏统计量 $T$ 在每个 $\theta\in\Theta$ 均达到下界，则在一个共同满测度样本集合上，
\[
 p_\theta^{(n)}(x)
 =h_n(x)\exp\{\eta(\theta)T(x)-\zeta(\eta(\theta))\},
\]
其中 $\eta'(\theta)=I_n(\theta)/g'(\theta)\ne0$，并有
$g(\theta)=\zeta'(\eta(\theta))$。
\end{definitionbox}

\paragraph{证明}
令 $k(\theta)=I_n(\theta)/g'(\theta)$。CR 等号给出
\[
 \partial_\theta\log p_\theta^{(n)}(x)
 =k(\theta)\{T(x)-g(\theta)\}
 \quad P_\theta\text{-a.s.}
\]
不能直接固定 $x$ 对整个参数区间积分，因为此时异常零集可能随参数改变。

取可数稠密集 $D\subset\Theta$。共同支撑上的密度严格为正，所以每个 $q\in D$ 下的概率零集也是支撑内的控制测度零集。把这些可数零集及密度连续可微的固定异常集一并剔除，得到共同满测度集合 $F$。对 $x\in F$，上式在所有 $q\in D$ 成立；由两边关于参数连续，遂在全部 $\Theta$ 成立。

固定参考点 $\theta_*$ 与 $x\in F$，对参数积分：
\[
 \log p_\theta^{(n)}(x)-\log p_{\theta_*}^{(n)}(x)
 =\left(\int_{\theta_*}^\theta k(t)\,\mathrm dt\right)T(x)
 -\int_{\theta_*}^\theta k(t)g(t)\,\mathrm dt.
\]
定义 $\eta(\theta)=\int_{\theta_*}^\theta k(t)\,\mathrm dt$、
$c(\theta)=\int_{\theta_*}^\theta k(t)g(t)\,\mathrm dt$，取
$h_n=p_{\theta_*}^{(n)}$，便得到指数形式。$k$ 连续非零，在区间上不能变号，故 $\eta$ 严格单调，其像为开区间。归一化给出
\[
 c(\theta)=\log\int h_n(x)e^{\eta(\theta)T(x)}\,\mathrm d\nu^n(x)
 =\zeta(\eta(\theta)).
\]
最后由 $c'=kg$、$\eta'=k\ne0$ 得 $\zeta'(\eta(\theta))=g(\theta)$。

\subsection{一点取等不足以推断整族结构}
在样本空间 $\{-1,0,1\}$ 上，设
\[
 p_\theta(x)=\frac13\left[1+\theta x+
 \theta^2\left(x^2-\frac23\right)\right],\qquad |\theta|<\frac12.
\]
三点概率均为正且总和为 $1$。取 $T=X$，则 $g(\theta)=2\theta/3$。在 $\theta=0$，得分为 $X$，$I(0)=\Var_0(X)=2/3$，故
\[
 \frac{g'(0)^2}{I(0)}=\frac23=\Var_0(X).
\]
但若该族能以同一个 $T(x)=x$ 写成一维指数形式，则
$p_\theta(1)p_\theta(-1)/p_\theta(0)^2$ 必为常数，而实际等于
\[
 \frac{1-\theta^2/3+\theta^4/9}{1-4\theta^2/3+4\theta^4/9},
\]
会随 $\theta$ 改变。这个反例否定的是“一点取等推出以 $X$ 为自然统计量的整族结构”，不是否定模型嵌入更高维指数族的可能性。

\paragraph{复习题}
\begin{enumerate}
 \item 逆向证明中，为什么不能直接合并全部参数点的异常零集？可数稠密集之后为何还要连续性？
 \item 联合 Poisson 模型中，$Y=\sum_iX_i$ 的自然参数归一化函数为 $\zeta(\eta)=ne^\eta$。说明 $Y$ 和 $Y/n$ 分别达到哪个目标的下界。
\end{enumerate}
\begin{notebox}[答案线索]
第一题：不可数个零集的并未必为零集；可数稠密集先统一样本异常集，连续性再把等式推广至整个区间，才能固定样本点积分。第二题：$Y$ 对 $n\lambda$ 有效，$Y/n$ 对 $\lambda$ 有效，方差分别为 $n\lambda$ 与 $\lambda/n$。
\end{notebox}

\section{非正则模型、效率与 MLE 展开}

\subsection{移动支撑为什么使正则恒等式失效}
\begin{examplebox}{均匀模型的边界项}
设 $X_i\sim U(0,\theta)$，$\theta>0$。只在内部 $0<x<\theta$ 对 $\log p_\theta(x)=-\log\theta$ 求导，会得到形式上的内部得分 $-1/\theta$，其均值不为零。实际上
\[
 \frac{\mathrm d}{\mathrm d\theta}\int_0^\theta\frac1\theta\,\mathrm dx
 =\int_0^\theta-\frac1{\theta^2}\,\mathrm dx+\frac1\theta=0.
\]
移动端点贡献的第二项不能省略。因此内部导数的平方均值 $1/\theta^2$ 不能直接代入正则 CR 公式。
\end{examplebox}

记 $M_n=\max_iX_i$。在 $0<m<\theta$ 上，
\[
 P_\theta(M_n\le m)=(m/\theta)^n,\qquad
 \E_\theta M_n^r=\frac n{n+r}\theta^r\quad(r>0).
\]
故 $T_n=(n+1)M_n/n$ 无偏，且
\[
 \Var_\theta(T_n)=\frac{\theta^2}{n(n+2)}.
\]
这比误用正则公式得到的 $\theta^2/n$ 小，并不构成对定理的反例。模型的正则条件已经失效。

它的自然误差尺度是 $n$。对固定 $z\ge0$，
\[
 P_\theta\{n(1-M_n/\theta)>z\}
 =(1-z/n)^n\to e^{-z}.
\]
因此 $n(1-M_n/\theta)\dto E$、$E\sim\operatorname{Exp}(1)$，并有
\[
 n(T_n-\theta)\dto\theta(1-E).
\]
更快速度来自支撑边界携带的信息；这里的非正态极限不应混同于正则模型中孤立参数点的超效率现象。

\subsection{有限样本效率与渐近效率}
无偏估计量达到正确的 CR 协方差界，称达到有限样本信息效率。这足以保证它在相应无偏类中无法被方差统一改进；未达到下界则不自动意味着可以改进，因为下界可能不可达。

正则模型中，估计可微目标 $g(\theta)$ 的一阶信息基准为
\[
 \sqrt n\{T_n-g(\theta)\}
 \dto N_k\!\left(0,Dg(\theta)I(\theta)^{-1}Dg(\theta)^\top\right).
\]
渐近效率的比较涉及适当的局部正则估计量类，不要求每个有限样本量都精确无偏。不能把有限样本无偏下界直接令 $n\to\infty$，便声称对一切估计量序列成立；弱收敛不自动给出风险收敛，偏差和局部行为也须控制。

\subsection{从得分方程到渐近正态}
设真参数为 $\theta_0$，$\widehat\theta_n$ 是内部 MLE 或一个选定的得分根，并假设 $S_n(\widehat\theta_n)=0$。多维时使用积分 Hessian，避免给所有分量强设同一个中间点：
\[
 \widetilde H_n=
 \int_0^1H_n\{\theta_0+t(\widehat\theta_n-\theta_0)\}\,\mathrm dt.
\]
在线段位于可微参数邻域的事件上，微积分基本定理给出精确等式
\[
 0=S_n(\theta_0)+\widetilde H_n(\widehat\theta_n-\theta_0).
\]
若沿线段矩阵最终可逆，则
\[
 \sqrt n(\widehat\theta_n-\theta_0)
 =\left(-\frac{\widetilde H_n}{n}\right)^{-1}
 \frac{S_n(\theta_0)}{\sqrt n}.
\]
得分的多元中心极限定理给出
\[
 S_n(\theta_0)/\sqrt n\dto N_d(0,I(\theta_0)).
\]
若另有 $-\widetilde H_n/n\pto I(\theta_0)$，且 $I(\theta_0)$ 正定，则由矩阵求逆的连续性与 Slutsky 定理，
\[
 \sqrt n(\widehat\theta_n-\theta_0)
 \dto N_d(0,I(\theta_0)^{-1}).
\]

\begin{notebox}[展开之外必须补的条件]
固定点的大数定律 $-H_n(\theta_0)/n\pto I(\theta_0)$ 尚不足以控制随机线段上的 Hessian。一组可用的充分条件是：所选根存在且一致；真值附近有凸邻域；该邻域内 $-H_n(v)/n$ 一致收敛到连续函数 $J_{\theta_0}(v)$，并满足 $J_{\theta_0}(\theta_0)=I(\theta_0)$ 正定。于是随机线段收缩至真值，得到所需收敛。

MLE 是否存在、所选根是否一致、最终是否位于内部、多个根如何选择，都不是得分方程本身保证的。边界极值也未必满足得分为零。本段给出的是在这些条件落实后的证明，不是无条件的一般 MLE 定理。
\end{notebox}

\begin{examplebox}{正态均值中的展开是精确恒等式}
已知方差时，$S_n(\mu)=n(\bar X-\mu)/\sigma^2$，$H_n=-n/\sigma^2$ 为常数，MLE 为 $\bar X$。于是
\[
 \sqrt n(\bar X-\mu)
 =\sigma^2\frac{S_n(\mu)}{\sqrt n}
 \sim N(0,\sigma^2).
\]
这里没有 Taylor 余项，且极限方差 $I(\mu)^{-1}=\sigma^2$ 同时也是精确分布的方差。
\end{examplebox}

\paragraph{复习题}
\begin{enumerate}
 \item 在均匀模型中，证明 $n(T_n-\theta)=O_{P_\theta}(1)$，并说明它为何不违背正则模型的信息下界。
 \item MLE 渐近正态性的证明中，分别指出使用中心极限定理、大数型控制及 Slutsky 定理的位置。固定点 Hessian 收敛还缺少什么？
\end{enumerate}
\begin{notebox}[答案线索]
第一题：$n^2\Var(T_n)=n\theta^2/(n+2)\le\theta^2$，由 Chebyshev 即得随机有界；该模型支撑依赖参数，正则得分恒等式失效。第二题：CLT 用于真值处的得分，局部一致控制用于随机线段上的 Hessian，Slutsky 用于两者的乘积；仅有固定点收敛不能替代随机邻域控制。
\end{notebox}

\chapter{最优无偏估计与 U 统计量}

本章的前半部分在全部无偏估计中比较方差。充分性使条件化后的规则能够由数据计算，完备性把充分统计量的无偏函数确定下来。后半部分研究另一种构造：若总体量可以写成若干独立观测的函数的期望，就把同一个函数应用于所有同样大小的子样本，再取平均。所得 U 统计量通常含有相互依赖的核项，其方差与极限分布须由重叠结构和条件投影分析。

除特别说明外，估计量均为不含未知参数的可测数据函数；涉及方差的比较，均要求在每个参数分布下二阶矩有限。记 $\E_\theta$、$\Var_\theta$ 为参数 $\theta$ 下的期望与方差。

\section{条件化与一致最小方差无偏估计}

\subsection{UMVUE 的比较范围}

\begin{definitionbox}{一致最小方差无偏估计}
设目标为 $g(\theta)$。若同一个统计量 $U$ 满足
\[
 \E_\theta U=g(\theta),\qquad
 \Var_\theta(U)\leq\Var_\theta(W)
\]
对每个 $\theta\in\Theta$ 及每个有限方差无偏估计 $W$ 均成立，则称 $U$ 为 $g(\theta)$ 的一致最小方差无偏估计，简称 UMVUE。
\end{definitionbox}

这里的“一致”指同一规则在全部参数点上最优，并非一致收敛。统计量不能用到未知参数，因此“对每个 $\theta$ 单独找一个最优规则”不等于找到了 UMVUE。无偏条件下，平方误差风险为
\[
 \E_\theta\{[U-g(\theta)]^2\}=\Var_\theta(U).
\]
因此方差最小等价于在无偏类中均方误差最小；该结论没有把有偏估计纳入比较。

\subsection{条件期望的正交性质}

先固定一个概率分布。对 $\E U^2<\infty$，令
\[
 H=\E(U\mid T),\qquad R=U-H.
\]
塔式法则给出 $\E H=\E U$；条件 Jensen 不等式给出 $\E H^2\leq\E U^2$。又因 $\E(R\mid T)=0$，对每个二阶可积的 $a(T)$，
\begin{align*}
 \E[Ra(T)]
 &=\E\{\E[Ra(T)\mid T]\}\\
 &=\E\{a(T)\E(R\mid T)\}=0.
\end{align*}
乘积的可积性由 Cauchy--Schwarz 不等式保证。取 $a(T)=H-\E U$，有
\begin{equation}
 \Var(U)=\Var(H)+\E R^2
 =\Var\{\E(U\mid T)\}+\E\{\Var(U\mid T)\}.
 \label{eq:ch4-total-var}
\end{equation}
更一般地，对任意二阶可积的 $b(T)$，
\[
 \E\{[U-b(T)]^2\}=\E(U-H)^2+\E\{[H-b(T)]^2\}.
\]
所以条件期望是所有 $T$ 的函数中对 $U$ 最好的均方逼近。残差与这些函数正交；正交并不要求独立。

\subsection{Rao--Blackwell 定理及共同版本}

\begin{definitionbox}{Rao--Blackwell 定理}
设 $T$ 充分，$\E_\theta U^2<\infty$ 对所有参数成立。在通常的支配统计模型中，可以选取同一个参数无关的可测函数 $u^*$，使
\[
 U^*=u^*(T)=\E_\theta(U\mid T)
 \quad P_\theta\text{-几乎处处},\quad\forall\theta.
\]
它满足
\[
 \E_\theta U^*=\E_\theta U,\qquad
 \Var_\theta(U^*)\leq\Var_\theta(U).
\]
对指定的 $\theta$，方差取等当且仅当 $U=U^*$，$P_\theta$-几乎处处。
\end{definitionbox}

充分性的作用先于方差比较。一般的 $\E_\theta(U\mid T=t)$ 可能含未知参数，不能直接作为估计规则；充分性保证可选择整个模型共用的条件期望版本。在标准 Borel 样本空间上，用不依赖 $\theta$ 的样本条件分布核 $K(t,dx)$ 表示这一点，即
\[
 u^*(t)=\int U(x)K(t,dx).
\]
对积分未定义或不有限的 $t$ 可统一赋予任一固定值；由于 $U$ 在每个参数下可积，这些例外在每个 $P_\theta^T$ 下均为零测集。如此得到的是同一个数据函数，而非为每个参数分别选择的公式。在不直接使用正则条件分布的表述中，只需使用充分性的共同条件期望版本。

固定任意参数，由塔式法则有 $\E_\theta U^*=\E_\theta U$。条件 Jensen 不等式保证 $U^*$ 二阶可积，再由式~\eqref{eq:ch4-total-var}，
\[
 \Var_\theta(U)-\Var_\theta(U^*)
 =\E_\theta(U-U^*)^2\geq0.
\]
差为零等价于残差几乎处处为零，定理得证。定理本身不要求 $U$ 无偏，只保持其期望；若原估计无偏，条件化后才随之无偏。若要断言每个参数点都严格改进，还须逐一证明 $P_\theta(U\neq U^*)>0$。

\begin{examplebox}{Poisson 尾概率的条件化}
设 $X_i\overset{\mathrm{i.i.d.}}\sim\operatorname{Poisson}(\lambda)$，$\lambda>0$，$n\geq2$，欲估计 $e^{-\lambda}$。取
\[
 U=\ind_{\{X_1=0\}},\qquad T=\sum_{i=1}^nX_i.
\]
联合概率质量函数为 $e^{-n\lambda}\lambda^{\sum x_i}/\prod x_i!$，故 $T$ 充分。利用独立的 $X_1$ 与 $T-X_1$，对 $0\leq j\leq t$，
\[
 P(X_1=j\mid T=t)
 =\binom tj\left(\frac1n\right)^j
       \left(1-\frac1n\right)^{t-j}.
\]
所以 $X_1\mid T=t\sim\operatorname{Bin}(t,1/n)$，
\[
 U^*=\left(1-\frac1n\right)^T.
\]
Poisson 概率母函数 $\E s^T=\exp\{n\lambda(s-1)\}$ 给出
\[
 \E U^*=e^{-\lambda},\qquad
 \Var(U^*)=e^{-2\lambda+\lambda/n}-e^{-2\lambda}.
\]
原指标函数的方差为 $e^{-\lambda}-e^{-2\lambda}$，故这里确有严格改进。
\end{examplebox}

\subsection{Lehmann--Scheff\'e 定理}

\begin{definitionbox}{完备充分统计量给出的 UMVUE}
设 $T$ 完备且充分，$g(\theta)$ 存在每个参数下均有有限方差的无偏估计。则存在 $T$ 的无偏函数 $h(T)$，它是 $g(\theta)$ 的唯一 UMVUE。

更准确地，$T$ 的任何可积无偏函数都与 $h(T)$ 在每个 $P_\theta$ 下几乎处处相等；所有估计量中的任何另一个 UMVUE 也具有同样的几乎处处唯一性。
\end{definitionbox}

\textbf{证明。} 从任一有限方差无偏估计 $U_0$ 出发，令 $h(T)=\E(U_0\mid T)$。Rao--Blackwell 定理保证这是参数无关、二阶可积的无偏统计量。对任意无偏竞争者 $W$，同样令 $h_W(T)=\E(W\mid T)$。两者均估计 $g(\theta)$，因此
\[
 \E_\theta\{h_W(T)-h(T)\}=0\qquad\forall\theta.
\]
差是 $T$ 的可积函数，完备性给出 $h_W(T)=h(T)$，对每个参数几乎处处成立。于是
\[
 \Var_\theta\{h(T)\}
 =\Var_\theta\{h_W(T)\}
 \leq\Var_\theta(W).
\]
由于 $W$ 任意，$h(T)$ 为 UMVUE。若 $W$ 本身也是 UMVUE，则上述方差取等。由 Rao--Blackwell 的等号条件，$W=h_W(T)=h(T)$，$P_\theta$-几乎处处。

最后，若 $\widetilde h(T)$ 只是可积无偏函数，则其与 $h(T)$ 的差仍满足完备性的条件，因而同样相等；不必预先对 $\widetilde h$ 另加二阶矩条件。\hfill$\square$

\begin{notebox}[唯一性的范围]
“$T$ 的无偏函数唯一”不等于“所有无偏估计唯一”。例如 Poisson 例中的指标函数和条件化结果都无偏，但只有后者是 $T$ 的函数。唯一性均按每个 $P_\theta$ 几乎处处理解，不是说公式在样本空间每一点都必须相同，也不自动承诺存在一个适用于不可数参数集的共同满测度集合。
\end{notebox}

\section{UMVUE 的几种构造}

已有完备充分统计量时，一种办法是把简单无偏估计条件化，另一种是直接解无偏方程 $\E_\theta h(T)=g(\theta)$。前者多用于事件概率，后者常用积分求导、幂级数比较或矩的偏差修正。两种办法都须核对可积性与有限方差。

\subsection{均匀分布：积分方程与端点条件}

设 $X_i\overset{\mathrm{i.i.d.}}\sim U(0,\theta)$，$\theta>0$，令 $T=X_{(n)}$。联合密度的参数相关部分为 $\theta^{-n}\ind_{\{T<\theta\}}$，故 $T$ 充分。其分布为
\[
 P_\theta(T\leq t)=\left(\frac t\theta\right)^n,
 \qquad
 f_T(t;\theta)=\frac{nt^{n-1}}{\theta^n}\ind_{(0,\theta)}(t).
\]
若可积函数 $v(T)$ 在全部参数下期望为零，则
\[
 \int_0^\theta v(t)t^{n-1}\,dt=0\qquad\forall\theta>0.
\]
被积函数在每个 $(0,b)$ 上可积，故左侧作为上限的函数绝对连续，导数几乎处处为 $v(\theta)\theta^{n-1}$。因此 $v(t)=0$ 在 $(0,\infty)$ 上 Lebesgue-几乎处处，各个 $P_\theta^T$ 都绝对连续，遂有 $v(T)=0$，$P_\theta$-几乎处处。这里不需要假设 $v$ 连续。

由
\[
 \E_\theta T^r=\frac n{n+r}\theta^r\qquad(r>-n),
\]
可知对 $r>0$，
\[
 \widehat{\theta^r}=\frac{n+r}{n}T^r,
 \qquad
 \Var_\theta(\widehat{\theta^r})
 =\frac{r^2\theta^{2r}}{n(n+2r)}.
\]
它是完备充分统计量的有限方差无偏函数，故为 UMVUE。特别地，$\widehat\theta=(n+1)T/n$。

对一般连续可微目标 $g$，无偏方程为
\[
 n\int_0^\theta h(t)t^{n-1}\,dt=\theta^ng(\theta).
\]
求导得到候选公式
\begin{equation}
 h(t)=g(t)+\frac tn g'(t).
 \label{eq:ch4-uniform-target}
\end{equation}
但求导会丢失积分常数。对 $0<\varepsilon<\theta$，代回候选函数可得
\[
 n\int_\varepsilon^\theta h(t)t^{n-1}\,dt
 =\theta^ng(\theta)-\varepsilon^ng(\varepsilon).
\]
所以若 $t^ng(t)\to0$（$t\downarrow0$），且 $\E_\theta h(T)^2<\infty$ 对每个参数成立，才能令 $\varepsilon\downarrow0$ 得到所需无偏性。

\begin{examplebox}{不能省略的端点核验}
若把 $g(\theta)=\theta^{-n}$ 机械地代入式~\eqref{eq:ch4-uniform-target}，所得 $h$ 恒为零，显然不无偏估计正的 $\theta^{-n}$。原因是 $t^ng(t)=1$，没有满足零端点条件。候选公式来自微分方程，并不自动满足原积分方程。
\end{examplebox}

\subsection{Poisson 分布：下降阶乘与幂级数}

设 $X_i\sim\operatorname{Poisson}(\theta)$ 独立同分布，$T=\sum_iX_i\sim\operatorname{Poisson}(n\theta)$。自然参数 $\log\theta$ 遍历 $\R$，所以 $T$ 完备且充分。记下降阶乘
\[
 (t)_{\underline r}=t(t-1)\cdots(t-r+1),
 \quad (t)_{\underline0}=1,
\]
当非负整数 $t<r$ 时令其为零。直接计算可得
\begin{align*}
 \E_\theta\frac{(T)_{\underline r}}{n^r}
 &=\sum_{t=r}^\infty
 \frac{t!}{n^r(t-r)!}e^{-n\theta}\frac{(n\theta)^t}{t!}\\
 &=\theta^r e^{-n\theta}
   \sum_{u=0}^\infty\frac{(n\theta)^u}{u!}
 =\theta^r.
\end{align*}
该统计量是 $T$ 的多项式，Poisson 的所有整数阶矩有限。因此
\[
 \widehat{\theta^r}=\frac{(T)_{\underline r}}{n^r}
\]
为 UMVUE。例如 $T(T-1)/n^2$ 无偏估计 $\theta^2$，而 $(T/n)^2$ 的期望是 $\theta^2+\theta/n$。

若目标有在全部正实数处收敛的幂级数 $g(\theta)=\sum_{j\geq0}a_j\theta^j$，无偏方程乘以 $e^{n\theta}$ 后，比较 $\theta^t$ 的系数，得到
\[
 h(t)=\frac{t!}{n^t}\sum_{j=0}^t\frac{n^{t-j}a_j}{(t-j)!}
 =\sum_{j=0}^t a_j\frac{(t)_{\underline j}}{n^j}.
\]
这里收敛半径为无穷，故 $\sum_j|a_j|\theta^j<\infty$。由 Tonelli 及阶乘矩公式，
\[
 \E_\theta\sum_{j\geq0}|a_j|\frac{(T)_{\underline j}}{n^j}
 =\sum_{j\geq0}|a_j|\theta^j<\infty,
\]
所以允许交换求和与期望，证得无偏性。还须另验 $\E_\theta h(T)^2<\infty$，才能按本章的有限方差定义称其为 UMVUE；幂级数绝对收敛不能替代二阶矩条件。

\subsection{正态分布：偏差修正}

设 $X_i\overset{\mathrm{i.i.d.}}\sim N(\mu,\sigma^2)$，$n\geq2$，$\mu\in\R$、$\sigma^2>0$ 均未知。自然统计量与自然参数分别为
\[
 \left(\sum_iX_i,\sum_iX_i^2\right),\qquad
 \left(\frac\mu{\sigma^2},-\frac1{2\sigma^2}\right).
\]
自然参数遍历开集 $\R\times(-\infty,0)$，从而统计量完备且充分。令
\[
 \bar X=\frac1n\sum_iX_i,\qquad
 S^2=\frac1{n-1}\sum_i(X_i-\bar X)^2.
\]
两组统计量相互确定，因为 $\sum_iX_i^2=(n-1)S^2+n\bar X^2$，故 $(\bar X,S^2)$ 也完备且充分。

由平方和分解，
\[
 \sum_i(X_i-\bar X)^2
 =\sum_i(X_i-\mu)^2-n(\bar X-\mu)^2,
\]
取期望得 $\E S^2=\sigma^2$。又 $\E\bar X^2=\mu^2+\sigma^2/n$，因此
\[
 \widehat\mu=\bar X,\qquad
 \widehat{\sigma^2}=S^2,\qquad
 \widehat{\mu^2}=\bar X^2-\frac{S^2}{n}
\]
均为对应目标的 UMVUE。正态分布的矩保证其方差有限。最后一个估计可能为负，无偏性约束的是期望，并不强迫每次输出都落在目标值域内。

\subsection{指数分布：条件概率}

设 $X_i\overset{\mathrm{i.i.d.}}\sim\operatorname{Exp}(\theta)$，这里 $\theta>0$ 为尺度参数，密度为 $\theta^{-1}e^{-x/\theta}\ind_{\{x>0\}}$。固定 $a>0$，欲估计 $e^{-a/\theta}=P_\theta(X_1>a)$。对 $n\geq2$，取
\[
 U=\ind_{\{X_1>a\}},\qquad T=\sum_iX_i.
\]
$T\sim\operatorname{Gamma}(n,\text{scale }\theta)$，其自然参数 $-1/\theta$ 遍历 $(-\infty,0)$，故 $T$ 完备充分。令 $Y=T-X_1$，它与 $X_1$ 独立，并服从形状 $n-1$ 的同尺度 Gamma 分布。由联合密度除以边缘密度，
\begin{align*}
 f_{X_1\mid T}(x\mid t)
 &=\frac{f_{X_1}(x)f_Y(t-x)}{f_T(t)}\\
 &=\frac{n-1}{t}\left(1-\frac xt\right)^{n-2},
 \qquad0<x<t.
\end{align*}
因此 $X_1/T\mid T=t\sim\operatorname{Beta}(1,n-1)$，且
\[
 h(T)=\E(U\mid T)=
 \begin{cases}
  0,&T\leq a,\\
  (1-a/T)^{n-1},&T>a.
 \end{cases}
\]
第二种情形来自对条件密度在 $(a,t)$ 上积分。塔式法则给出无偏性，$0\leq h\leq1$ 保证有限方差，Lehmann--Scheff\'e 定理遂给出唯一 UMVUE。

\begin{notebox}[练习与核验]
\begin{enumerate}
 \item 在 $U(0,\theta)$ 模型中求 $\theta^{-r}$ 的 UMVUE，并说明有限方差条件。\textbf{线索：}用 $\E T^s=n\theta^s/(n+s)$；当 $0<r<n/2$ 时，$(n-r)T^{-r}/n$ 有有限方差。
 \item 在 Poisson 模型中无偏估计 $\theta^3$。\textbf{答案：}$T(T-1)(T-2)/n^3$；比较它与 $\bar X^3$ 的期望。
 \item 指数分布例中若 $n=1$，原指标函数是否为 UMVUE？\textbf{线索：}$X_1$ 本身完备充分；上述 Beta 条件密度的推导不适用于 $n=1$，但结论可直接由 Lehmann--Scheff\'e 得出。
\end{enumerate}
\end{notebox}

\section{正交判别与参数空间的影响}

\subsection{与所有零均值方向正交}

定义
\[
 \mathcal U_0=\{V(X):\E_\theta V=0,
 \E_\theta V^2<\infty\text{ 对所有 }\theta\in\Theta\}.
\]
这里 $V$ 必须是固定的数据函数。虽然 $X_1-\theta$ 在位置模型中可能有零均值，它含未知参数，不能据此列入 $\mathcal U_0$。若 $U,W$ 是同一目标的无偏估计，则 $W-U\in\mathcal U_0$；反之，$U+V$ 对每个 $V\in\mathcal U_0$ 都是无偏竞争者。

\begin{definitionbox}{UMVUE 的正交判别}
设 $U$ 为有限方差无偏估计，则
\[
 U\text{ 为 UMVUE}
 \quad\Longleftrightarrow\quad
 \E_\theta(UV)=0
 \quad\forall V\in\mathcal U_0,\ \forall\theta\in\Theta.
\]
由于 $\E_\theta V=0$，右侧等价于 $\Cov_\theta(U,V)=0$。该判别不要求存在完备充分统计量。
\end{definitionbox}

\textbf{证明。} 若 $U$ 为 UMVUE，则对任意实数 $c$，$U+cV$ 仍无偏，因此
\[
 0\leq\Var_\theta(U+cV)-\Var_\theta(U)
 =c^2A+2cB,
\]
其中 $A=\Var_\theta(V)$、$B=\Cov_\theta(U,V)$。若 $A>0$ 而 $B\neq0$，取 $c=-B/A$ 得矛盾；若 $A=0$，则 $V=0$，$P_\theta$-几乎处处，亦有 $B=0$。

反之，任一竞争者可写为 $W=U+V$，$V\in\mathcal U_0$。正交性给出
\[
 \Var_\theta(W)=\Var_\theta(U)+\Var_\theta(V)
 \geq\Var_\theta(U),
\]
故 $U$ 为 UMVUE。\hfill$\square$

若 $U_j$ 分别为 $g_j(\theta)$ 的 UMVUE，则对固定常数 $c_j$，$\sum_jc_jU_j$ 与所有零均值方向仍正交，因而是 $\sum_jc_jg_j(\theta)$ 的 UMVUE。这一性质只保证有限线性组合，不保证平方或其他非线性变换。若 $U_1,U_2$ 为同一目标的 UMVUE，取 $V=U_1-U_2$，则
\[
 \E_\theta V^2=\E_\theta(U_1V)-\E_\theta(U_2V)=0,
\]
再次得到 UMVUE 的几乎处处唯一性。

若已有充分统计量 $T$，判断其无偏函数 $h(T)$ 是否为 UMVUE，只需检验
\[
 \E_\theta[h(T)v(T)]=0
\]
对所有零均值、二阶可积的 $v(T)$ 成立。事实上，对一般 $V\in\mathcal U_0$，令 $v(T)=\E(V\mid T)$。充分性、塔式法则和条件 Jensen 保证它属于所检验的类，且
\[
 \E_\theta[h(T)V]
 =\E_\theta[h(T)\E(V\mid T)]
 =\E_\theta[h(T)v(T)]=0.
\]

\subsection{限制均匀模型中的改进}

仍考虑 $U(0,\theta)$，现在参数只允许 $\theta\geq1$。最大值 $T$ 的密度不变，也仍充分。但 $\E_\theta v(T)=0$ 对全部 $\theta\geq1$ 成立，仅能推出
\begin{equation}
 \int_0^1v(t)t^{n-1}\,dt=0,
 \qquad v(t)=0\ \text{在 }(1,\infty)\text{ 上几乎处处}.
 \label{eq:ch4-restricted-zero}
\end{equation}
第一项由 $\theta=1$ 得到，第二项由积分上限 $\theta>1$ 的几乎处处求导得到；两条件也反过来保证零均值。区间 $(0,1)$ 上只剩一个积分约束，因此 $T$ 不再完备。例如
\[
 v(t)=\left(\frac{n+1}{n}t-1\right)\ind_{\{t\leq1\}}
\]
是非零的有界函数，却满足式~\eqref{eq:ch4-restricted-zero}。

要使估计与所有此类函数正交，可以在 $t\leq1$ 上取常数。尝试
\[
 h(t)=\begin{cases}c,&t\leq1,\\bt,&t>1.\end{cases}
\]
积分得
\[
 \E_\theta h(T)=\frac{bn}{n+1}\theta
 +\left(c-\frac{bn}{n+1}\right)\theta^{-n}.
\]
取 $b=(n+1)/n$、$c=1$，得到
\begin{equation}
 \widehat\theta_*=\begin{cases}
 1,&T\leq1,\\
 (n+1)T/n,&T>1.
 \end{cases}
 \label{eq:ch4-restricted-umvue}
\end{equation}
它无偏且在每个参数下有界。对零均值 $v(T)$，只有 $T\leq1$ 部分有贡献，而该处 $h(T)=1$，故
\[
 \E_\theta[h(T)v(T)]
 =\frac n{\theta^n}\int_0^1v(t)t^{n-1}\,dt=0.
\]
充分统计量上的正交判别证明式~\eqref{eq:ch4-restricted-umvue} 为 UMVUE，无须完备性。

原估计 $U_0=(n+1)T/n$ 在限制模型中仍无偏。其与新估计的差为上述 $v(T)$，并与新估计正交。因此
\begin{align*}
 \Var_\theta(U_0)-\Var_\theta(\widehat\theta_*)
 &=\E_\theta v(T)^2\\
 &=\frac n{\theta^n}\int_0^1
   \left(\frac{n+1}{n}t-1\right)^2t^{n-1}\,dt\\
 &=\frac{\theta^{-n}}{n(n+2)}>0.
\end{align*}
于是
\[
 \Var_\theta(\widehat\theta_*)
 =\frac{\theta^2-\theta^{-n}}{n(n+2)}.
\]
在 $\theta=1$ 处，新估计恒为 $1$，方差为零。缩小参数空间会放宽“对所有参数均无偏”的要求，可能出现更多零均值方向与竞争者；UMVUE 因而依赖整个模型，不能只看密度的形式。

\begin{notebox}[练习与证明线索]
\begin{enumerate}
 \item 用正交判别证明两个 UMVUE 的固定线性组合仍为 UMVUE，再解释为什么 $U^2$ 未必无偏估计 $g(\theta)^2$。
 \item 将限制改成 $\theta\geq a>0$，求 $\theta$ 的 UMVUE。\textbf{答案：}$T\leq a$ 时取 $a$，$T>a$ 时取 $(n+1)T/n$；方差为 $[\theta^2-a^{n+2}\theta^{-n}]/[n(n+2)]$。
 \item 证明限制模型中的任一 UMVUE 若写成 $h(T)$，则它在 $(0,1)$ 上必须几乎处处为常数。\textbf{线索：}令 $c=\E_1h(T)$，用 $(h(T)-c)\ind_{\{T\leq1\}}$ 作零均值方向，在 $\theta=1$ 处应用正交性。
\end{enumerate}
\end{notebox}

\section{U 统计量：在所有子样本上平均}

\subsection{核与总体泛函}

设 $X_1,\ldots,X_n\overset{\mathrm{i.i.d.}}\sim P$，总体量具有表示
\[
 \vartheta(P)=\E_P h(X_1,\ldots,X_m),\qquad \E_P|h|<\infty,
\]
其中 $m\leq n$ 固定，$h$ 是不含未知分布的可测函数，称为核。“泛函”表示它以分布 $P$ 为输入，输出一个数；例如总体方差。这里的核既不是密度，也不要求为正定函数。

若原核不对称，可用
\[
 h_s(x_1,\ldots,x_m)=\frac1{m!}\sum_\pi
 h(x_{\pi(1)},\ldots,x_{\pi(m)})
\]
代替，和式遍历全部排列。独立同分布保证排列后的向量同分布，所以 $\E h_s=\E h$；若有二阶矩，Jensen 还给出 $\E h_s^2\leq\E h^2$。以下采用对称核。

\begin{definitionbox}{单样本 U 统计量}
由 $m$ 元对称核 $h$ 定义
\[
 U_n=\binom nm^{-1}
 \sum_{1\leq i_1<\cdots<i_m\leq n}
 h(X_{i_1},\ldots,X_{i_m}).
\]
每个核项期望均为 $\vartheta(P)$，故 $\E_PU_n=\vartheta(P)$。阶数 $m$ 是每项使用的观测数，并非多项式次数。
\end{definitionbox}

\subsection{独立随机子集的条件化解释}

在数据之外引入独立随机变量 $J$，它从全部 $m$ 元下标子集中均匀选取一个，令 $Z=h(X_j:j\in J)$。独立性是这里的条件：若根据观测大小偏向选择某些子集，以下等式一般不成立。

给定整份样本 $X$ 后，唯一剩下的随机性是选中哪个子集。因 $P(J=A\mid X)=\binom nm^{-1}$，
\[
 \E(Z\mid X)=\sum_{|A|=m}h(X_i:i\in A)P(J=A\mid X)=U_n.
\]
每个子集均有 $m$ 个独立同分布观测，故 $\E Z=\vartheta$，塔式法则给出 $\E U_n=\vartheta$。若 $\E h^2<\infty$，则
\[
 \Var(Z)=\Var(U_n)+\E\{\Var(Z\mid X)\}\geq\Var(U_n).
\]
这是消除额外抽选波动的方差比较，没有单凭此式断言 $U_n$ 在任意参数模型中都是 UMVUE。

\begin{examplebox}{一份数据固定以后还剩什么随机性}
取 $n=3,m=2$、$h(x,y)=(x+y)/2$。若观测为 $(1,3,8)$，三个可能子集给出 $2$、$4.5$、$5.5$，各以概率 $1/3$ 被选中。条件平均为
\[
 \E(Z\mid X=(1,3,8))=\frac{2+4.5+5.5}{3}=4.
\]
一般地，三个二元均值的平均正好为 $\bar X$。单次随机子集估计的方差为 $\sigma^2/2$，$U_3=\bar X$ 的方差为 $\sigma^2/3$；平均消除了子集抽选造成的那部分波动。
\end{examplebox}

\subsection{均值的幂与样本方差}

若 $\E|X|<\infty$、$\mu=\E X$，独立性给出 $\E(X_1\cdots X_m)=\mu^m$。因此乘积核产生 $\mu^m$ 的无偏估计。二阶时，
\[
 U_n=\frac{(\sum_iX_i)^2-\sum_iX_i^2}{n(n-1)}.
\]
若二阶矩有限，还可写为 $\bar X^2-S^2/n$。这个无偏性并不限于正态总体；它在何种参数模型下为 UMVUE，要另行证明。

若 $\E X^2<\infty$，取 $h(x,y)=(x-y)^2/2$，则 $\E h(X_1,X_2)=\sigma^2$。利用
\begin{align*}
 \sum_{i<j}(X_i-X_j)^2
 &=n\sum_iX_i^2-\left(\sum_iX_i\right)^2\\
 &=n\sum_i(X_i-\bar X)^2,
\end{align*}
对应的 U 统计量恰为
\[
 U_n=\frac1{n(n-1)}\sum_{i<j}(X_i-X_j)^2=S^2.
\]
二阶矩足以保证此无偏性；分析其方差和通常的方差 CLT 则需四阶矩。包含重复下标的平均并非同一个对象，例如 $n^{-2}\sum_{i,j}X_iX_j=\bar X^2$ 含有对角项，通常有偏估计 $\mu^2$，属于 V 统计量形式。

\section{条件核与精确方差}

本节及后续单样本理论均假定 $m$ 与 $P$ 固定，$\E h^2<\infty$，令 $n\to\infty$。若核阶数也随样本量变化，下面固定阶的余项估计不能直接沿用。

\subsection{重叠下标决定协方差}

定义条件核
\[
 h_k(x_1,\ldots,x_k)
 =\E h(x_1,\ldots,x_k,X_{k+1},\ldots,X_m),
 \quad 1\leq k\leq m,
\]
以及 $h_0=\vartheta$、$h_m=h$，并记
\[
 \zeta_k=\Var\{h_k(X_1,\ldots,X_k)\},\qquad\zeta_0=0.
\]
塔式法则和条件 Jensen 给出 $\E h_k=\vartheta$、$\E h_k^2\leq\E h^2$。

对两个 $m$ 元下标集 $A,B$，设 $C=A\cap B$、$|C|=k$。给定公共观测 $X_C$，两核其余输入来自互不相交的观测，因而条件独立。其中心化条件均值均为 $h_k(X_C)-\vartheta$，所以
\begin{align*}
 \Cov\{h(X_A),h(X_B)\}
 &=\E\bigl[\E\{(h(X_A)-\vartheta)
          (h(X_B)-\vartheta)\mid X_C\}\bigr]\\
 &=\E\{[h_k(X_C)-\vartheta]^2\}=\zeta_k.
\end{align*}
无公共下标时协方差为零；$A=B$ 时协方差为 $\zeta_m=\Var(h)$。

\begin{definitionbox}{Hoeffding 精确方差公式}
对 $n\geq m$，
\begin{equation}
 \Var(U_n)=\binom nm^{-1}
 \sum_{k=1}^m\binom mk\binom{n-m}{m-k}\zeta_k.
 \label{eq:ch4-exact-variance}
\end{equation}
选取数目不合法时，组合数按零处理。特别地，二阶核满足
\[
 \Var(U_n)=\frac{4(n-2)}{n(n-1)}\zeta_1
          +\frac2{n(n-1)}\zeta_2.
\]
\end{definitionbox}

\textbf{计数证明。} 令 $N=\binom nm$，$H_A=h(X_A)-\vartheta$。展开
\[
 \Var(U_n)=\frac1{N^2}\sum_{|A|=m}\sum_{|B|=m}\E H_AH_B.
\]
固定 $A$，要使 $|A\cap B|=k$，须从 $A$ 内选 $k$ 个下标，再从 $A$ 外选 $m-k$ 个下标，恰有 $\binom mk\binom{n-m}{m-k}$ 种。外层有 $N$ 个 $A$，故相应有序对数目为 $N\binom mk\binom{n-m}{m-k}$。双和原本计算有序对，不能再除以 $2$。代入协方差 $\zeta_k$，并利用 $\zeta_0=0$，即得公式。当 $n=m$ 时只剩一项，公式为 $\Var(U_n)=\zeta_m$，与定义一致。\hfill$\square$

\subsection{方差速度与退化阶数}

令 $H_k=h_k(X_1,\ldots,X_k)$，$\mathcal F_k=\sigma(X_1,\ldots,X_k)$。逐次条件化给出 $\E(H_{k+1}\mid\mathcal F_k)=H_k$，因此
\[
 \zeta_{k+1}=\zeta_k+
 \E\{\Var(H_{k+1}\mid\mathcal F_k)\}\geq\zeta_k.
\]
于是 $0\leq\zeta_1\leq\cdots\leq\zeta_m$。对固定 $m,k$，精确方差系数满足
\begin{align*}
 c_{n,k}
 &:=\frac{\binom mk\binom{n-m}{m-k}}{\binom nm}\\
 &=\frac{k!\binom mk^2}{n^k}
 \frac{\prod_{r=0}^{m-k-1}(1-(m+r)/n)}
      {\prod_{r=0}^{m-1}(1-r/n)}\\
 &=\frac{k!\binom mk^2}{n^k}+O(n^{-k-1}).
\end{align*}
空乘积取 $1$。由于每个乘积的因子数固定，最后的比值为 $1+O(n^{-1})$。若第一个非零条件核方差为 $\zeta_k$，则
\begin{equation}
 \Var(U_n)=\frac{k!\binom mk^2\zeta_k}{n^k}
            +O(n^{-k-1}).
 \label{eq:ch4-var-rate}
\end{equation}
特别地，$\zeta_1>0$ 时方差为 $m^2\zeta_1/n+O(n^{-2})$；$\zeta_1=0$、$\zeta_2>0$ 时为 $n^{-2}$ 量级。

各系数趋零，故 $\E(U_n-\vartheta)^2\to0$，即 $U_n\to\vartheta$ 于 $L^2$，从而也依概率收敛。如果全部 $\zeta_k=0$，则核本身几乎处处为常数，结论仍成立。方差的速度给出可能的放大尺度，却不单独保证极限分布存在。

\begin{notebox}[练习与答案线索]
\begin{enumerate}
 \item 对 $h(x,y)=xy$、$\E X=\mu$、$\Var(X)=\sigma^2<\infty$，求 $\vartheta,h_1,\zeta_1,\zeta_2$。\textbf{答案：}$\vartheta=\mu^2$，$h_1(x)=\mu x$，$\zeta_1=\mu^2\sigma^2$，$\zeta_2=\sigma^4+2\mu^2\sigma^2$。
 \item 将上题代入精确公式，证明 $\Var(U_n)=4\mu^2\sigma^2/n+2\sigma^4/[n(n-1)]$。
 \item 为什么不能把 $\binom nm$ 个核项当成独立项？\textbf{线索：}即使样本观测独立，两个核项仍可能共享下标。逐一写出 $h(X_1,X_2)$ 与 $h(X_1,X_3)$ 的协方差。
\end{enumerate}
\end{notebox}

\section{H\'ajek 投影与非退化极限}

\subsection{投影到单变量加性函数}

对独立观测 $X_1,\ldots,X_n$ 及任意二阶可积 $W$，定义
\[
 a_i(X_i)=\E(W\mid X_i)-\E W,\qquad
 \widehat W=\E W+\sum_{i=1}^na_i(X_i).
\]
独立性保证当 $j\neq i$ 时 $\E[a_j(X_j)\mid X_i]=0$，从而
\[
 \E(W-\widehat W\mid X_i)=0.
\]
所以残差与每个二阶可积单变量函数 $b_i(X_i)$ 正交，也与其有限和及常数正交。对任何加性函数 $Q=c+\sum_ib_i(X_i)$，
\[
 \E(W-Q)^2=\E(W-\widehat W)^2+\E(\widehat W-Q)^2.
\]
这称为 H\'ajek 投影：允许每个单变量函数非线性，但不保留不同观测间的交互项。它并不是对完整数据取条件期望；后者对统计量 $W$ 仍给出 $W$ 本身。例如零均值独立变量的乘积 $X_1X_2$，对两个单变量的条件期望均为零，所以其 H\'ajek 投影为零，而完整数据仍确定该乘积。

\subsection{U 统计量的一阶项与余项}

令
\[
 \varphi(x)=h_1(x)-\vartheta,
 \qquad \E\varphi(X)=0,\quad\Var\{\varphi(X)\}=\zeta_1.
\]
固定下标 $i$，包含它的 $m$ 元子集有 $\binom{n-1}{m-1}$ 个，不包含它的有 $\binom{n-1}{m}$ 个。因此
\[
 \E(U_n\mid X_i)-\vartheta
 =\frac{\binom{n-1}{m-1}}{\binom nm}\varphi(X_i)
 =\frac mn\varphi(X_i).
\]
故投影为
\begin{equation}
 \widehat U_n=\vartheta+L_n,
 \qquad L_n=\frac mn\sum_{i=1}^n\varphi(X_i),
 \qquad R_n=U_n-\widehat U_n.
 \label{eq:ch4-hajek}
\end{equation}
由独立性，$\Var(L_n)=m^2\zeta_1/n$。同时
\begin{align*}
 \E[(U_n-\vartheta)\varphi(X_i)]
 &=\E[\E(U_n-\vartheta\mid X_i)\varphi(X_i)]\\
 &=\frac mn\zeta_1,
\end{align*}
因此 $\Cov(U_n,L_n)=m^2\zeta_1/n=\Var(L_n)$。展开残差平方得到精确恒等式
\begin{equation}
 \E R_n^2=\Var(U_n)-\frac{m^2\zeta_1}{n}.
 \label{eq:ch4-residual}
\end{equation}
这也就是 $R_n$ 与 $L_n$ 正交后的勾股关系，不需要二者独立。将方差展开代入，得 $\E R_n^2=O(n^{-2})$，于是
\[
 P(|\sqrt nR_n|>\varepsilon)
 \leq\frac{n\E R_n^2}{\varepsilon^2}\longrightarrow0.
\]
所以 $R_n=O_P(n^{-1})$，且
\begin{equation}
 \sqrt n(U_n-\vartheta)
 =\frac m{\sqrt n}\sum_{i=1}^n\varphi(X_i)+o_P(1).
 \label{eq:ch4-asym-linear}
\end{equation}

\begin{definitionbox}{非退化 U 统计量的中心极限定理}
若核阶数 $m$ 固定、$\E h^2<\infty$ 且 $\zeta_1>0$，则
\[
 \sqrt n(U_n-\vartheta)\dto N(0,m^2\zeta_1).
\]
\end{definitionbox}

\textbf{证明。} $\varphi(X_i)$ 独立同分布，均值为零、有限方差为 $\zeta_1$。普通中心极限定理给出 $n^{-1/2}\sum_i\varphi(X_i)\dto N(0,\zeta_1)$；将式~\eqref{eq:ch4-asym-linear} 代入 Slutsky 定理即可。\hfill$\square$

这里的极限方差来自一阶条件核，不是原核方差 $\zeta_m$；核阶数还贡献系数 $m^2$。若 $\zeta_1=0$，同一推导只能得到 $\sqrt n(U_n-\vartheta)\pto0$，更细的波动须换用其他尺度。

\subsection{样本方差：与 Delta 方法的核对}

对 $h(x,y)=(x-y)^2/2$，$\vartheta=\sigma^2$。固定 $x$ 并平均另一个观测，得
\[
 h_1(x)=\frac12[(x-\mu)^2+\sigma^2],\qquad
 \varphi(x)=\frac12[(x-\mu)^2-\sigma^2].
\]
设 $\mu_4=\E(X-\mu)^4<\infty$。由于
\[
 (X_1-X_2)^4\leq8[(X_1-\mu)^4+(X_2-\mu)^4],
\]
核二阶矩有限，并且
\[
 4\zeta_1=\Var\{(X-\mu)^2\}=\mu_4-\sigma^4.
\]
若 $\mu_4>\sigma^4$，则
\begin{equation}
 \sqrt n(S^2-\sigma^2)\dto N(0,\mu_4-\sigma^4).
 \label{eq:ch4-variance-clt}
\end{equation}
它与第一章通过样本一、二阶矩及 Delta 方法所得结果一致：两种方法找到的主导项都是 $(X-\mu)^2-\sigma^2$。正态总体有 $\mu_4=3\sigma^4$，极限方差为 $2\sigma^4$。

还可计算精确方差。令 $Z_i=X_i-\mu$，独立性与 $\E Z_i=0$ 给出
\[
 \E(Z_1-Z_2)^4=2\mu_4+6\sigma^4,
 \qquad\zeta_2=\frac12(\mu_4+\sigma^4).
\]
代入二阶精确公式，得
\[
 \Var(S^2)=\frac1n
 \left[\mu_4-\frac{n-3}{n-1}\sigma^4\right].
\]
正态时为 $2\sigma^4/(n-1)$，与样本方差的精确卡方分布相符。

\begin{examplebox}{样本方差也可能一阶退化}
若 $P(X=\mu+a)=P(X=\mu-a)=1/2$，$a>0$，则 $(X-\mu)^2=a^2$ 恒定，故 $\mu_4-\sigma^4=0$。式~\eqref{eq:ch4-variance-clt} 的 $\sqrt n$ 极限退化为零，不能再除以该极限标准差做正态标准化。

写 $X_i=\mu+a\varepsilon_i$，其中 $P(\varepsilon_i=\pm1)=1/2$。则
\[
 S^2=\frac{na^2}{n-1}(1-\bar\varepsilon^2),
\]
所以
\[
 n(S^2-a^2)=\frac n{n-1}a^2
 \left[1-\left(\frac1{\sqrt n}\sum_i\varepsilon_i\right)^2\right]
 \dto a^2(1-Z^2),\quad Z\sim N(0,1).
\]
退化后出现的是中心化卡方型极限，而非非零方差的正态极限。
\end{examplebox}

\begin{notebox}[练习与答案线索]
\begin{enumerate}
 \item 对乘积核 $h(x,y)=xy$，在 $\mu\neq0$、$\sigma^2>0$ 时求渐近分布。\textbf{答案：}$\sqrt n(U_n-\mu^2)\dto N(0,4\mu^2\sigma^2)$。
 \item 对 $h(x,y)=(x+y)/2$ 求 H\'ajek 投影和余项。\textbf{答案：}$U_n=\bar X$，投影等于统计量本身，余项恒为零；虽是二阶核，却没有真正的二阶交互。
 \item 对 $X\sim U(0,1)$，求 $\sqrt n(S^2-1/12)$ 的极限方差。\textbf{答案：}$\mu_4=1/80$，故方差为 $1/80-1/144=1/180$。
\end{enumerate}
\end{notebox}

\section{两样本秩统计量}

\subsection{Mann--Whitney 与 Wilcoxon 秩和}

本节的 $m,n$ 均表示样本量，不再表示单样本核阶数。设 $X_1,\ldots,X_m\sim F$、$Y_1,\ldots,Y_n\sim G$，两组内部独立同分布，组间独立，$F,G$ 连续。约定
\[
 h(x,y)=\ind_{\{x>y\}},\qquad
 \vartheta=P(X>Y),\qquad
 U_{m,n}=\frac1{mn}\sum_{i=1}^m\sum_{j=1}^n\ind_{\{X_i>Y_j\}}.
\]
它是所有跨组比较中 $X$ 大于 $Y$ 的比例，无偏估计 $\vartheta$。部分教材或软件把未归一化的胜出次数 $mnU_{m,n}$ 称为 Mann--Whitney $U$，另一些使用相反的比较方向，核对约定后才能直接引用公式。

把 $m+n$ 个观测合并排序，无并列值时，记 $X_i$ 的合并秩为 $R_i$，Wilcoxon 秩和为 $R_X=\sum_iR_i$。把 $X$ 组内部排序为 $X_{(1)}<\cdots<X_{(m)}$，相应合并秩为 $r_a$。位于 $X_{(a)}$ 之前的 $r_a-1$ 个观测中，$a-1$ 个来自 $X$ 组，所以
\[
 r_a-a=\#\{j:Y_j<X_{(a)}\}.
\]
对 $a$ 求和即得
\begin{equation}
 R_X-\frac{m(m+1)}2=mnU_{m,n}.
 \label{eq:ch4-rank-equivalence}
\end{equation}
因此秩和与跨组胜出次数是已知的线性变换，携带相同信息；U 统计量的分析可以直接转给秩和。若合并次序为 $Y,X,Y,Y,X,Y,X$，则 $m=3,n=4$，$X$ 的秩为 $(2,5,7)$，胜出次数为 $1+3+4=8$，故 $R_X=14$、$U_{3,4}=2/3$。

\subsection{两个一阶投影与余项控制}

连续性给出 $\E[h(x,Y)]=G(x)$、$\E[h(X,y)]=1-F(y)$。定义
\[
 a(x)=G(x)-\vartheta,\qquad
 b(y)=1-F(y)-\vartheta,
\]
以及
\[
 r(x,y)=h(x,y)-\vartheta-a(x)-b(y).
\]
两个投影均值为零，且 $\E[r(X,Y)\mid X]=\E[r(X,Y)\mid Y]=0$。于是有精确分解
\begin{equation}
 U_{m,n}-\vartheta
 =\frac1m\sum_i a(X_i)+\frac1n\sum_j b(Y_j)+R_{m,n},
 \label{eq:ch4-twosample-projection}
\end{equation}
其中 $R_{m,n}=(mn)^{-1}\sum_{i,j}r(X_i,Y_j)$。

残差核与两个单变量项均正交。两个不同配对若没有共享下标，则独立；若只共享 $X_i$，给定 $X_i$ 后两残差核条件独立且条件均值零，协方差也为零；只共享 $Y_j$ 同理。故仅完全相同的下标对贡献方差，
\[
 \E R_{m,n}^2=\frac{\E r(X,Y)^2}{mn}.
\]
若 $N=m+n$、$m/N\to\rho\in(0,1)$，则 $N/(mn)\to0$，因而 $\sqrt N R_{m,n}\pto0$。两组主项相互独立，分别使用 CLT 得到
\[
 \sqrt N(U_{m,n}-\vartheta)\dto
 N\left(0,\frac{\Var(G(X))}{\rho}
          +\frac{\Var(F(Y))}{1-\rho}\right).
\]
若显示的极限方差为零，应另作退化分析，而不能使用非零方差的标准化。

\subsection{相同分布原假设下的检验}

在 $H_0:F=G$ 下，连续性与交换性给出 $\vartheta=1/2$，概率积分变换给出 $F(X),F(Y)\sim U(0,1)$。因此两个投影方差都是 $1/12$，有
\[
 \sqrt{m+n}\left(U_{m,n}-\frac12\right)
 \dto N\left(0,\frac1{12\rho(1-\rho)}\right),
\]
等价地，
\[
 Z_{m,n}=\sqrt{\frac{12mn}{m+n}}
 \left(U_{m,n}-\frac12\right)\dto N(0,1).
\]
于是 $|Z_{m,n}|>z_{1-\alpha/2}$ 为渐近双侧水平 $\alpha$ 的拒绝规则。渐近水平表示拒绝概率趋于 $\alpha$，不表示每个有限样本都精确等于 $\alpha$。

在原假设下，$\Var(h)=1/4$。由核的正交分解，$\E r^2=1/4-1/12-1/12=1/12$，所以精确方差为
\[
 \Var_0(U_{m,n})=\frac1{12m}+\frac1{12n}+\frac1{12mn}
 =\frac{m+n+1}{12mn}.
\]
相应地，
\[
 \E_0R_X=\frac{m(m+n+1)}2,\qquad
 \Var_0(R_X)=\frac{mn(m+n+1)}{12}.
\]
用精确方差作正态近似时，标准化可采用 $\sqrt{12mn/(m+n+1)}$，与上述渐近形式等价。有并列值时，须明确比较核、平均秩或其他并列处理，并重新修正方差；这里的无并列值公式不能直接照搬。

\begin{notebox}[检验目标与练习]
$F=G$ 蕴含 $P(X>Y)=1/2$，反之不成立。例如 $X\sim N(0,1)$、$Y\sim N(0,4)$ 独立时，$X-Y$ 关于零对称，仍有 $P(X>Y)=1/2$。因此该统计量不能保证对所有 $F\neq G$ 的备择都具有趋于 $1$ 的功效。
\begin{enumerate}
 \item 若改用 $\ind_{\{X<Y\}}$，秩和等式如何变化？\textbf{答案：}新统计量为 $1-U_{m,n}$，故 $R_X=m(m+1)/2+mn(1-U_{m,n}^{\mathrm{new}})$。
 \item 当 $m=n$ 时写出 $U_{m,n}$ 与秩和的原假设精确方差，再比较两种标准化。
 \item 用条件均值零逐一证明残差核共享一个下标时协方差仍为零；这一步解释了为何两样本余项只保留 $1/(mn)$ 量级。
\end{enumerate}
\end{notebox}

\section{一阶退化与中心化卡方极限}

本节保留一般证明，也先列出可用普通 CLT 与 LLN 独立核验的例子。需要额外借用的泛函分析结论只有平方可积对称核的谱展开；其余概率步骤均在此前工具范围内。

\subsection{乘积核的直接计算}

设 $\E X=0$、$\E X^2=\sigma^2>0$，取 $h(x,y)=xy$。此时 $\vartheta=0$、$h_1(x)=0$、$\zeta_2=\sigma^4$，精确方差为 $2\sigma^4/[n(n-1)]$。由乘积恒等式，
\[
 nU_n=\frac n{n-1}
 \left[\left(\frac1{\sqrt n}\sum_iX_i\right)^2
       -\frac1n\sum_iX_i^2\right].
\]
CLT 给出括号第一项趋于 $\sigma^2Z^2$，LLN 给出第二项依概率趋于 $\sigma^2$。因此
\[
 nU_n\dto\sigma^2(Z^2-1),\qquad Z\sim N(0,1).
\]
这里只需二阶矩，因为 $\E h(X_1,X_2)^2=(\E X^2)^2$。中心化卡方来自“和的平方减去平均平方”，而不是把退化情形继续当作非退化 CLT 使用。

\subsection{积分算子是矩阵作用的推广}

回到固定 $m\geq2$ 的一般核，假设 $\zeta_1=0$、$\zeta_2>0$。定义
\[
 K(x,y)=h_2(x,y)-\vartheta.
\]
因为 $h_1(X)=\vartheta$ 几乎处处，
\[
 \E[K(X_1,X_2)\mid X_1]=0,
 \qquad\E K(X_1,X_2)^2=\zeta_2.
\]
一般非退化情形的二阶正交核应为 $h_2(x,y)-h_1(x)-h_1(y)+\vartheta$；在本节条件下，它才等于 $K$。

考虑 $L^2(P)$ 上的积分算子
\[
 (\mathcal Kf)(x)=\int K(x,y)f(y)\,dP(y).
\]
若 $P$ 在 $r$ 个点 $x_1,\ldots,x_r$ 等概率取值，积分变成
\[
 (\mathcal Kf)(x_i)=\frac1r\sum_{j=1}^rK(x_i,x_j)f(x_j),
\]
即矩阵 $(K(x_i,x_j)/r)_{i,j}$ 乘上函数值向量。一般分布只是把离散求和换成积分，把向量换成函数。特征方程为
\[
 \int K(x,y)e_j(y)\,dP(y)=\lambda_je_j(x).
\]
对称核给出自伴算子，核平方可积给出 Hilbert--Schmidt 性质。在通常的实值观测（标准 Borel）设置中，谱定理允许在 $L^2(P\times P)$ 意义下写为
\begin{equation}
 K(x,y)=\sum_{j\geq1}\lambda_je_j(x)e_j(y),
 \qquad\sum_j\lambda_j^2=\zeta_2,
 \label{eq:ch4-spectrum}
\end{equation}
其非零特征值均实，特征函数正交归一，即 $\E[e_j(X)e_\ell(X)]=\ind_{\{j=\ell\}}$。等式是均方收敛，不要求每个 $(x,y)$ 都逐点成立。

退化性和对称性给出 $\int K(x,y)\,dP(x)=0$ 对几乎所有 $y$ 成立。对特征方程积分，利用 Cauchy--Schwarz 保证绝对可积，可得
\[
 \lambda_j\E e_j(X)
 =\iint K(x,y)e_j(y)\,dP(y)dP(x)=0.
\]
所以非零特征值对应的 $e_j$ 均值为零。乘积例中只有一个方向：$e(x)=x/\sigma$、$\lambda=\sigma^2$，于是 $K(x,y)=\sigma^2e(x)e(y)$，正好恢复上一小节的极限。

\begin{definitionbox}{一阶退化 U 统计量的极限}
若 $m\geq2$ 固定、$\E h^2<\infty$，且 $\zeta_1=0$、$\zeta_2>0$，则
\begin{equation}
 n(U_n-\vartheta)\dto
 \binom m2\sum_{j\geq1}\lambda_j(Z_j^2-1),
 \label{eq:ch4-degenerate-limit}
\end{equation}
其中 $Z_j$ 相互独立且服从 $N(0,1)$，$\lambda_j$ 为式~\eqref{eq:ch4-spectrum} 的非零特征值。无限和按 $L^2$ 收敛理解，权重可以为正也可以为负。
\end{definitionbox}

\subsection{第一步：严格降到二阶核}

令
\[
 V_n=\binom n2^{-1}\sum_{i<j}K(X_i,X_j),
\]
并定义扣除所有成对贡献后的核
\[
 q(x_1,\ldots,x_m)=h(x_1,\ldots,x_m)-\vartheta
                  -\sum_{1\leq a<b\leq m}K(x_a,x_b).
\]
它对称、均值为零且二阶可积。给定 $X_1,X_2$，原核的条件均值为 $h_2(X_1,X_2)$；成对和中只有 $K(X_1,X_2)$ 保留。只含一个已知下标的项因退化性条件均值零，两个下标都未知的项也均值零。因此
\[
 \E[q(X_1,\ldots,X_m)\mid X_1,X_2]
 =h_2(X_1,X_2)-\vartheta-K(X_1,X_2)=0.
\]
再取条件期望得一阶投影也为零。对称性保证任意一个或一对下标均有该性质。

固定样本中的一对下标，它在 $\binom{n-2}{m-2}$ 个 $m$ 元子集中出现。于是
\begin{align*}
 U_n-\vartheta
 &=\frac{\binom{n-2}{m-2}}{\binom nm}
       \sum_{i<j}K(X_i,X_j)+U_n(q)\\
 &=\binom m2V_n+U_n(q).
\end{align*}
这给出了式~\eqref{eq:ch4-degenerate-limit} 中不能遗漏的 $\binom m2$。若 $m\geq3$，核 $q$ 的前两阶条件方差均为零，精确方差公式只剩 $k\geq3$ 项，故 $\Var\{U_n(q)\}=O(n^{-3})$，
\[
 \E\{[nU_n(q)]^2\}=O(n^{-1})\to0.
\]
若 $m=2$，$q$ 恒为零。因此一般问题归结为证明 $nV_n$ 的极限。

\subsection{第二步：固定有限谱应用 CLT}

令 $K_J(x,y)=\sum_{j=1}^J\lambda_je_j(x)e_j(y)$，并记对应二阶 U 统计量为 $V_{n,J}$。利用成对乘积恒等式，
\begin{equation}
 \begin{split}
 nV_{n,J}=\frac n{n-1}\sum_{j=1}^J\lambda_j
 \bigg[&\left(\frac1{\sqrt n}\sum_{i=1}^ne_j(X_i)\right)^2\\
       &-\frac1n\sum_{i=1}^ne_j(X_i)^2\bigg].
 \end{split}
 \label{eq:ch4-finite-spectrum}
\end{equation}
对固定 $J$，向量 $(e_1(X_i),\ldots,e_J(X_i))$ 独立同分布，均值零、协方差为单位阵。多元 CLT 给出
\[
 \frac1{\sqrt n}\sum_{i=1}^n
 (e_1(X_i),\ldots,e_J(X_i))^\mathsf T
 \dto (Z_1,\ldots,Z_J)^\mathsf T\sim N_J(0,I_J).
\]
有限样本下，同一个 $X_i$ 产生的各个 $e_j(X_i)$ 不必独立；正交性给出协方差阵，联合正态极限才使 $Z_j$ 独立。又因 $\E e_j(X)^2=1$，LLN 给出 $n^{-1}\sum_ie_j(X_i)^2\pto1$。这里只需 $e_j^2$ 可积，不要求特征函数有四阶矩。连续映射与 Slutsky 因而给出
\[
 nV_{n,J}\dto L_J:=\sum_{j=1}^J\lambda_j(Z_j^2-1).
\]

\subsection{第三步：以统一均方界处理无限尾项}

由于 $\Var(Z_j^2-1)=2$，$\sum_j\lambda_j^2<\infty$ 保证 $L_J$ 在 $L^2$ 中为 Cauchy 列，故存在
\[
 L=\sum_{j\geq1}\lambda_j(Z_j^2-1),\qquad
 \E(L-L_J)^2=2\sum_{j>J}\lambda_j^2\to0.
\]
平方可和不保证绝对可和。因此上式须按中心化级数的 $L^2$ 极限理解，不能拆成两个未必各自收敛的级数之差：
\[
 \sum_j\lambda_j Z_j^2-\sum_j\lambda_j.
\]

样本端同样需要尾项控制。核 $K-K_J$ 仍退化，其二阶矩为 $\sum_{j>J}\lambda_j^2$。由二阶退化精确方差公式，
\begin{equation}
 \E\{[n(V_n-V_{n,J})]^2\}
 =\frac{2n}{n-1}\sum_{j>J}\lambda_j^2
 \leq4\sum_{j>J}\lambda_j^2,\qquad n\geq2.
 \label{eq:ch4-spectral-tail}
\end{equation}
界不依赖 $n$，因此可以先固定 $J$ 令 $n$ 趋于无穷，再令 $J$ 趋于无穷。严格地，对任一有界且 Lipschitz 常数不超过 $1$ 的函数 $b$，
\begin{align*}
 |\E b(nV_n)-\E b(L)|
 &\leq\E|n(V_n-V_{n,J})|\\
 &\quad+|\E b(nV_{n,J})-\E b(L_J)|
       +\E|L_J-L|.
\end{align*}
固定 $J$ 后，中间项随 $n\to\infty$ 趋零。两边项由 Cauchy--Schwarz 及式~\eqref{eq:ch4-spectral-tail}，分别不超过
\[
 2\left(\sum_{j>J}\lambda_j^2\right)^{1/2},\qquad
 \left(2\sum_{j>J}\lambda_j^2\right)^{1/2},
\]
再令 $J\to\infty$，二者均趋零。有界 Lipschitz 函数刻画依分布收敛，故 $nV_n\dto L$。结合第一步的 $nU_n(q)\pto0$，使用 Slutsky 即得式~\eqref{eq:ch4-degenerate-limit}。\hfill$\square$

\begin{notebox}[退化情形的核验与练习]
\begin{enumerate}
 \item 对一般 $m$，用式~\eqref{eq:ch4-var-rate} 核对极限的方差系数。\textbf{答案：}$n^2\Var(U_n)\to2\binom m2^2\zeta_2$，而极限和的方差正是 $2\binom m2^2\sum_j\lambda_j^2$。
 \item 取三阶核 $h(x,y,z)=xy+xz+yz$，总体均值零、方差 $\sigma^2$。求 $h_2$ 与极限。\textbf{答案：}$h_2(x,y)=xy$，故 $nU_n\dto3\sigma^2(Z^2-1)$；直接计数也有 $U_n=3\binom n2^{-1}\sum_{i<j}X_iX_j$。
 \item 设 $a,b$ 为均值零、正交归一的平方可积函数，核为 $K(x,y)=a(x)a(y)-2b(x)b(y)$。\textbf{答案：}二阶 U 统计量满足 $nU_n\dto(Z_1^2-1)-2(Z_2^2-1)$。负权重并不妨碍核对称或平方可积，极限也不必非负。
 \item 指出无限谱证明中若只写“固定 $J$ 的 CLT，再令 $J=\infty$”，缺了哪一步。\textbf{线索：}必须同时控制样本端尾项与极限端尾项，且样本端的界须对 $n$ 一致。
\end{enumerate}
\end{notebox}

\section{综合练习}

\begin{enumerate}
 \item \textbf{条件化与完备性。} 设 $X_i\sim\operatorname{Bernoulli}(p)$ 独立同分布，$0<p<1$，$n\geq2$。从 $\ind_{\{X_1=1,X_2=0\}}$ 出发，求 $p(1-p)$ 的 UMVUE，要求写出条件分布与有限方差核验。\par
 \textbf{解答线索：}给定 $T=\sum_iX_i=t$，全部含 $t$ 个 $1$ 的序列等概率，所求条件概率为 $t(n-t)/[n(n-1)]$。$T$ 的自然参数遍历 $\R$，完备充分；估计有界，故适用 Lehmann--Scheff\'e。

 \item \textbf{同一核在不同参数点的变化。} 设 $\E X^2<\infty$、$\Var(X)=\sigma^2>0$。对乘积核 $h(x,y)=xy$，比较 $\mu\neq0$ 与 $\mu=0$ 时的方差速度、放大尺度和极限。\par
 \textbf{解答线索：}$\mu\neq0$ 时为 $n^{-1}$ 方差及 $\sqrt n$ 正态极限；$\mu=0$ 时为 $n^{-2}$ 方差及 $nU_n\dto\sigma^2(Z^2-1)$。退化由分布和核共同决定，不是核公式本身的固定标签。

 \item \textbf{构造但不误称最优。} 独立同分布总体具有二阶矩，证明 $\bar X^2-S^2/n$ 无偏估计 $\mu^2$；再说明在哪个已讨论的模型中可以进一步称它为 UMVUE。\par
 \textbf{解答线索：}一般总体由乘积核即可得到无偏性；正态双参数模型中它是完备充分统计量的函数，故为 UMVUE。无偏或 U 统计量形式本身不足以替代模型内的最优性论证。

 \item \textbf{由秩恢复胜出次数。} 两组各有三个观测，$X$ 的合并秩为 $(2,4,6)$。求 $U_{3,3}$、原假设下 $R_X$ 的均值与方差，并说明为什么仅此观测结果不能证明两总体相同。\par
 \textbf{答案：}$R_X=12$，胜出次数 $12-6=6$，$U_{3,3}=2/3$；原假设均值为 $21/2$、方差为 $21/4$。检验需参照零分布或近似分布，且不拒绝不等于证实原假设。

 \item \textbf{证明题。} 不直接引用全方差公式，从 $R=U-\E(U\mid T)$ 的条件均值为零出发，证明 Rao--Blackwell 方差比较及其等号条件；说明充分性用在何处。\par
 \textbf{解答线索：}先证残差与所有二阶可积 $T$ 的函数正交，再展开平方。方差恒等式对任意 $T$ 成立，充分性用来保证条件期望可取同一个参数无关的版本。

 \item \textbf{退化证明题。} 在 $\zeta_1=0$、$\zeta_2>0$ 时，不依赖完整的高阶 Hoeffding 分解，证明扣除全部成对核后的余项在 $n$ 尺度可忽略。\par
 \textbf{解答线索：}构造本章的 $q$，证明它的一、二阶条件投影均为零，再用重叠计数方差公式得到 $O(n^{-3})$；$m=2$ 时单独观察 $q=0$。最后仍须用有限谱 CLT 和统一尾界，才能确定极限分布。
\end{enumerate}

\chapter{假设检验与置信集合}

估计问题要求给出参数的数值，检验问题则要求根据样本在两个参数集合之间作出判断。本章先固定第一类错误的上界，再比较备择下的拒绝概率。似然比决定证据的排序；充分性与完备性负责处理干扰参数；置信集合由一族检验反演得到。

\begin{notebox}[阅读范围]
本章按完整课件整理。检验、$p$ 值、Neyman--Pearson 引理、单侧 UMP 及双侧 UMP 不存在的内容，对应原课件第 1--27 页；无偏检验与 UMPU 从第 28 页开始。这里的收录范围不表示已独立完成全部证明和习题。
\end{notebox}

\section{检验、功效与比较的范围}

设样本 $X$ 的分布属于 $\{P_\theta:\theta\in\Theta\}$，检验
\[
 H_0:\theta\in\Theta_0,\qquad H_1:\theta\in\Theta_1,
 \qquad \Theta_0\cap\Theta_1=\varnothing.
\]
简单假设只指定一个分布；复合假设包含多个分布。例如正态模型中，若方差未知，即使把均值固定为 $\mu_0$，原假设仍是复合假设。

\begin{definitionbox}{检验、功效、大小与水平}
检验是可测函数 $\varphi:\mathcal X\to[0,1]$，其值表示观察到 $x$ 后拒绝原假设的概率。功效函数和大小分别为
\[
 \beta_\varphi(\theta)=\E_\theta\varphi(X),\qquad
 \operatorname{size}(\varphi)=\sup_{\theta\in\Theta_0}\beta_\varphi(\theta).
\]
若大小不超过 $\alpha$，称为水平 $\alpha$ 检验；若大小恰为 $\alpha$，称为精确大小 $\alpha$ 的检验。在 $\Theta_0$ 上，$\beta_\varphi$ 是第一类错误概率；在 $\Theta_1$ 上，$1-\beta_\varphi$ 是第二类错误概率。
\end{definitionbox}

$\varphi(x)\in\{0,1\}$ 时规则不随机化。一般情形取与样本独立的 $V\sim U(0,1)$，在 $V\le\varphi(x)$ 时拒绝。对固定 $x$，有
\[
 P_V\{V\le\varphi(x)\}=\varphi(x),
\]
故这个实施步骤确实产生规定的条件拒绝概率。$\varphi$ 是规则，$V$ 用于实施规则，两者不能与功效函数混淆。

总拒绝的功效恒为 1，但大小也为 1；总不拒绝没有第一类错误，却不能检出任何备择。因此最优性只能在规定的比较类中讨论。在水平约束下增大备择功效，等价于减小第二类错误，并不表示两个错误概率可以同时无约束地最小化。

\begin{definitionbox}{MP 与 UMP}
对固定简单备择，功效最大的水平 $\alpha$ 检验称为最有力检验（MP）。若同一个水平 $\alpha$ 检验 $\varphi^*$ 满足
\[
 \E_\theta\varphi^*\ge\E_\theta\psi,
 \qquad \theta\in\Theta_1,
\]
对每个水平 $\alpha$ 检验 $\psi$ 都成立，则称为一致最有力检验（UMP）。这里的“一致”指同一个规则对所有备择都最优，与一致收敛无关。
\end{definitionbox}

精确大小是原假设上的上确界等于 $\alpha$，不要求每个原假设参数处的拒绝率都是 $\alpha$。这种逐点相等的性质称为相似性，后文另行讨论。

\begin{examplebox}{练习 5.1：水平与功效}
设 $Z\sim N(\delta,1)$，按 $Z>z_{1-\alpha}$ 拒绝，其中 $z_q=\Phi^{-1}(q)$。写出功效，并判断它是 $H_0:\delta\le0$ 的什么水平。

\textbf{答案线索。} $\beta(\delta)=1-\Phi(z_{1-\alpha}-\delta)$ 随 $\delta$ 递增，故 $\sup_{\delta\le0}\beta(\delta)=\beta(0)=\alpha$。参数越负，拒绝率越低，说明精确大小不等于各点拒绝率相同。
\end{examplebox}

\section{p 值与离散边界随机化}

\subsection{观察值与重复抽样}

先考虑简单原假设 $P_0$，且统计量 $T$ 越大越支持备择。对已经观察到的数据 $x$，常规右尾 $p$ 值为
\[
 p(x)=P_0\{T(X')\ge T(x)\},\qquad X'\sim P_0.
\]
此时 $x$ 和 $T(x)$ 已固定，概率针对假想的重新抽样 $X'$。这个数衡量原假设下出现至少同样偏向备择的结果有多罕见，不是原假设为真的概率。实验前的 $p(X)$ 则是一个随机变量。

\begin{examplebox}{正态右侧检验}
设 $X_1,\ldots,X_{25}$ 独立服从 $N(\mu,10^2)$，检验 $\mu=100$ 对 $\mu>100$。若观察到 $\bar x=104$，则
\[
 z=\frac{\bar x-100}{10/\sqrt{25}}=2,\qquad
 p(x)=P(Z'\ge2)=1-\Phi(2)\approx0.0228.
\]
在水平 $0.05$ 下拒绝，在水平 $0.01$ 下不拒绝。连续型统计量恰好等于观察值的概率为零，故用整块尾部概率衡量极端程度。
\end{examplebox}

\subsection{与嵌套拒绝域的联系}

若拒绝域 $R_\alpha$ 随 $\alpha$ 单调扩张，定义
\[
 p_R(x)=\inf\{\alpha\in(0,1]:x\in R_\alpha\},\qquad R_1=\mathcal X.
\]
假设对 $0<u<1$ 有集合意义的右连续性
\[
 R_u=\bigcap_{v>u}R_v.
\]
则
\[
 \{x:p_R(x)\le u\}=R_u.
\]
证明如下：若 $x\in R_u$，下确界自然不超过 $u$。反之，若下确界不超过 $u$，则对每个 $v>u$ 都能找到某个 $\alpha<v$ 使 $x\in R_\alpha$；由嵌套性 $x\in R_v$，再用右连续性得到 $x\in R_u$。因此若 $P_0(R_u)\le u$，便有 $P_0\{p_R(X)\le u\}\le u$。

正态右尾取 $R_\alpha=\{z:z\ge\Phi^{-1}(1-\alpha)\}$，直接计算
\[
 z\in R_\alpha
 \iff \alpha\ge1-\Phi(z),
\]
所以 $p_R(x)=1-\Phi(z)=p(x)$。任意嵌套族的下确界不一定等于上述右尾概率；两者相等还依赖拒绝域按同一统计量的尾部排序构造。离散模型尤其要说明采用的边界与随机化约定。

\subsection{均匀性、超均匀性与复合原假设}

若 $T$ 在简单原假设下的分布函数 $F_0$ 连续，则概率积分变换给出
\[
 F_0(T)\sim U(0,1),\qquad p(X)=1-F_0(T)\sim U(0,1).
\]
故 $P_0\{p(X)\le u\}=u$。连续性已经足够，不必额外要求 $F_0$ 严格递增。

一般称 $p$ 值有效，若
\[
 P_\theta\{p(X)\le u\}\le u,
 \qquad 0\le u\le1,\quad\theta\in\Theta_0.
\]
这种性质称为超均匀性，使规则 $p\le\alpha$ 控制水平。复合原假设必须对其中每个参数都成立，不能任取一个原假设参数计算尾概率。

一种通用构造是：若每个简单原假设下都有有效的 $p_\theta$，且上确界可测，则
\[
 p_{\rm sup}(x)=\sup_{\theta\in\Theta_0}p_\theta(x)
\]
在复合原假设下有效。因为真值为 $\theta$ 时，$p_{\rm sup}\ge p_\theta$，故
\[
 P_\theta(p_{\rm sup}\le u)\le P_\theta(p_\theta\le u)\le u.
\]
存在最不利边界分布时，这个上确界常可简化为边界处的尾概率。

\subsection{离散概率质量怎样补足}

对离散 $T$，记
\[
 a_t=P_0(T>t),\qquad m_t=P_0(T=t).
\]
若选 $c$ 使 $m_c>0$ 且 $a_c\le\alpha\le a_c+m_c$，则取
\[
 \gamma=\frac{\alpha-a_c}{m_c},\qquad
 \varphi(x)=\ind_{\{T(x)>c\}}+\gamma\ind_{\{T(x)=c\}}.
\]
于是 $\E_0\varphi=a_c+\gamma m_c=\alpha$。分子是尚缺的拒绝概率，分母是边界整块概率；二者的比值正是边界上应拒绝的比例。

\begin{examplebox}{二项分布：$15\%$ 的边界拒绝与 $10\%$ 的总拒绝}
设 $T\sim\operatorname{Bin}(4,1/2)$，要求大小 $\alpha=0.10$。只拒绝 $T=4$ 的大小为 $1/16$；再把 $T=3$ 全部加入便成为 $5/16$，超过预算。故
\[
 c=3,\qquad
 \gamma=\frac{1/10-1/16}{4/16}=\frac3{20}=0.15.
\]
$T=4$ 时全拒绝，$T=3$ 时取独立 $V\sim U(0,1)$，按 $V\le0.15$ 拒绝，其他值不拒绝。总拒绝概率为
\[
 \frac1{16}+\frac4{16}\times0.15=0.10.
\]
$0.15$ 是给定边界事件后的拒绝概率；$0.10$ 是对全部可能数据平均后的拒绝概率。
\end{examplebox}

定义随机化 $p$ 值
\[
 p_{\rm rand}(x,v)=a_{T(x)}+v m_{T(x)},\qquad
 V\sim U(0,1),\quad V\perp X.
\]
观察到 $T=t$ 后，它在区间 $[a_t,a_t+m_t]$ 内均匀分布。这个区间的长度恰为该取值的概率 $m_t$。对任意
$u\in[a_c,a_c+m_c]$，由独立性，
\begin{align*}
 P_0\{p_{\rm rand}(X,V)\le u\}
 &=P_0(T>c)+P_0(T=c)
       P\left(V\le\frac{u-a_c}{m_c}\right)\\
 &=a_c+m_c\frac{u-a_c}{m_c}=u.
\end{align*}
各原子区间的内部互不相交，长度总和为 $\sum_t m_t=1$。给定 $T=t$ 时的条件密度为 $1/m_t$，再乘上选中该区间的概率 $m_t$，得到区间内无条件密度为 1；其余部分的 Lebesgue 测度为零。因此在简单离散原假设下
\[
 p_{\rm rand}(X,V)\sim U(0,1).
\]
按 $p_{\rm rand}\le\alpha$ 拒绝，恰好等价于前述边界随机化规则。常规尾概率 $p(X)=a_T+m_T$ 满足 $p(X)\ge p_{\rm rand}(X,V)$，从而
\[
 P_0\{p(X)\le u\}\le u.
\]
这同时说明离散常规 $p$ 值为何通常保守。

\begin{examplebox}{练习 5.2：区分三种概率}
在上述二项例子中，观察到 $T=3$。求常规 $p$ 值、随机化 $p$ 值，并写出其在水平 $0.10$ 下拒绝的条件。再证明常规 $p$ 值对应的非随机检验大小为 $1/16$。

\textbf{答案线索。} $p=5/16$，$p_{\rm rand}=1/16+V/4$；后者不超过 $0.10$ 当且仅当 $V\le0.15$。常规规则只在 $T=4$ 时拒绝，所以大小为 $1/16<0.10$。
\end{examplebox}

\section{Neyman--Pearson 引理}

\subsection{定理及临界值的存在}

设 $P_0,P_1$ 对同一 $\sigma$-有限测度 $\nu$ 有密度 $f_0,f_1$。目标是
\[
 \max_{0\le\varphi\le1}\E_1\varphi,
 \qquad \E_0\varphi\le\alpha.
\]
用密度差比直接相除更方便：当 $f_0=0<f_1$ 时，该点只支持备择，应拒绝。

\begin{definitionbox}{Neyman--Pearson 引理}
对 $0<\alpha<1$，存在有限 $c\ge0$ 及 $\gamma\in[0,1]$，使
\[
 \varphi^*=\ind_{\{f_1>cf_0\}}
              +\gamma\ind_{\{f_1=cf_0\}},
 \qquad \E_0\varphi^*=\alpha.
\]
它是检验 $P_0$ 对 $P_1$ 的 MP 水平 $\alpha$ 检验。任何同样最优的水平检验，在 $\{f_1\ne cf_0\}$ 上都与它相等，$\nu$-几乎处处。若 $c>0$，所有最优检验都必须具有精确大小 $\alpha$；$c=0$ 时最优功效为 1，不一定需要用满大小。
\end{definitionbox}

\textbf{临界值存在的证明。}
在 $P_0$ 下，$f_0(X)>0$ 且两个密度可取有限版本，故似然比 $L=f_1(X)/f_0(X)$ 几乎处处有限且非负。令
\[
 q(t)=P_0(L>t),\qquad
 c=\inf\{t:q(t)\le\alpha\}.
\]
$q$ 单调不增，$t<0$ 时为 1，$t\to\infty$ 时趋于 0；它右连续，且左极限为 $q(t-)=P_0(L\ge t)$。所以 $c\in[0,\infty)$，并有
\[
 q(c)\le\alpha\le q(c-).
\]
若差为正，令 $\gamma=[\alpha-q(c)]/[q(c-)-q(c)]$；若差为零，则两端均为 $\alpha$，取 $\gamma=0$。这使
\[
 \E_0\varphi^*=q(c)+\gamma P_0(L=c)=\alpha.
\]
在 $P_0$ 零概率但 $f_1>0$ 的区域按密度差规则全拒绝，不影响大小。

\subsection{最优性与等号条件的证明}

任取 $0\le\psi\le1$ 且 $\E_0\psi\le\alpha$。逐点有
\[
 D=(\varphi^*-\psi)(f_1-cf_0)\ge0.
\]
因为密度差为正处 $\varphi^*=1$，为负处 $\varphi^*=0$，等于零处乘积为零。因此
\begin{align*}
 \E_1\varphi^*-\E_1\psi
 &=\int(\varphi^*-\psi)(f_1-cf_0)\,d\nu
        +c\int(\varphi^*-\psi)f_0\,d\nu\\
 &=\int D\,d\nu+c(\alpha-\E_0\psi)\ge0.
\end{align*}
若 $\psi$ 也最优，两个非负项都必须为零。$\int D=0$ 给出 $D=0$ 几乎处处，于是密度差严格为正处必须 $\psi=1$，严格为负处必须 $\psi=0$。等号边界的安排不受这一乘积限制，所以最优规则不一定唯一。若 $c>0$，第二项为零还强迫 $\E_0\psi=\alpha$。

若 $c=0$，$\varphi^*=1$ 于 $\{f_1>0\}$，故其功效已为 1。更一般地，若规则 $\psi$ 的大小 $a<\alpha$ 且功效 $b<1$，令
\[
 \varepsilon=\frac{\alpha-a}{1-a},\qquad
 \widetilde\psi=\psi+\varepsilon(1-\psi).
\]
则新规则大小恰为 $\alpha$，功效为 $b+\varepsilon(1-b)>b$。所以功效尚未到 1 的最优规则必须用满大小；这个结论不要求唯一性。

\subsection{似然比平台与不同的最优规则}

\begin{examplebox}{均匀分布端点：两个不同的 MP 检验}
设 $X_i$ 独立服从 $U(0,\theta)$，$M=X_{(n)}$，检验 $\theta_0$ 对 $\theta_1>\theta_0$。在原假设支持上，似然比恒为
$k=(\theta_0/\theta_1)^n$；在 $\theta_0\le M<\theta_1$ 的新增支持上为 $+\infty$。任意水平检验都满足
\begin{align*}
 \E_{\theta_1}\varphi
 &=k\E_{\theta_0}\varphi
   +\E_{\theta_1}[\varphi\ind_{\{M\ge\theta_0\}}]\\
 &\le k\alpha+1-k
 =1-(1-\alpha)(\theta_0/\theta_1)^n.
\end{align*}
令 $c_\alpha=\theta_0(1-\alpha)^{1/n}$，则
\[
 \varphi_A=\ind_{\{M>c_\alpha\}},\qquad
 \varphi_B=\ind_{\{M>\theta_0\}}+
              \alpha\ind_{\{M\le\theta_0\}}
\]
在 $\theta_0$ 下大小都为 $\alpha$，在 $\theta_1$ 下功效都达到上述上界。它们在原假设支持内采取不同安排，却没有改变任何相对证据的排序，因为该处似然比完全相同。
\end{examplebox}

\begin{examplebox}{练习 5.3：最优性与预算例外}
设 $P_0=U(0,1)$，$P_1=U(0,\delta)$，其中 $0<\delta<\alpha$。证明 $\varphi=\ind_{\{X<\delta\}}$ 最优，但大小严格小于 $\alpha$。再指出 NP 引理的哪个条件允许这一现象。

\textbf{答案线索。} 大小为 $\delta$，功效为 1，已达到上限；可取密度差阈值 $c=0$，所以不能使用 $c>0$ 时的精确大小结论。
\end{examplebox}

\section{单调似然比与单侧 UMP}

\subsection{单调似然比与随机单调性}

\begin{definitionbox}{单调似然比（MLR）}
对任意 $\theta_1<\theta_2$，若 $f_{\theta_2}/f_{\theta_1}$ 在相关支持上可写成 $T$ 的非递减函数，则称模型在 $T$ 中有 MLR；分母为零、分子为正时按 $+\infty$ 处理。其交叉乘积形式要求：若 $T(x)\le T(y)$，则
\[
 f_{\theta_2}(y)f_{\theta_1}(x)
 \ge f_{\theta_2}(x)f_{\theta_1}(y).
\]
各参数对的密度可选择适当版本，使上述比较在相关支持上成立。当 $T(x)=T(y)$ 时须同时作两个方向的比较，因而交叉乘积相等，保证同一个 $T$ 值对应同一个扩展似然比；只要求严格小于时的不等式不足以表达这一点。它表示较大的 $T$ 相对更支持较大的参数。共同支持是常用的充分背景条件，并非 MLR 的必要条件。
\end{definitionbox}

\textbf{MLR 推出随机单调性的证明。}
在 $A=\{T\le t\}$、$B=\{T>t\}$ 上积分交叉乘积不等式，Tonelli 定理给出
\[
 P_{\theta_1}(A)P_{\theta_2}(B)
 \ge P_{\theta_2}(A)P_{\theta_1}(B).
\]
令 $p_i=P_{\theta_i}(T>t)$，化为 $(1-p_1)p_2\ge(1-p_2)p_1$，故 $p_2\ge p_1$。同理 $P_\theta(T\ge t)$ 也随参数非递减。因此对固定 $c,\gamma$，
\[
 \E_\theta[\ind_{\{T>c\}}+\gamma\ind_{\{T=c\}}]
 =(1-\gamma)P_\theta(T>c)+\gamma P_\theta(T\ge c)
\]
随参数非递减。这是分布的尾概率比较，不是不同参数下样本值的逐点比较。

对共同支持上的指数族
\[
 f_\theta(x)=\exp\{\eta(\theta)T(x)-A(\theta)\}h(x),
\]
若 $\eta$ 非递减，则两参数的似然比为
\[
 \exp\{[\eta(\theta_2)-\eta(\theta_1)]T
                  -A(\theta_2)+A(\theta_1)\},
\]
所以在 $T$ 中有 MLR。$\eta$ 递减时应改用 $-T$。Cauchy 位置族在 $T=X$ 中不是 MLR：其两参数似然比在正、负无穷处都趋于 1，却在两个位置参数处分别小于和大于 1，不可能单调。

\subsection{Karlin--Rubin 定理及证明}

\begin{definitionbox}{单侧 UMP 定理}
若模型在 $T$ 中有 MLR，取
\[
 \varphi^*=\ind_{\{T>c\}}+\gamma\ind_{\{T=c\}},
 \qquad \E_{\theta_0}\varphi^*=\alpha\in(0,1).
\]
则它是 $H_0:\theta\le\theta_0$ 对 $H_1:\theta>\theta_0$ 的 UMP 水平 $\alpha$ 检验；对点原假设 $\theta=\theta_0$ 也成立。其功效随参数非递减。
\end{definitionbox}

\textbf{证明。}
随机单调性给出
\[
 \sup_{\theta\le\theta_0}\E_\theta\varphi^*
 =\E_{\theta_0}\varphi^*=\alpha,
\]
所以复合原假设水平确实受控。固定任意 $\theta>\theta_0$，MLR 保证上尾规则按 $f_\theta/f_{\theta_0}$ 的大小拒绝：在比值平台上的部分选择属于 NP 等号边界的允许安排。因此它是 $P_{\theta_0}$ 对 $P_\theta$ 的 MP 检验。每个复合原假设水平检验 $\psi$ 都满足 $\E_{\theta_0}\psi\le\alpha$，由 NP 引理
\[
 \E_\theta\varphi^*\ge\E_\theta\psi.
\]
临界值只由边界分布和 $\alpha$ 决定，不依赖所选的具体备择，故对每个 $\theta>\theta_0$ 都成立，得到 UMP。

同一规则若在另一参数 $\theta'$ 处的功效为 $a\in(0,1)$，重新以 $\theta'$ 为边界，就成为水平 $a$ 的右侧 UMP。对 $\theta<\theta_0$，补检验 $1-\varphi^*$ 是从边界检验较小参数的 NP 最优规则，所以在所有满足 $\E_{\theta_0}\psi=\alpha$ 的规则中，
\[
 \E_\theta\varphi^*\le\E_\theta\psi.
\]
即该上尾规则在较小参数处还具有最小拒绝概率。

若各参数分布具有共同正支持，$0<\alpha<1$ 保证所有参数处功效均严格介于 0 与 1。若两个参数的似然比不是常数，功效还严格增加。否则令 $a=\beta(\theta_1)=\beta(\theta_2)$，常数检验 $a$ 在该简单检验问题中也最优；NP 等号条件却要求它在严格大于阈值处等于 1、严格小于阈值处等于 0。因 $0<a<1$，只能是似然比几乎处处恒定，矛盾。

\subsection{例子与方向}

\begin{examplebox}{Poisson、二项与指数分布}
若 $X_i\sim\operatorname{Poisson}(\theta)$ 独立，则 $S=\sum X_i\sim\operatorname{Poisson}(n\theta)$，自然参数 $\log\theta$ 递增。检验 $\theta\le\theta_0$ 对 $\theta>\theta_0$ 时选
\[
 P_{\theta_0}(S>c)\le\alpha\le P_{\theta_0}(S\ge c),\qquad
 \gamma=\frac{\alpha-P_{\theta_0}(S>c)}{P_{\theta_0}(S=c)}.
\]
按大 $S$ 拒绝即为 UMP。二项族的似然比同样随成功数递增，前节 $\operatorname{Bin}(4,1/2)$ 的边界规则因而对所有 $p>1/2$ 同时最优。

若 $X_i\sim\operatorname{Exp}(\lambda)$，以率参数 $\lambda$ 表示，则密度中的自然参数是 $-\lambda$。较大的率对应较短的等待时间，检验较大 $\lambda$ 应拒绝较小的 $\sum X_i$，方向不能照搬。
\end{examplebox}

均匀端点族在 $M$ 中有 MLR。前节的 $\varphi_A,\varphi_B$ 对所有 $\theta>\theta_0$ 都达到 NP 功效上界，又都满足复合原假设水平，因此都是 UMP。两者完整功效为
\begin{align*}
 \beta_A(\theta)&=
 \begin{cases}0,&\theta\le c_\alpha,\\1-(c_\alpha/\theta)^n,&\theta>c_\alpha,
 \end{cases}\\[3pt]
 \beta_B(\theta)&=
 \begin{cases}\alpha,&\theta\le\theta_0,\\
 1-(1-\alpha)(\theta_0/\theta)^n,&\theta>\theta_0.
 \end{cases}
\end{align*}
备择一侧功效相同，不要求原假设一侧也相同；支持随参数改变时，某些参数处功效可为零。

\begin{examplebox}{练习 5.4：把物理参数换成自然参数}
对率参数指数样本，检验 $\lambda\le\lambda_0$ 对 $\lambda>\lambda_0$。利用 $2\lambda_0\sum X_i\sim\chi^2_{2n}$ 写出 UMP 规则，并说明最坏原假设参数。

\textbf{答案线索。} 以 $-\sum X_i$ 为 MLR 统计量，拒绝 $2\lambda_0\sum X_i<\chi^2_{2n,\alpha}$。拒绝率随 $\lambda$ 增大，原假设最大值在 $\lambda_0$ 处取得；连续分布不需要随机化。
\end{examplebox}

\section{双侧困难与无偏约束}

\subsection{正态双侧问题不存在 UMP}

设 $Z\sim N(\delta,1)$，检验 $\delta=0$ 对 $\delta\ne0$。似然比
\[
 f_\delta(z)/f_0(z)=e^{\delta z-\delta^2/2}
\]
在 $\delta>0$ 时严格递增，在 $\delta<0$ 时严格递减。因连续分布的边界概率为零，NP 引理的唯一性给出：对任意固定正备择，唯一 MP 规则是
$\varphi_+=\ind_{\{Z>z_{1-\alpha}\}}$；对任意固定负备择，唯一 MP 规则是
$\varphi_-=\ind_{\{Z<z_\alpha\}}$。

若双侧 UMP 存在，它必须在 $P_0$ 下同时等于这两个规则。但在
\[
 \{Z>\max(z_\alpha,z_{1-\alpha})\}
\]
上，前者拒绝、后者不拒绝，该集合的原假设概率为 $\min(\alpha,1-\alpha)>0$，矛盾。此证明对所有 $0<\alpha<1$ 成立。把左右两尾各改为 $\alpha/2$ 只能控制总大小，不能恢复对每个固定单侧备择的 NP 最优性。

\subsection{无偏性给出的两个矩方程}

\begin{definitionbox}{无偏检验与 UMPU}
水平 $\alpha$ 检验若对每个备择满足 $\beta_\varphi(\theta)\ge\alpha$，称为无偏检验。若它在所有无偏水平 $\alpha$ 检验中对每个备择都最有力，称为 UMPU。常数规则 $\varphi\equiv\alpha$ 是比较基线。检验无偏性不是估计量期望等于参数。
\end{definitionbox}

设 $\theta_0$ 为参数区间内点，功效在该点连续可微，并允许对固定检验在积分号下求导。无偏性与水平条件先给出
\[
 \beta(\theta_0)=\alpha.
\]
因为邻近备择的功效至少为 $\alpha$，连续性给出边界处至少为 $\alpha$，水平又给出不超过 $\alpha$。又因该点是功效的局部最小点，$\beta'(\theta_0)=0$。

令原假设得分 $S_0=\dot f_{\theta_0}/f_{\theta_0}$，由密度归一化求导及功效求导，
\[
 \E_0S_0=0,\qquad
 \E_0\varphi=\alpha,\qquad \E_0(\varphi S_0)=0.
\]
共同正支持、导数绝对可积及相应的积分求导条件需实际核验；含参数支持边界的模型不能直接套用。

在正则自然指数族 $f_\eta=e^{\eta T-A(\eta)}h$ 中，
$S_0=T-m$，其中 $m=\E_0T=A'(\eta_0)$，故必要条件化为
\[
 \boxed{\E_0\varphi=\alpha,\qquad \E_0(T\varphi)=\alpha m.}
\]
第一个方程固定总拒绝概率；第二个要求拒绝区域内的正负中心化一阶矩平衡。它们只是无偏性的必要局部条件，候选规则的全局无偏性仍要证明。

对连续双尾规则，第二个方程可写为
\[
 \int_{t<c_1}(m-t)\,dP_0^T(t)
 =\int_{t>c_2}(t-m)\,dP_0^T(t).
\]
它平衡的是距均值的加权概率，不是两尾面积。对称分布常得到等尾，偏态分布通常不是每尾 $\alpha/2$。

\begin{examplebox}{练习 5.5：必要条件为什么不充分}
设 $Z\sim N(\delta,1)$。选择 $c>0$ 使 $P_0(|Z|\le c)=\alpha$，令 $\varphi=\ind_{\{|Z|\le c\}}$。验证它满足大小和原假设处的导数条件，却不是双侧无偏检验。

\textbf{答案线索。} 对称性给出 $\E_0(Z\varphi)=0$，大小为 $\alpha$；但 $|\delta|\to\infty$ 时 $P_\delta(|Z|\le c)\to0<\alpha$。在原点导数为零并不能控制全部备择。
\end{examplebox}

\section{广义 Neyman--Pearson 引理}

\subsection{多个矩约束下的符号规则}

设 $g,g_1,\ldots,g_k$ 对 $\sigma$-有限测度 $\nu$ 绝对可积，允许取正负值。记
\[
 \Phi=\{\varphi:0\le\varphi\le1,\ \varphi\text{ 可测}\},\qquad
 \mathcal C_b=\left\{\varphi\in\Phi:
       \int\varphi g_j\,d\nu=b_j,\ j=1,\ldots,k\right\}.
\]
目标是在非空的 $\mathcal C_b$ 中最大化 $\int\varphi g\,d\nu$。

\begin{definitionbox}{广义 NP 引理}
最大值可以取到。若某个可行 $\varphi^*$ 和有限实数乘子 $a_j$ 满足：令
\[
 \Delta=g-\sum_{j=1}^k a_jg_j,
\]
则 $\varphi^*=1$ 于 $\{\Delta>0\}$、$\varphi^*=0$ 于 $\{\Delta<0\}$，它就最优。等号边界可随机化。

进一步令
\[
 M=\left\{\left(\int\varphi g_j\,d\nu\right)_{j=1}^k:
                    \varphi\in\Phi\right\}.
\]
若 $b$ 为 $M$ 的内点，则存在有限实数乘子，使所有最优解都具有上述符号形式，几乎处处成立。
\end{definitionbox}

\textbf{充分性的证明。}
对任意 $\psi\in\mathcal C_b$，逐点有 $(\varphi^*-\psi)\Delta\ge0$，而矩约束相同，故
\begin{align*}
 \int(\varphi^*-\psi)g\,d\nu
 &=\int(\varphi^*-\psi)\Delta\,d\nu
   +\sum_j a_j\int(\varphi^*-\psi)g_j\,d\nu\\
 &=\int(\varphi^*-\psi)\Delta\,d\nu\ge0.
\end{align*}
相同功效时，非负乘积几乎处处为零，故 $\Delta\ne0$ 处具有同样的唯一性。

若约束改成 $\int\varphi g_j\,d\nu\le b_j$，充分条件相应为 $a_j\ge0$、符号规则、可行性和互补松弛
\[
 a_j\left(b_j-\int\varphi^*g_j\,d\nu\right)=0.
\]
此时上式中乘子项等于 $\sum_j a_j(b_j-\int\psi g_j\,d\nu)\ge0$。等式约束允许任意实数乘子，不等式上界约束才要求非负。

\subsection{存在性与逆结论的证明}

这部分使用两个标准分析结果：Hilbert 空间有界序列存在弱收敛子序列；有限维闭凸集的边界点存在支撑超平面。它们分别保证最优解取到，以及有限乘子的存在。

记录每个检验的矩向量和目标值：
\[
 u_\varphi=\left(\int\varphi g_j\,d\nu\right)_j,
 \quad v_\varphi=\int\varphi g\,d\nu,
 \quad K=\{(u_\varphi,v_\varphi):\varphi\in\Phi\}.
\]
$K$ 显然凸。为证紧，利用 $\nu$ 的 $\sigma$-有限性选取严格正的可积 $q$，设
\[
 w=q+|g|+\sum_j|g_j|,\qquad
 C=\int w\,d\nu,\qquad d\mu=(w/C)\,d\nu.
\]
$\mu$ 是与 $\nu$ 等价的概率测度。每个检验的 $L^2(\mu)$ 范数不超过 1，所以任意检验序列有弱收敛子序列 $\varphi_{n_\ell}\rightharpoonup\varphi$。对任意可测集 $A$，以 $\ind_A$ 为测试函数，有
\[
 0\le\int_A\varphi\,d\mu\le\mu(A).
\]
分别取 $A=\{\varphi<0\}$ 和 $A=\{\varphi>1\}$，得 $0\le\varphi\le1$ 几乎处处。又因 $Cg_j/w$、$Cg/w$ 有界，均属于 $L^2(\mu)$，弱收敛给出
\[
 \int\varphi_{n_\ell}g_j\,d\nu\longrightarrow\int\varphi g_j\,d\nu,
 \qquad
 \int\varphi_{n_\ell}g\,d\nu\longrightarrow\int\varphi g\,d\nu.
\]
所以每列可达向量都有仍可达的收敛子序列，$K$ 在有限维空间中紧。选可行检验的目标值趋于上确界，紧性给出某个可行检验达到最大值 $v^*$。

再定义向下延伸的闭凸集
\[
 H=\{(u,v-t):(u,v)\in K,\ t\ge0\}.
\]
闭性由 $K$ 的紧性直接验证。最优点 $(b,v^*)$ 在其上边界，故存在非零 $(p,s)$，使
\[
 p^\top(u-b)+s(v-v^*)\le0,\qquad (u,v)\in H.
\]
因 $H$ 可向下无限延伸，必有 $s\ge0$。若 $s=0$，则 $p\ne0$；$b\in\operatorname{int}M$ 允许取 $u=b+\varepsilon p\in M$，将对应点代入便得 $\varepsilon\|p\|^2\le0$，矛盾。因此 $s>0$。令 $a=-p/s$，便有
\[
 v_\varphi-a^\top u_\varphi\le v^*-a^\top b,
 \qquad\varphi\in\Phi.
\]
最优解达到等号，因此它在全部 $\Phi$ 中最大化
$\int\varphi\Delta\,d\nu$。这个无约束积分逐点最大化时，正系数处必须取 1，负系数处必须取 0，得到逆结论。内点条件正是排除 $s=0$ 的竖直支撑面。

\begin{examplebox}{练习 5.6：从两个约束得到一条直线}
在 $\E_0\varphi=\alpha$、$\E_0(T\varphi)=\alpha m$ 下最大化 $\E_\eta\varphi$。写出广义 NP 的 $g,g_1,g_2$ 与密度差，并解释乘子为何不需要非负。

\textbf{答案线索。} 取 $g=f_\eta$、$g_1=f_0$、$g_2=Tf_0$，于是 $\Delta=f_\eta-(a+bT)f_0$。原约束为等式，故 $a,b$ 为任意有限实数。
\end{examplebox}

\section{一参数指数族的双侧 UMPU}

\subsection{定理、可行性与双尾结构}

考虑共同正支持上的自然指数族
\[
 f_\eta(x)=\exp\{\eta T(x)-A(\eta)\}h(x),\qquad \eta\in I,
\]
其中 $I$ 为开区间，$\eta_0\in I$，$T$ 非恒定，$0<\alpha<1$。开自然参数区间保证在 $\eta_0$ 两侧有指数矩，从而可交换局部求导与积分，并有 $m=\E_0T=A'(\eta_0)$。

\begin{definitionbox}{一参数指数族双侧 UMPU}
检验 $\eta=\eta_0$ 对 $\eta\ne\eta_0$，存在满足
\[
 \E_0\varphi^*=\alpha,\qquad
 \E_0(T\varphi^*)=\alpha m
\]
的双尾规则：
\[
 \varphi^*(x)=
 \begin{cases}
 1,&T<c_1\text{ 或 }T>c_2,\\
 \gamma_i,&T=c_i,\quad i=1,2,\\
 0,&c_1<T<c_2.
 \end{cases}
\]
每个这样标定的双尾规则都是 UMPU。若临界点重合，使用一个边界随机化值；若 $T$ 恒定，模型不随 $\eta$ 改变，常数规则 $\alpha$ 即可。
\end{definitionbox}

\textbf{第一步：验证广义 NP 的内点条件。}
设 $\mathcal C_b$ 为上述两矩约束类，常数规则 $\alpha$ 可行。考虑
\[
 \varphi_{s,t}=\alpha+s+t\ind_{\{T>m\}},\qquad
 |s|+|t|<\min(\alpha,1-\alpha).
\]
它是合法检验，矩向量的改变量为
\[
 \begin{pmatrix}
 1&P_0(T>m)\\m&\E_0[T\ind_{\{T>m\}}]
 \end{pmatrix}
 \begin{pmatrix}s\\t\end{pmatrix}.
\]
矩阵行列式为 $\E_0[(T-m)\ind_{\{T>m\}}]>0$，因为非恒定可积变量不可能仅在均值的一侧取值。因此小扰动覆盖 $b=(\alpha,\alpha m)^\top$ 附近的二维邻域，内点条件成立。

\textbf{第二步：由严格凸性得到双尾。}
固定 $\eta\ne\eta_0$，广义 NP 保证有最优解，且按
\[
 d_\eta(t)=r_\eta(t)-(a+bt),\qquad
 r_\eta(t)=e^{(\eta-\eta_0)t-A(\eta)+A(\eta_0)}
\]
的正负拒绝。因
\[
 d_\eta''(t)=(\eta-\eta_0)^2r_\eta(t)>0,
\]
负值区间是一个区间，规则在其外拒绝，内部接受，至多在两个零点随机化。

单尾情形不能满足两个矩约束。因为对任何常数 $c$ 都有
\[
 \E_0[(T-c)(\varphi-\alpha)]=0.
\]
若 $\varphi$ 为上尾规则，则乘积在 $T<c$ 时为 $\alpha(c-T)>0$，在 $T>c$ 时为 $(1-\alpha)(T-c)>0$，在 $T=c$ 时为零。非恒定性使其期望严格为正，矛盾；下尾规则则严格为负。恒为 0 或 1 也不符合大小约束。因此可取双尾形式，包含在有限支持边界上随机化和两个临界点重合的情形。

\subsection{同一个规则对全部备择都最优}

\textbf{第三步：固定临界点，用割线逐一比较。}
固定已得到的两矩可行双尾规则，设 $c_1<c_2$。对任意新的备择 $\eta\ne\eta_0$，连接指数曲线上的
$(c_1,r_\eta(c_1))$、$(c_2,r_\eta(c_2))$，得到直线 $\ell_\eta(t)=a_\eta+b_\eta t$。严格凸性给出
\[
 r_\eta(t)-\ell_\eta(t)
 \begin{cases}
 >0,&t<c_1\text{ 或 }t>c_2,\\
 <0,&c_1<t<c_2,\\
 =0,&t=c_1,c_2.
 \end{cases}
\]
区间内曲线低于割线；区间外的严格不等式由割线斜率的单调性得到。例如 $t<c_1$ 时，
\[
 \frac{r_\eta(c_1)-r_\eta(t)}{c_1-t}
 <\frac{r_\eta(c_2)-r_\eta(c_1)}{c_2-c_1}=b_\eta,
\]
移项即得曲线高于延长线。于是固定的 $\varphi^*$ 对每个备择都满足广义 NP 的充分符号规则，所以
\[
 \E_\eta\varphi^*\ge\E_\eta\psi,
 \qquad\psi\in\mathcal C_b,\quad\eta\ne\eta_0.
\]
改变备择时，改变的是割线的斜率和截距，临界点与检验本身均不变。

\textbf{第四步：证明全局无偏并限制比较类。}
取 $\psi\equiv\alpha$ 得 $\E_\eta\varphi^*\ge\alpha$，证明候选确实无偏。每个无偏水平检验又都满足前文两个必要矩方程，因此属于 $\mathcal C_b$，上述功效比较覆盖所有无偏竞争者，UMPU 得证。

若 $c_1=c_2=c$，将割线改为 $r_\eta$ 在 $c$ 处的切线。严格凸性保证除 $c$ 外差均为正，规则因而在 $T\ne c$ 全拒绝，在 $T=c$ 随机化。记 $p=P_0(T=c)$，矩方程给出
\[
 (1-\gamma)p=1-\alpha,
 \qquad m-(1-\gamma)cp=\alpha m.
\]
所以 $c=m$ 且 $\gamma=1-(1-\alpha)/p$，要求 $p\ge1-\alpha$。连续分布没有这样的重合情形。

\begin{examplebox}{已知方差的双侧正态检验}
对独立正态样本 $N(\mu,\sigma^2)$，$\sigma$ 已知，检验 $\mu=\mu_0$。令 $Z=\sqrt n(\bar X-\mu_0)/\sigma$，则
\[
 \varphi^*=\ind_{\{|Z|>z_{1-\alpha/2}\}}
\]
在原假设下大小为 $\alpha$。又因中心化正态分布对称，
\[
 \E_{\mu_0}[(\bar X-\mu_0)\varphi^*]=0,
\]
所以两个矩方程都成立，双侧 $z$ 检验为 UMPU。这与双侧不存在 UMP 不矛盾：两者的竞争类不同。
\end{examplebox}

\begin{examplebox}{练习 5.7：两尾为什么不一定等概率}
设 $T$ 在原假设下服从偏态连续分布，均值为 $m$。分别写出双侧 UMPU 临界值的面积方程与一阶矩方程。解释仅令两尾各为 $\alpha/2$ 为何通常不足。

\textbf{答案线索。} 要求 $P_0(T<c_1)+P_0(T>c_2)=\alpha$ 及
$\E_0[(T-m)\ind_{\{T<c_1\text{ 或 }T>c_2\}}]=0$。尾概率未包含距均值的权重，等面积一般不能推出矩平衡。
\end{examplebox}

\section{干扰参数与条件 UMPU}

\subsection{相似性、Neyman 结构与完备性}

考虑参数 $(\eta,\lambda)$，$\eta$ 为研究对象，$\lambda$ 为干扰参数。固定 $\eta=\eta_0$ 的边界族记为
$\mathcal P_b=\{P_{\eta_0,\lambda}:\lambda\in\Lambda\}$。

\begin{definitionbox}{相似检验与 Neyman 结构}
若对每个 $\lambda$ 都有 $\E_{\eta_0,\lambda}\varphi=\alpha$，称检验在边界族上相似。设 $U$ 对边界族充分，若
\[
 \E_{\eta_0,\lambda}(\varphi\mid U)=\alpha
 \quad P_{\eta_0,\lambda}\text{-几乎处处，对每个 }\lambda,
\]
则称具有 Neyman 结构。塔式法则立即给出后者蕴含相似性。
\end{definitionbox}

若 $U$ 对边界族充分，则以下两项等价：$U$ 有界完备；每个大小 $\alpha\in(0,1)$ 的相似检验都具有 Neyman 结构。

\textbf{证明。}
充分性允许选取对所有 $\lambda$ 共用的条件期望函数。对相似检验令
\[
 h(U)=\E_{\eta_0,\lambda}(\varphi\mid U)-\alpha.
\]
它有界，且对每个 $\lambda$ 期望为零，有界完备性推出 $h=0$，即 Neyman 结构。反之，若有界完备性失败，存在有界 $q(U)$，对所有 $\lambda$ 期望为零，但并非在每个分布下几乎处处为零。若 $|q|\le M$，取
\[
 0<\varepsilon<\min(\alpha,1-\alpha)/M.
\]
则 $\varphi=\alpha+\varepsilon q(U)$ 是相似检验，但它本身为 $U$ 的函数，故条件期望就是自身，不具有 Neyman 结构，矛盾。

\subsection{条件似然比与两个条件矩}

设模型为共同正支持上的正则指数族
\[
 f_{\eta,\lambda}(x)
 =\exp\{\eta T(x)+\lambda^\top U(x)-A(\eta,\lambda)\}h(x),
\]
其中 $(\eta,\lambda)\in I\times\Lambda$，$I$ 为含 $\eta_0$ 的开区间，$\Lambda$ 为开集；假设存在通常样本空间上的正则条件分布。开自然参数邻域保证所需的局部积分求导。边界族中 $U$ 充分。给定 $U=u$ 后，$\lambda$ 的因子在归一化中消去，条件似然比为
\[
 \frac{dQ_{\eta,u}}{dQ_{\eta_0,u}}
 =\frac{\exp\{(\eta-\eta_0)T\}}
 {\E_{\eta_0}[\exp\{(\eta-\eta_0)T\}\mid U=u]}.
\]
右侧不含干扰参数，条件问题成为一参数指数倾斜。分母必须保留，它使条件密度积分为 1。

任取全模型中的无偏水平检验 $\psi$。固定 $\lambda$，由原假设处功效的连续性与可微性，
\[
 \E_{\eta_0,\lambda}\psi=\alpha,\qquad
 \E_{\eta_0,\lambda}[T(\psi-\alpha)]=0.
\]
充分性允许选择共同函数
\[
 h(U)=\E_0(\psi-\alpha\mid U),\qquad
 g(U)=\E_0[T(\psi-\alpha)\mid U].
\]
二者对所有 $\lambda$ 的期望为零。$h$ 有界，有界完备性足以消去；$g$ 未必有界，仅有
\[
 \E_0|g(U)|\le\E_0|T|<\infty.
\]
所以要把第二个约束也传入条件层，需要完整的完备性。若 $U$ 对边界族完备充分，便有
\[
 \boxed{\E_0(\psi\mid U)=\alpha,\qquad
 \E_0(T\psi\mid U)=\alpha\E_0(T\mid U).}
\]
这里只为第一个条件使用有界完备性，不能据此省略第二个条件的可积性和完备性要求。

\begin{definitionbox}{条件 UMPU 定理}
在上述正则共同正支持模型中，取 $0<\alpha<1$，设 $U$ 对边界族完备充分。若可测的临界值及端点随机化给出逐层双尾规则 $\varphi^*$，并满足
\[
 \E_0(\varphi^*\mid U)=\alpha,\qquad
 \E_0(T\varphi^*\mid U)=\alpha\E_0(T\mid U),
\]
则 $\varphi^*$ 是 $\eta=\eta_0$ 对 $\eta\ne\eta_0$、干扰参数不受限制时的 UMPU 检验。条件层上 $T$ 恒定时用常数规则 $\alpha$。
\end{definitionbox}

\textbf{证明。}
任意无偏竞争者 $\psi$ 都满足上述两个条件矩。条件似然比不含干扰参数，故在几乎每个 $U=u$ 层内，可以用一参数结果的割线证明得到
\[
 \E_{\eta,\lambda}(\varphi^*-\psi\mid U=u)\ge0.
\]
共同正支持使原假设与备择的分布相互绝对连续，因而它们的 $U$ 分布也具有相同零测集，原假设下忽略的条件层不会在备择下突然有正概率。按\emph{备择下} $U$ 的分布平均，
\[
 \E_{\eta,\lambda}(\varphi^*-\psi)
 =\E_{\eta,\lambda}
       [\E_{\eta,\lambda}(\varphi^*-\psi\mid U)]\ge0.
\]
候选的条件大小给出每个干扰参数处的大小为 $\alpha$；再取也满足条件矩的常数规则 $\psi\equiv\alpha$，得到候选在每个备择下功效至少为 $\alpha$。它自身无偏，且优于全部无偏竞争者，结论成立。临界值的可测性是拼成合法检验的必要条件，不能只逐个条件层抽象选取而不加说明。

\subsection{未知方差的正态均值检验}

设 $X_i\sim N(\mu,\sigma^2)$ 独立，$n\ge2$，检验 $\mu=\mu_0$。令
\[
 Y=\bar X-\mu_0,\qquad U=\sum_i(X_i-\mu_0)^2.
\]
平方展开后，联合密度可写成
\[
 f_{\mu,\sigma^2}(x)=
 (2\pi\sigma^2)^{-n/2}e^{-n(\mu-\mu_0)^2/(2\sigma^2)}
 \exp\left\{\frac{n(\mu-\mu_0)}{\sigma^2}Y
                         -\frac{U}{2\sigma^2}\right\}.
\]
目标自然参数 $\eta=n(\mu-\mu_0)/\sigma^2\in\R$，干扰自然参数 $\lambda=-1/(2\sigma^2)<0$；原假设等价于 $\eta=0$。在边界族中 $U/\sigma^2\sim\chi_n^2$，因子分解与开自然参数区间给出 $U$ 对未知方差完备充分。

原假设下同时把各 $X_i-\mu_0$ 变号，$U$ 不变而 $Y$ 变为 $-Y$，故给定 $U$ 时 $Y$ 关于零对称。因此按 $|Y|>c(U)$ 拒绝，只要条件大小为 $\alpha$，就自动满足第二个条件矩。

令 $V_0=Y/\sqrt U$。边界族下写 $X_i-\mu_0=\sigma Z_i$，则
\[
 V_0=\frac{\bar Z}{\sqrt{\sum_i Z_i^2}},
\]
分布不依赖 $\sigma^2$，所以它对\emph{边界族}辅助。Basu 定理给出 $V_0\perp U$，可取常数 $d$ 使 $P_0(|V_0|>d)=\alpha$，则
\[
 P_0(|V_0|>d\mid U)=\alpha,
 \qquad c(u)=d\sqrt u.
\]

设 $S^2=(n-1)^{-1}\sum_i(X_i-\bar X)^2$。由 $U=(n-1)S^2+nY^2$，有
\[
 t=\frac{\sqrt nY}{S}
 =\frac{\sqrt{n(n-1)}V_0}{\sqrt{1-nV_0^2}}.
\]
$|V_0|<1/\sqrt n$ 几乎处处成立，右侧是严格递增的奇函数，其导数为
$\sqrt{n(n-1)}(1-nV_0^2)^{-3/2}$。故 $|V_0|>d$ 等价于 $|t|>k$。正态均值方向与正交残差独立，给出原假设下 $t\sim t_{n-1}$，从而
\[
 \boxed{\varphi^*=
 \ind_{\{|\sqrt n(\bar X-\mu_0)/S|>t_{n-1,1-\alpha/2}\}}}
\]
是对所有 $\sigma^2>0$ 都成立的双侧 UMPU 检验。Basu 在这里用于固定 $\mu_0$ 的边界族，不是整个双参数模型。

\subsection{两个 Poisson 强度的条件检验}

设 $X_j\sim\operatorname{Poisson}(\rho_j)$ 独立，$\rho_1,\rho_2>0$，检验 $\rho_1=\rho_2$。令
\[
 T=X_1,\quad U=X_1+X_2,\quad
 \eta=\log(\rho_1/\rho_2),\quad\lambda=\log\rho_2.
\]
联合概率函数为
\[
 p(x_1,x_2)=\frac{\exp\{\eta T+\lambda U-e^\lambda(1+e^\eta)\}}
 {x_1!x_2!}.
\]
给定 $U=u$，直接相除得到
\[
 X_1\mid U=u\sim\operatorname{Bin}\left(u,
               \frac{\rho_1}{\rho_1+\rho_2}\right).
\]
原假设下成功概率为 $1/2$，干扰强度消失，且 $U\sim\operatorname{Poisson}(2\rho)$ 对边界族完备充分。在每层选择关于 $u/2$ 对称的二项双尾规则，并作对称边界随机化达到条件大小 $\alpha$，就有
\[
 \E_0[(X_1-u/2)\varphi\mid U=u]=0.
\]
两个条件矩均成立，故得到强度不等的 UMPU 检验。$u=0$ 时条件统计量恒定；$u=1$ 时两个矩方程迫使两个可能取值处的拒绝概率均为 $\alpha$，这两层都用常数随机化。

\begin{examplebox}{练习 5.8：小计数层的精确检验}
取 $u=2$、$\alpha=0.05$，在 $X_1=0,2$ 时以同一概率 $\gamma$ 拒绝，$X_1=1$ 时不拒绝。求 $\gamma$，验证两个条件矩。

\textbf{答案线索。} 原假设条件概率为 $(1/4,1/2,1/4)$，故 $\gamma/2=0.05$，得到 $\gamma=0.10$；$\E_0(X_1\varphi\mid U=2)=2\gamma/4=0.05=\alpha\,\E_0(X_1\mid U=2)$。
\end{examplebox}

\section{检验反演与最优置信集合}

\subsection{覆盖率与错误候选的纳入概率}

对每个候选 $\theta_0$，给出检验该点原假设的非随机水平 $\alpha$ 规则 $\varphi_{\theta_0}$，定义
\[
 C(X)=\{\theta_0:\varphi_{\theta_0}(X)=0\}.
\]
即固定同一份数据，保留所有未被拒绝的参数候选。结果一般是集合，不必是一个区间。假设有关集合事件可测，则真值为 $\theta$ 时
\[
 P_\theta\{\theta\in C(X)\}
 =1-\E_\theta\varphi_\theta(X)\ge1-\alpha.
\]
若每个点原假设下的拒绝概率恰为 $\alpha$，覆盖率就精确为 $1-\alpha$。对另一个候选 $\theta_0\ne\theta$，
\[
 P_\theta\{\theta_0\in C(X)\}
 =1-\E_\theta\varphi_{\theta_0}(X).
\]
$\theta$ 决定数据的真分布，$\theta_0$ 决定被检验的候选，两个下标各有作用。

随机化检验先取统一的独立 $V\sim U(0,1)$，定义
$C(X,V)=\{\theta_0:V>\varphi_{\theta_0}(X)\}$。相同公式仍成立，但概率同时对数据和辅助随机数计算。

\subsection{UMP 与 UMA 的对偶}

一致最准确（UMA）在本章指：满足覆盖保证的前提下，对规定方向上的每个错误候选，其纳入概率最小。

\textbf{对偶证明。}
设 $D(X)$ 为另一满足覆盖率至少 $1-\alpha$ 的置信集合。对固定候选 $\theta_0$ 定义
\[
 \psi_{\theta_0}(X)=\ind_{\{\theta_0\notin D(X)\}}.
\]
覆盖保证给出 $\E_{\theta_0}\psi_{\theta_0}\le\alpha$，所以它是一个合法水平检验。若 $\varphi_{\theta_0}$ 在所考虑的备择方向为 UMP，则
\[
 \E_\theta\varphi_{\theta_0}\ge\E_\theta\psi_{\theta_0}.
\]
两边从 1 中减去，便有
\[
 P_\theta(\theta_0\in C)\le P_\theta(\theta_0\in D).
\]
最有力的拒绝转成了对错误候选最小的纳入概率。若比较对象本身随机化，用其条件拒绝概率代替上述指示函数，论证相同。

\begin{examplebox}{正态均值的单侧下置信界}
已知方差时，检验 $\mu=\mu_0$ 对 $\mu>\mu_0$ 的 UMP 规则在
$\sqrt n(\bar X-\mu_0)/\sigma>z_{1-\alpha}$ 时拒绝。接受候选的条件等价于
\[
 \mu_0\ge\bar X-z_{1-\alpha}\frac\sigma{\sqrt n},
\]
所以反演为
\[
 C(X)=\left[\bar X-z_{1-\alpha}\frac\sigma{\sqrt n},\ \infty\right).
\]
其 UMA 结论针对 $\mu_0<\mu$ 的错误候选，即对应检验的备择方向。不能据此声称它对真值两侧的所有错误候选都最准确。
\end{examplebox}

\subsection{UMPU 与无偏置信集合}

本章称置信集合无偏，若它对每个真参数有覆盖率至少 $1-\alpha$，并且对每个错误候选 $\theta_0\ne\theta$ 都有
\[
 P_{\theta,\lambda}(\theta_0\in C)\le1-\alpha.
\]
反演双侧无偏检验立即得到这个性质，因为错误候选对应备择下的拒绝概率至少为 $\alpha$。反之，对任意这样的竞争集合 $D$，规则
$\psi_{\theta_0}=\ind_{\{\theta_0\notin D\}}$ 既有水平保证，也在每个备择下功效至少为 $\alpha$，所以属于无偏检验类。

因此反演 UMPU 检验，得到在无偏置信集合类中对每个错误候选纳入概率都最小的集合，称为一致最准确无偏（UMAU）置信集合。这里的最优指标是逐个错误候选的纳入概率，不是笼统的“长度最短”。

\begin{examplebox}{双侧 $t$ 区间}
正态未知方差时，接受候选 $\mu_0$ 的条件是
\[
 \left|\frac{\sqrt n(\bar X-\mu_0)}S\right|
 \le t_{n-1,1-\alpha/2}.
\]
故反演得到
\[
 C(X)=\left[\bar X-t_{n-1,1-\alpha/2}\frac S{\sqrt n},\quad
                 \bar X+t_{n-1,1-\alpha/2}\frac S{\sqrt n}\right].
\]
对所有 $\mu\in\R$、$\sigma^2>0$，其覆盖率都恰为 $1-\alpha$；由条件 UMPU 与反演对偶，它在无偏置信集合类中具有 UMAU 性质。已知方差时将 $t$ 分位数与 $S$ 分别换成 $z$ 分位数与 $\sigma$，结论相同。

覆盖率描述固定真参数下的重复抽样行为。观测数据后的区间端点已固定，频率覆盖结论本身不赋予参数一个后验分布。
\end{examplebox}

\begin{examplebox}{练习 5.9：反演的方向与最优性的含义}
从已知方差正态均值检验 $H_0:\mu=\mu_0$ 对 $H_1:\mu<\mu_0$ 出发，求单侧置信集合，并说明其 UMA 性质针对哪一侧错误候选。再说明双侧 $t$ 区间的 UMAU 性质为何不能直接译成“所有置信区间中长度最短”。

\textbf{答案线索。} 接受条件为 $\sqrt n(\bar X-\mu_0)/\sigma\ge z_\alpha$，所以
$C=(-\infty,\bar X+z_{1-\alpha}\sigma/\sqrt n]$；比较针对 $\mu_0>\mu$。UMAU 限于无偏集合类，比较的是错误候选的纳入概率，并非直接比较每次样本的长度或期望长度。
\end{examplebox}

\section{本章结果的使用条件}

\begin{center}
\small
\begin{tabularx}{\textwidth}{@{}p{0.23\textwidth}XX@{}}
\toprule
问题 & 主要工具 & 必须核验的条件\\
\midrule
简单对简单 & NP 引理 & 大小、似然比边界与随机化\\
单侧复合假设 & MLR 与单侧 UMP & 自然参数方向、边界水平、规则不随备择改变\\
一参数双侧问题 & 广义 NP 与 UMPU & 正则求导、两个矩方程、全局无偏\\
存在干扰参数 & 条件 UMPU & 边界族完备充分、条件两矩、可测性与零测集转移\\
构造置信集合 & 检验反演 & 覆盖率、候选与真值的下标、比较方向和集合类\\
\bottomrule
\end{tabularx}
\end{center}

一项功效比较只有同时说明假设、水平和竞争类，结论才完整。NP 在全部水平检验中处理固定备择；MLR 使同一排序适用于所有单侧备择；双侧问题借无偏性限定竞争者，再以矩约束和凸性完成统一比较。条件化与反演保留的正是这些约束，因此最优性名称不能脱离适用范围单独使用。

\appendix
\renewcommand{\chaptermark}[1]{\markboth{附录 \thechapter}{}}
\chapter{充分性与分布变换习题}

本组对应第一次作业第 7--12 题。第 1--6 题在本轮学习讨论中没有完整解答，不计入已整理范围。以下重述数学问题并给出解答；几乎必然等式均指对每个模型分布分别成立，所用统计量和表示函数不含未知参数。

\section{柯西位置尺度族的最小充分统计量}
设 $X_1,\ldots,X_n$ 独立同分布于 $C(\mu,\sigma)$，$\mu\in\R$、$\sigma>0$ 均未知。证明全部次序统计量组成的向量最小充分。

\textbf{解。} 联合密度处处为正，且
\[
 p_{\mu,\sigma}(x)=\frac{\sigma^n}{\pi^n\prod_i[(x_i-\mu)^2+\sigma^2]}.
\]
排序不改变乘积，所以因子分解给出 $T=(X_{(1)},\ldots,X_{(n)})$ 充分。若两条似然成比例，令 $\sigma\to\infty$，两边分母同为首项系数 $1$ 的 $\sigma^{2n}$ 多项式，比例常数必为 $1$。于是对所有 $\mu\in\R$、$\sigma>0$，
\[
 \prod_i[(x_i-\mu)^2+\sigma^2]
 =\prod_i[(y_i-\mu)^2+\sigma^2].
\]
固定 $\mu$ 并令 $\sigma\downarrow0$，得 $P_x(\mu)^2=P_y(\mu)^2$，其中
$P_x(z)=\prod_i(z-x_i)$、$P_y(z)=\prod_i(z-y_i)$。因此多项式
$(P_x-P_y)(P_x+P_y)$ 恒为零。实多项式环无零因子，且两个多项式均首一，所以 $P_x=P_y$。实根及其重数相同，故排序后的样本相同。反之，排序相同显然使似然相同。由第二章的似然成比例判别，$T$ 最小充分。

\section{均匀位置族：最小充分但不完备}
设 $n\ge2$、$X_i\sim U(\theta,\theta+1)$ 独立同分布。证明 $T=(X_{(1)},X_{(n)})$ 最小充分而不完备。

\textbf{解。} 对正似然样本，允许的参数集合为
\[
 (x_{(n)}-1,x_{(1)}),\qquad
 p_\theta(x)=\ind_{\{x_{(n)}-1<\theta<x_{(1)}\}}.
\]
因子分解给出充分性。两条非零似然成比例，等价于上述非空开区间相同，因而两个端点分别相同，得到最小充分性。比较的是完整似然，不是可能出现 $0/0$ 的密度比。

令 $R=X_{(n)}-X_{(1)}$。由 $X_i=\theta+U_i$，$R$ 的分布不含 $\theta$。其密度及期望为
\[
 f_R(r)=n(n-1)r^{n-2}(1-r)\ind_{(0,1)}(r),\qquad
 \E R=\frac{n-1}{n+1}.
\]
于是固定有界函数 $g(T)=R-(n-1)/(n+1)$ 对所有参数均值为零，却不几乎必然为零。因此 $T$ 不有界完备，也不完备。第二章给出了该密度及 Basu 反证的详细推导。

\section{可以恢复充分统计量的统计量仍充分}
在本讲义的受支配模型框架内，设 $T$ 充分且 $T=\psi(S)$，$\psi$ 为不含参数的可测函数，证明 $S$ 充分。

\textbf{解。} 在受支配模型中，由因子分解，$p_\theta(x)=a_\theta(T(x))h(x)$。代入得
\[
 p_\theta(x)=a_\theta(\psi(S(x)))h(x).
\]
右端仍是 $S(x)$ 的函数与参数无关因子的乘积，故 $S$ 充分。方向不可颠倒：充分统计量的任意函数未必充分，例如把它映射成常数。

\section{完备充分统计量与最小充分统计量的关系}
设 $T$ 完备（或有界完备）且充分，最小充分统计量 $S$ 存在。证明 $T$ 最小充分，且 $S$ 也完备（或有界完备）。

\textbf{解。} 最小充分性先给出 $S=\psi(T)$。若 $a(S)$ 可积（在有界版本中要求有界）且所有参数下均值为零，则 $a(\psi(T))$ 具有相同性质；由 $T$ 的完备性，它几乎必然为零，所以 $S$ 完备。

为证明反向可恢复性，任取有界可测 $b$。$S$ 充分，故可统一选取不含参数的 $k$，使
\[
 k(S)=\E_\theta[b(T)\mid S].
\]
若 $|b|\le M$，可取 $|k|\le M$。残差 $b(T)-k(\psi(T))$ 是 $T$ 的固定有界函数，所有参数下均值为零，故由有界完备性
\[
 b(T)=k(S)\quad P_\theta\text{-a.s.}
\]
实值时取 $b(t)=\arctan t$，反解得到 $T=\tan(k(S))$；逆函数在有效像外可补定义。有限维统计量对每个坐标分别如此处理。一般标准 Borel 值域可用有界 Borel 编码。任意充分统计量 $U$ 都能恢复 $S$，从而恢复 $T$，所以 $T$ 最小充分。不能直接对未保证可积的 $T$ 作条件期望，也不能把一般无界残差用于有界完备性。

\section{平移指数分布的间距与独立性}
设 $n>2$，样本密度为
\[
 f_{a,\theta}(x)=\theta^{-1}e^{-(x-a)/\theta}\ind_{\{x>a\}},
 \quad a\in\R,\quad\theta>0.
\]
证明 $X_{(1)}$ 与 $Q=\sum_i(X_i-X_{(1)})$ 独立；再证明向量
\[
 Z_i=\frac{X_{(n)}-X_{(i)}}{X_{(n)}-X_{(n-1)}},\quad 1\le i\le n-2,
\]
与 $(X_{(1)},Q)$ 独立。

\textbf{解。} 定义缩放间距
\[
 W_1=n(X_{(1)}-a),\quad
 W_j=(n-j+1)(X_{(j)}-X_{(j-1)}),\quad 2\le j\le n.
\]
反变换为 $X_{(j)}=a+W_1/n+\sum_{k=2}^jW_k/(n-k+1)$，其 Jacobian 绝对值为 $1/n!$。次序统计量密度为
\[
 n!\theta^{-n}\exp\{-\textstyle\sum_j(x_j-a)/\theta\}
 \ind_{\{a<x_1<\cdots<x_n\}}.
\]
又有 $\sum_j(x_j-a)=\sum_jw_j$，故
\[
 f_W(w)=\prod_{j=1}^n\theta^{-1}e^{-w_j/\theta}\ind_{\{w_j>0\}}.
\]
即 $W_j$ 独立同分布于尺度为 $\theta$ 的指数分布。计数每个间距出现的次数，得
\[
 Q=\sum_{j=2}^nW_j,\qquad X_{(1)}=a+W_1/n,
\]
所以第一问成立。

令 $P_j=W_j/Q$，$j=2,\ldots,n$；只需 $P_2,\ldots,P_{n-1}$ 为自由变量。变换 $w_j=qp_j$ 的 Jacobian 为 $q^{n-2}$，于是
\[
 f_{Q,P}(q,p)=
 \frac{q^{n-2}e^{-q/\theta}}{(n-2)!\theta^{n-1}}\ind_{\{q>0\}}
 \cdot (n-2)!\ind_{\Delta}(p),
\]
其中 $\Delta=\{p_j>0,\sum_{j=2}^{n-1}p_j<1\}$。前者是 Gamma 密度，后者在体积为 $1/(n-2)!$ 的单纯形上积分为 $1$，故 $Q\perp\!\!\!\perp P$。再由 $W_1$ 与 $(W_2,\ldots,W_n)$ 独立，联合密度分解为
$f_{W_1}f_Qf_P$，得到 $P\perp\!\!\!\perp(W_1,Q)$。这是联合独立性，不能仅由两次两两独立直接推断。

最后
\[
 Z_i=\frac{\sum_{k=i+1}^nP_k/(n-k+1)}{P_n},\quad P_n>0\text{ a.s.}
\]
整个 $Z$ 向量是 $P$ 的函数，$(X_{(1)},Q)$ 是 $(W_1,Q)$ 的函数，故第二问成立。该结论不声称 $Z_i$ 之间相互独立。

\section{两段支撑模型的最小充分统计量}
设 $\theta>0$，$X_i$ 独立同分布，密度为
\[
 f_\theta(x)=\frac1{2\theta}
 [\ind_{(0,\theta)}(x)+\ind_{(2\theta,3\theta)}(x)].
\]
求一个最小充分统计量。

\textbf{解。} 只需考虑共同满测集 $(0,\infty)^n$。令
\[
\begin{aligned}
 M&=\max_i x_i,\\
 A&=\max(\{x_i:x_i<M/2\}\cup\{0\}),\\
 B&=\min\{x_i:x_i\ge M/2\},\\
 L&=\max(M/3,A),\qquad U=B/2.
\end{aligned}
\]
$B$ 总有定义，因为最大值属于相应集合。令 $(L^*,U^*)=(L,U)$ 若 $L<U$，否则令其为 $(0,0)$。下面证明 $T=(M,L^*,U^*)$ 最小充分。

由最大值所在区间，正似然参数必满足 $\theta>M$ 或 $M/3<\theta<M/2$。前者使全部观测落入第一段。后者中，$x_i<M/2$ 时只能属于第一段，要求 $x_i<\theta$；$x_i\ge M/2$ 时只能属于第二段，要求 $2\theta<x_i$，而 $x_i<3\theta$ 已由 $x_i\le M<3\theta$ 保证。故正似然参数集合恰为
\[
 C(x)=(M,\infty)\cup(L,U).
\]
空区间用 $(L^*,U^*)=(0,0)$ 统一编码。联合密度为
\[
 p_\theta(x)=(2\theta)^{-n}
 [\ind_{\{\theta>M\}}+\ind_{\{L^*<\theta<U^*\}}],
\]
给出充分性。

若 $p_\theta(x)=c\,p_\theta(y)$ 对所有参数成立，取 $\theta>\max(M_x,M_y)$，有 $c=1$，所以 $C(x)=C(y)$。该集合唯一无界连通分支的左端点恢复 $M$；剩余非空有界分支恢复 $L,U$，剩余为空则恢复编码 $(0,0)$。因此 $T(x)=T(y)$。反之 $T$ 相同使似然相同。由似然成比例判别，$T$ 最小充分。

\begin{notebox}[支撑与空区间]
记录区间端点时必须把所有空区间统一编码。否则不同的 $L\ge U$ 可能对应同一条似然，额外保留无用端点便不能由上述判别得到最小性。
\end{notebox}

\chapter{估计、信息与 U 统计量习题}

本组对应第二次作业的全部 11 题。题意按数学条件重新表述。风险比较均在平方损失下进行；最优性总是相对于明确给定的模型和规则类。第 10 题沿用第 9 题的投影记号。

\section{第 1 题 二项分布中不存在 1/p 的无偏估计}

\textbf{题意。} $X\sim\operatorname{Bin}(n,p)$，$n$ 已知，$p\in(0,1)$ 未知，证明 $p^{-1}$ 不存在无偏估计。

 样本空间只有 $n+1$ 个点，任何统计量的期望都是关于 $p$ 的至多 $n$ 次多项式，不可能在整个 $(0,1)$ 上等于 $1/p$。

\textbf{证明。} 反设存在统计量 $\delta(X)$，使得

\[
E_p\delta(X)=\frac1p,\qquad 0<p<1.
\]

记 $\delta(j)=d_j$。由于每个点 $j=0,\ldots,n$ 在每个 $p\in(0,1)$ 下都有正概率，而无偏估计要求期望存在且有限，故每个 $d_j$ 都是有限实数。于是

\[
P(p):=E_p\delta(X)
=\sum_{j=0}^n d_j\binom nj p^j(1-p)^{n-j}
\]

是至多 $n$ 次多项式。无偏性给出

\[
pP(p)-1=0,\qquad 0<p<1.
\]

左边是一个多项式，却在区间 $(0,1)$ 上有无穷多个零点，因此必为零多项式。代入 $p=0$，得到 $-1=0$，矛盾。因此

\[
\boxed{\text{不存在 }p^{-1}\text{ 的无偏估计。}}
\]

 这里把 $p=0$ 代入的是已经证明恒为零的多项式，不是擅自把参数空间扩大到 $p=0$。也可直接令 $p\downarrow0$：左边期望趋于有限的 $d_0$，右边 $1/p\to\infty$。

\section{第 2 题 平方误差下的最优规则与最优线性无偏估计}

\textbf{题意。} $X_1,\ldots,X_n$ 有共同有限均值 $\mu$ 和有限方差，估计 $\mu$，损失为 $(\delta-\mu)^2$。

\subsection{(a) 所有估计量组成的规则类中不存在统一最优规则}

 任意常数规则都允许时，假如存在一个处处最优的规则，它必须在每个参数点都达到零风险，即从一份样本精确知道未知均值。正态子模型已经能排除这种可能。

\textbf{证明。} 假设 $\delta^*(X)$ 在所有估计量中统一最优。对每个固定实数 $c$，常数规则 $\delta_c(X)\equiv c$ 都属于允许的规则类。因此在均值为 $\mu$ 的任意模型分布下，

\[
R_\mu(\delta^*)\le R_\mu(\delta_\mu)=0.
\]

这里 $\delta_\mu$ 的意思是：在比较真值 $\mu$ 处风险时，从全部常数规则中取数值恰为 $\mu$ 的那一个；并未允许实际估计规则依赖未知参数。

风险非负，故上式强迫

\[
E_\mu(\delta^*-\mu)^2=0,
\qquad
\delta^*=\mu\quad P_\mu\text{-a.s.}
\]

现在只考虑题目允许的子模型

\[
X_1,\ldots,X_n\overset{\mathrm{i.i.d.}}{\sim}N(\mu,1),
\qquad \mu\in\mathbb R.
\]

取两个不同的均值 $\mu_0\ne\mu_1$。这两个样本分布在 $\mathbb R^n$ 上的密度均处处为正，因此具有相同的零测集。由

\[
\delta^*=\mu_0\quad P_{\mu_0}\text{-a.s.}
\]

可推出

\[
\delta^*=\mu_0\quad P_{\mu_1}\text{-a.s.}
\]

但最优性同时要求 $\delta^*=\mu_1$，$P_{\mu_1}$-几乎处处成立，与 $\mu_0\ne\mu_1$ 矛盾。故不存在统一最优规则。

\textbf{适用范围。} 仅凭“共同均值、有限方差”不能对任意另外指定的小模型都断言无最优规则。例如若模型只有 $X_1=\cdots=X_n=\mu$ 几乎处处的退化分布，则 $X_1$ 风险恒为零。上述证明针对本题通常意指的完整分布类；它包含正态子模型，所以结论成立。

\subsection{(b) 独立、方差不同时的最优权重}

设

\[
\operatorname{Var}(X_i)=\frac{\sigma^2}{a_i},
\qquad a_i>0,\quad \sigma^2>0,
\]

其中 $a_i$ 已知。考虑

\[
\mathcal I_2=
\left\{\sum_{i=1}^n c_iX_i:\ c_i\in\mathbb R,
\sum_{i=1}^n c_i=1\right\}.
\]

\textbf{解。} 对 $\delta_c=\sum_i c_iX_i$，

\[
E\delta_c=\sum_i c_i\mu=\mu,
\]

故类内所有规则均无偏。利用独立性，

\[
R(\delta_c)=\operatorname{Var}(\delta_c)
=\sigma^2\sum_{i=1}^n\frac{c_i^2}{a_i}.
\]

由 Cauchy--Schwarz 不等式，

\[
1=\left(\sum_i c_i\right)^2
=\left(\sum_i\frac{c_i}{\sqrt{a_i}}\sqrt{a_i}\right)^2
\le
\left(\sum_i\frac{c_i^2}{a_i}\right)\left(\sum_i a_i\right).
\]

因此

\[
R(\delta_c)\ge\frac{\sigma^2}{\sum_i a_i}.
\]

等号成立当且仅当存在常数 $\lambda$，使

\[
\frac{c_i}{\sqrt{a_i}}=\lambda\sqrt{a_i},
\qquad c_i=\lambda a_i.
\]

结合 $\sum_i c_i=1$，得 $\lambda=(\sum_i a_i)^{-1}$。所以

\[
\boxed{\delta^*=\frac{\sum_{i=1}^n a_iX_i}{\sum_{i=1}^n a_i}},
\qquad
\boxed{R(\delta^*)=\frac{\sigma^2}{\sum_{i=1}^n a_i}}.
\]

权重与方差的倒数成正比；公共未知因子 $\sigma^2$ 在权重归一化时消去，因此规则不依赖未知参数。

\subsection{(c) 同分布且任意两项具有相同相关系数}

设共同方差为 $\sigma^2>0$，且对所有 $i\ne j$，

\[
\operatorname{Cov}(X_i,X_j)=\rho\sigma^2.
\]

对 $\sum_i c_i=1$ 的规则仍有无偏性，并且

\[
\begin{aligned}
R(\delta_c)
&=\sigma^2\sum_i c_i^2
+2\rho\sigma^2\sum_{i<j}c_ic_j\\
&=\sigma^2\left[
(1-\rho)\sum_i c_i^2+\rho\left(\sum_i c_i\right)^2
\right]\\
&=\sigma^2\left[(1-\rho)\sum_i c_i^2+\rho\right].
\end{aligned}
\]

相关矩阵必须半正定，故当 $n\ge2$ 时，

\[
-\frac1{n-1}\le\rho\le1.
\]

若 $\rho<1$，最小化风险等价于最小化 $\sum_i c_i^2$。再次用 Cauchy--Schwarz，

\[
1=\left(\sum_i c_i\right)^2\le n\sum_i c_i^2,
\]

等号当且仅当 $c_1=\cdots=c_n=1/n$。因此

\[
\boxed{\delta^*=\bar X},
\qquad
\boxed{R(\bar X)=\frac{\sigma^2}{n}[1+(n-1)\rho]}.
\]

若 $\rho=1$，上式风险对任何满足 $\sum_i c_i=1$ 的权重都等于 $\sigma^2$，所以所有类内规则均最优，$\bar X$ 仍是其中之一。若 $\rho=-1/(n-1)$，$\bar X$ 的风险为零。

\section{第 3 题 有限参数空间中的一致性与任意收敛速度}

\textbf{题意。} 参数空间 $\Theta=\{\theta_1,\ldots,\theta_k\}$ 有限，估计量 $T_n$ 也只取这些值。证明一致性等价于选对参数的概率趋于 1，并证明一致估计量具有任意正数列所指定的 $O_P(1)$ 收敛速度。

\subsection{(a) 一致性当且仅当选对概率趋于 1}

若参数空间只有一个点，结论显然。以下设至少有两个不同参数，并记

\[
d=\min_{i\ne j}|\theta_i-\theta_j|>0.
\]

固定真值 $\theta\in\Theta$。因为 $T_n\in\Theta$，只要 $T_n\ne\theta$，就有 $|T_n-\theta|\ge d$，所以

\[
\{T_n\ne\theta\}=\{|T_n-\theta|>d/2\}.
\]

若 $T_n\xrightarrow{P_\theta}\theta$，则

\[
P_\theta(T_n\ne\theta)
=P_\theta(|T_n-\theta|>d/2)\to0.
\]

反过来，若 $P_\theta(T_n=\theta)\to1$，则对任意 $\varepsilon>0$，

\[
P_\theta(|T_n-\theta|>\varepsilon)
\le P_\theta(T_n\ne\theta)\to0.
\]

因而

\[
\boxed{T_n\xrightarrow{P_\theta}\theta
\iff P_\theta(T_n=\theta)\to1.}
\]

\subsection{(b) 任意正数列放大后的误差仍依概率趋于零}

取任意确定性正数列 $a_n$，无需限制增长速度。对每个 $\varepsilon>0$，

\[
\{|a_n(T_n-\theta)|>\varepsilon\}\subseteq\{T_n\ne\theta\}.
\]

因此由 (a)，

\[
P_\theta(|a_n(T_n-\theta)|>\varepsilon)
\le P_\theta(T_n\ne\theta)\to0.
\]

所以

\[
\boxed{a_n(T_n-\theta)\xrightarrow{P_\theta}0}.
\]

特别地，$a_n(T_n-\theta)=O_{P_\theta}(1)$，即通常定义下的 $a_n$-一致性。

 错误时，放大后的误差可能极大；但只要估计恰好正确，放大多少倍都仍是零。错误事件的概率趋于零已经足够控制依概率收敛。这里并没有推出放大误差的方差或期望也趋于零。

\section{第 4 题 样本均值与样本中程的风险比较}

\textbf{题意。} $n>2$，$X_i$ 独立同分布，二阶矩有限。记

\[
\bar X=\frac1n\sum_{i=1}^nX_i,\qquad
T=\frac{X_{(1)}+X_{(n)}}2.
\]

在平方误差下比较二者估计位置参数 $\theta$ 的风险。

\subsection{(a) 正态模型中样本均值严格优于中程}

设 $X_i\overset{\mathrm{i.i.d.}}{\sim}N(\theta,\sigma^2)$，$\sigma>0$。

\textbf{第一步：写出均值与残差分解。} 令

\[
R_i=X_i-\bar X,\qquad
D=\frac{\min_iR_i+\max_iR_i}{2}.
\]

将每个观测都减去同一数 $\bar X$ 不改变排序，故

\[
T=\bar X+D.
\]

由于 $|D|\le\max_i|R_i|$，有 $D^2\le\sum_iR_i^2$，故 $E D^2<\infty$。

\textbf{第二步：证明 D 与样本均值独立且均值为零。} 对每个 $i$，

\[
\operatorname{Cov}(\bar X,R_i)
=\operatorname{Cov}(\bar X,X_i)-\operatorname{Var}(\bar X)
=\frac{\sigma^2}{n}-\frac{\sigma^2}{n}=0.
\]

$(\bar X,R_1,\ldots,R_n)$ 联合正态，因此 $\bar X$ 与整个残差向量独立，从而与 $D$ 独立。

残差向量关于零对称，即 $(R_i)_i$ 与 $(-R_i)_i$ 同分布，而

\[
D(-R_1,\ldots,-R_n)=-D(R_1,\ldots,R_n).
\]

因此 $ED=0$，于是 $ET=\theta$。利用独立性，

\[
R_\theta(T)
=\operatorname{Var}(T)
=\frac{\sigma^2}{n}+\operatorname{Var}(D)
\ge\frac{\sigma^2}{n}
=R_\theta(\bar X).
\]

\textbf{第三步：证明严格不等号。} 函数

\[
d(x)=\frac{\min_i x_i+\max_i x_i}{2}
-\frac1n\sum_i x_i
\]

连续，并且在 $x=(0,\ldots,0,1)$ 处，

\[
d(x)=\frac12-\frac1n>0
\]

（这里使用了 $n>2$）。所以在该点的某个开邻域内，$d(x)>0$。正态样本的联合密度处处为正，该邻域具有正概率，故 $P(D\ne0)>0$。结合 $ED=0$，得到 $\operatorname{Var}(D)>0$。

因此对所有 $\theta\in\mathbb R,\sigma>0$，

\[
\boxed{R_\theta(\bar X)<R_\theta(T)}.
\]

上述分解也说明 $E(T\mid\bar X)=\bar X$，因此可以把结论理解为 Rao--Blackwell 改进；严格性仍需第三步，不能仅凭“条件期望不增方差”省略。

\subsection{(b) 单位长度均匀模型中中程严格优于样本均值}

设 $X_i\sim U(\theta-\tfrac12,\theta+\tfrac12)$。令

\[
Y_i=X_i-\theta+\frac12\sim U(0,1),\qquad
A=Y_{(1)},\quad B=Y_{(n)}.
\]

于是

\[
T-\theta=\frac{A+B-1}{2}.
\]

先计算所需次序统计量矩。最小值与最大值的密度分别为

\[
f_A(a)=n(1-a)^{n-1},\quad
f_B(b)=nb^{n-1},\qquad 0<a,b<1.
\]

利用积分 $\int_0^1u^r(1-u)^sdu=r!s!/(r+s+1)!$，得到

\[
EA=\frac1{n+1},\qquad EA^2=\frac2{(n+1)(n+2)},
\]

\[
EB=\frac n{n+1},\qquad EB^2=\frac n{n+2}.
\]

从而

\[
\operatorname{Var}(A)=\operatorname{Var}(B)
=\frac{n}{(n+1)^2(n+2)}.
\]

为计算协方差，使用二者联合密度

\[
f_{A,B}(a,b)=n(n-1)(b-a)^{n-2},
\qquad 0<a<b<1.
\]

该密度来自：选一个观测落在最小值附近、一个落在最大值附近，其余 $n-2$ 个都落在长度为 $b-a$ 的区间内。于是

\[
\begin{aligned}
E(AB)
&=n(n-1)\int_0^1\int_0^b ab(b-a)^{n-2}\,da\,db\\
&=n(n-1)\int_0^1 b^{n+1}\,db
\int_0^1u(1-u)^{n-2}\,du\\
&=n(n-1)\frac1{n+2}\frac1{n(n-1)}
=\frac1{n+2}.
\end{aligned}
\]

第二行用了 $a=bu$。因此

\[
\operatorname{Cov}(A,B)
=\frac1{n+2}-\frac n{(n+1)^2}
=\frac1{(n+1)^2(n+2)}.
\]

因为 $EA+EB=1$，$T$ 无偏，而且

\[
\begin{aligned}
R_\theta(T)
&=\frac14\left[\operatorname{Var}(A)+\operatorname{Var}(B)
+2\operatorname{Cov}(A,B)\right]\\
&=\boxed{\frac1{2(n+1)(n+2)}}.
\end{aligned}
\]

另一方面，$\bar X$ 无偏，且

\[
R_\theta(\bar X)=\frac{\operatorname{Var}(X_1)}n
=\frac1{12n}.
\]

二者之差为

\[
R_\theta(\bar X)-R_\theta(T)
=\frac{(n+1)(n+2)-6n}{12n(n+1)(n+2)}
=\frac{(n-1)(n-2)}{12n(n+1)(n+2)}>0.
\]

所以

\[
\boxed{R_\theta(T)<R_\theta(\bar X),\qquad n>2.}
\]

\subsection{(c) 构造互不支配的模型族}

取

\[
\boxed{\mathcal P=
\{N(\theta,1):\theta\in\mathbb R\}
\ \cup\
\{U(\theta-\tfrac12,\theta+\tfrac12):\theta\in\mathbb R\}.}
\]

估计目标统一定义为分布的均值 $\theta$。其中分布类型也未知，可以看成一个额外的干扰参数。

在正态子族内，(a) 给出 $\bar X$ 的风险严格更小；在均匀子族内，(b) 给出 $T$ 的风险严格更小。因此在整个 $\mathcal P$ 上，两种规则都不能处处支配另一种，即

\[
\boxed{\bar X\text{ 与 }T\text{ 互不支配。}}
\]

这并不意味着二者在所有规则中都可容许；这里只证明它们彼此不能支配对方。

\section{第 5 题 统计量不能增加 Fisher 信息}

\textbf{题意。} 样本 $X$ 的密度为 $f_\theta$，统计量 $T=T(X)$ 的密度为 $g_\theta$，$\theta\in\mathbb R^k$。两种密度关于参数可微，且在参数有界区域内，其一阶导数分别有相对于支配测度可积的包络。证明：

\begin{enumerate}
\item $I_X(\theta)-I_T(\theta)$ 半正定；
\item 若 $T$ 充分，则两者相等。

\end{enumerate}
以下在参数空间内点讨论，并按通常定义假设涉及的 Fisher 信息矩阵各项有限，以使矩阵差有意义。若某个方向的信息为无穷，应把结论理解为该方向上的扩展非负不等式，而不写“无穷减无穷”。

 先证明“统计量的得分＝样本得分在统计量条件下的平均”，再把信息之差写成残差的二阶矩。

\subsection{准备 得分恒等式}

将得分都写成列向量：

\[
S_X(X)=\nabla_\theta\log f_\theta(X),\qquad
S_T(T)=\nabla_\theta\log g_\theta(T).
\]

密度为零处可把得分定义为零。在参数内点，非负可微函数 $\theta\mapsto f_\theta(x)$ 若在当前点为零，其梯度也为零；对 $g_\theta$ 同理。因此

\[
S_X(x)f_\theta(x)=\nabla_\theta f_\theta(x),\qquad
S_T(t)g_\theta(t)=\nabla_\theta g_\theta(t)
\]

在积分意义下成立。题目的可积包络允许交换积分与求导，故

\[
E_\theta S_X
=\int\nabla_\theta f_\theta\,d\nu
=\nabla_\theta1=0,
\qquad E_\theta S_T=0.
\]

取任意有界可测标量函数 $b$。由 $T$ 的分布是 $X$ 经统计量映射得到的分布，

\[
\int b(T(x))f_\theta(x)\,d\nu(x)
=\int b(t)g_\theta(t)\,d\lambda(t).
\]

两侧关于 $\theta$ 求导，有

\[
\int b(T(x))\nabla_\theta f_\theta(x)\,d\nu(x)
=\int b(t)\nabla_\theta g_\theta(t)\,d\lambda(t).
\]

于是

\[
E_\theta[b(T)S_X]=E_\theta[b(T)S_T(T)].
\]

$S_T(T)$ 是 $T$ 的可测函数，上式又对所有有界可测 $b$ 成立，因此由条件期望的刻画，

\[
\boxed{E_\theta(S_X\mid T)=S_T(T).}
\tag{5.1}
\]

\subsection{(a) 信息差为非负定矩阵}

令 $D=S_X-S_T(T)$。由 (5.1)，

\[
E_\theta(D\mid T)=0.
\]

因此

\[
E_\theta[D S_T(T)^\top]
=E_\theta[E_\theta(D\mid T)S_T(T)^\top]=0,
\]

转置方向的交叉项也为零。展开得

\[
\begin{aligned}
I_X(\theta)
&=E_\theta(S_XS_X^\top)\\
&=E_\theta[(S_T+D)(S_T+D)^\top]\\
&=I_T(\theta)+E_\theta(DD^\top).
\end{aligned}
\]

所以

\[
\boxed{I_X(\theta)-I_T(\theta)
=E_\theta[(S_X-S_T)(S_X-S_T)^\top]\succeq0.}
\]

具体而言，对任意 $a\in\mathbb R^k$，

\[
a^\top[I_X-I_T]a
=E_\theta[(a^\top D)^2]\ge0,
\]

这正是半正定的定义。

\subsection{(b) 充分统计量保留全部 Fisher 信息}

\textbf{需要证明的关键点是 $S_X=S_T(T)$ 几乎处处。} 只知道 (5.1) 还不够，因为一般条件平均会消去一部分随机波动。

由充分性，对每个有界可测函数 $a(x)$，可以选择一个不依赖参数的函数 $K_a(t)$，使

\[
E_\theta[a(X)\mid T]=K_a(T)
\]

对所有 $\theta$ 成立，且 $|K_a|\le\sup|a|$。于是

\[
E_\theta a(X)=E_\theta K_a(T).
\]

对参数求导，题目的包络条件给出

\[
E_\theta[a(X)S_X(X)]
=E_\theta[K_a(T)S_T(T)].
\]

再利用条件期望，

\[
E_\theta[K_a(T)S_T(T)]
=E_\theta[E_\theta(a(X)\mid T)S_T(T)]
=E_\theta[a(X)S_T(T)].
\]

因此对所有有界可测 $a$，

\[
E_\theta[a(X)(S_X-S_T(T))]=0.
\]

对差向量的第 $j$ 个分量 $D_j$，分别取 $a(X)=\mathbf1_{\{D_j>0\}}$ 和 $a(X)=\mathbf1_{\{D_j<0\}}$，得到其正部、负部的期望均为零，故 $D_j=0$ 几乎处处。逐分量可得

\[
S_X=S_T(T)\quad P_\theta\text{-a.s.}
\]

代入 (a) 的信息差公式，

\[
\boxed{I_X(\theta)=I_T(\theta).}
\]

\textbf{证明中的分工。} 可积包络负责交换求导与积分；充分性负责使 $K_a$ 对所有参数是同一个函数。若 $K_a$ 也随参数变化，求导时就会多出一项，上述证明不能照搬。

\section{第 6 题 用似然比差分推导无偏估计的方差下界}

\textbf{题意。} $X_1,\ldots,X_n$ 独立同分布，密度 $f_\theta(x)>0$，$\theta\in\Theta\subset\mathbb R$。证明原题给出的差分信息下界，并在整数均值的正态模型中计算其最大值。

 Cramér--Rao 不等式用“参数的微小导数”与估计误差建立协方差。本题改用两个不同参数之间的似然比，不要求参数空间包含开区间。

\subsection{(a) 一般下界}

\textbf{证明。} 固定 $\theta$，取任意 $\Delta\ne0$ 且 $\theta+\Delta\in\Theta$。定义样本似然比

\[
L_\Delta
=\prod_{i=1}^n
\frac{f_{\theta+\Delta}(X_i)}{f_\theta(X_i)}.
\]

密度处处为正，故比值有定义。它是从 $P_\theta^{\otimes n}$ 到 $P_{\theta+\Delta}^{\otimes n}$ 的 Radon--Nikodym 导数，因此

\[
E_\theta L_\Delta=1,\qquad
E_\theta[T L_\Delta]=E_{\theta+\Delta}T.
\]

由 $T$ 对所有参数无偏，

\[
E_\theta T=\theta,\qquad
E_{\theta+\Delta}T=\theta+\Delta.
\]

所以

\[
\begin{aligned}
E_\theta[(T-\theta)(L_\Delta-1)]
&=E_\theta(TL_\Delta)-E_\theta T
-\theta E_\theta L_\Delta+\theta\\
&=(\theta+\Delta)-\theta-\theta+\theta\\
&=\Delta.
\end{aligned}
\]

若相关二阶矩有限，由 Cauchy--Schwarz，

\[
\Delta^2
\le E_\theta(T-\theta)^2\,E_\theta(L_\Delta-1)^2
=\operatorname{Var}_\theta(T)[E_\theta L_\Delta^2-1].
\]

独立性给出

\[
\begin{aligned}
E_\theta L_\Delta^2
&=\prod_{i=1}^n
E_\theta\left[
\frac{f_{\theta+\Delta}(X_i)^2}{f_\theta(X_i)^2}
\right]\\
&=\left[
\int_{\mathbb R}
\frac{f_{\theta+\Delta}(x)^2}{f_\theta(x)}\,dx
\right]^n.
\end{aligned}
\]

于是

\[
\operatorname{Var}_\theta(T)
\ge
\frac{\Delta^2}{
\left[\displaystyle\int_{\mathbb R}
\frac{f_{\theta+\Delta}(x)^2}{f_\theta(x)}\,dx\right]^n-1}.
\]

该不等式对每个允许的 $\Delta$ 都成立，取上确界得

\[
\boxed{
\operatorname{Var}_\theta(T)\ge
\sup_{\substack{\Delta\ne0\\\theta+\Delta\in\Theta}}
\frac{\Delta^2}{
\left[\displaystyle\int_{\mathbb R}
\frac{f_{\theta+\Delta}(x)^2}{f_\theta(x)}\,dx\right]^n-1}.}
\]

若 $\operatorname{Var}_\theta(T)=\infty$，结论显然；若某个 $\Delta$ 的似然比二阶矩为无穷，则该项下界按 0 理解。若分母为零，则 $L_\Delta=1$ 几乎处处，两个样本分布相同，不可能同时有期望 $\theta$ 和 $\theta+\Delta$；所以在确有无偏估计的前提下，不会出现这一冲突。

\subsection{(b) 整数均值的正态模型}

设 $f_\theta$ 是 $N(\theta,\sigma^2)$ 的密度，$\sigma^2>0$ 已知，$\Theta=\mathbb Z$。一份观测的密度比为

\[
\frac{f_{\theta+\Delta}(x)}{f_\theta(x)}
=\exp\left\{\frac{\Delta(x-\theta)}{\sigma^2}
-\frac{\Delta^2}{2\sigma^2}\right\}.
\]

利用 $X-\theta\sim N(0,\sigma^2)$ 的矩母函数，

\[
\begin{aligned}
\int\frac{f_{\theta+\Delta}(x)^2}{f_\theta(x)}\,dx
&=E_\theta\exp\left\{
\frac{2\Delta(X-\theta)}{\sigma^2}
-\frac{\Delta^2}{\sigma^2}\right\}\\
&=\exp\left\{
\frac12\sigma^2\left(\frac{2\Delta}{\sigma^2}\right)^2
-\frac{\Delta^2}{\sigma^2}\right\}\\
&=e^{\Delta^2/\sigma^2}.
\end{aligned}
\]

因此

\[
\operatorname{Var}_\theta(T)
\ge\sup_{\Delta\in\mathbb Z\setminus\{0\}}
\frac{\Delta^2}{e^{n\Delta^2/\sigma^2}-1}.
\]

下面确定上确界。令 $c=n/\sigma^2>0$，考虑

\[
\psi(u)=\frac{u}{e^{cu}-1},\qquad u>0.
\]

其导数为

\[
\psi'(u)=
\frac{e^{cu}(1-cu)-1}{(e^{cu}-1)^2}.
\]

令 $v=cu>0$，$r(v)=e^v(1-v)-1$，则 $r(0)=0$ 且

\[
r'(v)=-ve^v<0,
\]

故 $r(v)<0$，进而 $\psi'(u)<0$。所以目标函数随 $\Delta^2$ 严格减小。在非零整数中，最小的 $\Delta^2$ 是 1，恰由 $\Delta=\pm1$ 取得。因此

\[
\boxed{\operatorname{Var}_\theta(T)
\ge\frac1{e^{n/\sigma^2}-1},}
\]

且上确界在 $\Delta=\pm1$ 处取得。

这里“上确界取得”是指选择差分 $\Delta$ 的最优化，不能据此断言存在无偏估计量达到这个方差下界。

\section{第 7 题 一个有限离散模型中的 UMVUE 存在性}

\textbf{题意。} 单个观测 $X$ 的分布为

\[
P_p(X=-1)=2p(1-p),\qquad
P_p(X=k)=p^k(1-p)^{3-k},\ k=0,1,2,3,
\]

其中 $0<p<1$。分别判断 $p$ 与 $p(1-p)$ 是否存在 UMVUE。

记 $q=1-p$。先核对概率和：

\[
2pq+q^3+pq^2+p^2q+p^3
=2pq+(p^3+q^3)+pq(p+q)=1.
\]

 在有限样本空间中，列出所有“期望恒为零的统计量”，再使用 UMVUE 的正交判别。

\subsection{准备一 UMVUE 的正交判别及证明}

设 $\delta$ 是目标函数的无偏估计。则

\[
\boxed{\delta\text{ 是 UMVUE}
\iff E_p[\delta Z]=0
\text{ 对所有零均值无偏统计量 }Z\text{ 及所有 }p.}
\]

这里“零均值无偏统计量”是指同一个 $Z(X)$ 满足 $E_pZ=0$ 对全部 $p$ 成立。

\textbf{必要性。} $\delta+cZ$ 对任意实数 $c$ 仍然无偏。如果 $\delta$ 最优，则固定 $p$ 后，二次函数

\[
\operatorname{Var}_p(\delta+cZ)
=\operatorname{Var}_p(\delta)
+2c\operatorname{Cov}_p(\delta,Z)
+c^2\operatorname{Var}_p(Z)
\]

必须在 $c=0$ 取得最小值。其一次项必须为零，即 $\operatorname{Cov}_p(\delta,Z)=0$。因 $E_pZ=0$，这等价于 $E_p(\delta Z)=0$。

\textbf{充分性。} 任意其他无偏估计 $\tilde\delta$ 都可写为 $\delta+Z$，其中 $E_pZ=0$。由正交性，

\[
\operatorname{Var}_p(\tilde\delta)
=\operatorname{Var}_p(\delta)+\operatorname{Var}_p(Z)
\ge\operatorname{Var}_p(\delta).
\]

故 $\delta$ 是 UMVUE。

\subsection{准备二 求所有均值恒为零的统计量}

设 $Z$ 在 $-1,0,1,2,3$ 处的值分别为 $z_{-1},z_0,z_1,z_2,z_3$。则

\[
E_pZ=2z_{-1}p(1-p)+z_0(1-p)^3
+z_1p(1-p)^2+z_2p^2(1-p)+z_3p^3.
\]

它在整个 $(0,1)$ 上为零，当且仅当各次幂系数为零：

\[
\begin{cases}
z_0=0,\\
2z_{-1}-3z_0+z_1=0,\\
-2z_{-1}+3z_0-2z_1+z_2=0,\\
-z_0+z_1-z_2+z_3=0.
\end{cases}
\]

依次得到

\[
z_0=0,\qquad z_1=-2z_{-1},\qquad
z_2=-2z_{-1},\qquad z_3=0.
\]

所以所有零均值统计量均为下列 $Z_0$ 的常数倍：

\begin{center}
\begin{tabular}{cccccc}
\toprule
$x$ & $-1$ & $0$ & $1$ & $2$ & $3$ \\
\midrule
$Z_0(x)$ & $1$ & $0$ & $-2$ & $-2$ & $0$ \\
\bottomrule
\end{tabular}
\end{center}

\subsection{准备三 把正交条件化为三个函数值相等}

记 $\delta_x=\delta(x)$。则

\[
\begin{aligned}
E_p(\delta Z_0)
&=2pq\delta_{-1}-2pq^2\delta_1-2p^2q\delta_2\\
&=2pq[\delta_{-1}-(1-p)\delta_1-p\delta_2].
\end{aligned}
\]

由于 $pq>0$，该式对所有 $p$ 为零，当且仅当

\[
\boxed{\delta_{-1}=\delta_1=\delta_2.}
\tag{7.1}
\]

因此，一个无偏估计是 UMVUE，当且仅当它满足 (7.1)。

\subsection{(a) p 不存在 UMVUE}

若 $\delta$ 对 $p$ 无偏，即 $E_p\delta=p$，比较上述期望多项式与 $p$ 的系数，可得全部无偏估计为

\begin{center}
\begin{tabular}{cccccc}
\toprule
$x$ & $-1$ & $0$ & $1$ & $2$ & $3$ \\
\midrule
$\delta_a(x)$ & $a$ & $0$ & $1-2a$ & $2-2a$ & $1$ \\
\bottomrule
\end{tabular}
\end{center}

其中 $a\in\mathbb R$ 是不依赖未知参数的常数。

验证过程为：常数项给出 $\delta_0=0$；一次项给出 $\delta_1=1-2a$；二次项给出 $\delta_2=2-2a$；三次项给出 $\delta_3=1$。

但对所有 $a$，

\[
\delta_a(2)-\delta_a(1)=1,
\]

不可能满足 (7.1)。因此

\[
\boxed{p\text{ 有无偏估计，但不存在 UMVUE。}}
\]

例如取 $a=0$ 就给出一个无偏估计。不存在 UMVUE 不是因为无偏估计类为空，而是不存在一个无偏规则能在所有 $p$ 下同时达到该类的最小方差。

还可直接看出矛盾来源：固定 $p$ 时，$\delta_a=\delta_0+aZ_0$ 的方差最小点为

\[
a^*(p)=-\frac{E_p(\delta_0 Z_0)}{E_pZ_0^2}
=\frac{2pq(1+p)}{6pq}
=\frac{1+p}{3}.
\]

它随未知 $p$ 变化，不能用一个固定的 $a$ 同时实现。

\subsection{(b) p(1-p) 存在 UMVUE}

按照 (7.1)，寻找在 $-1,1,2$ 处同取常数 $c$ 的无偏估计。令在 $0,3$ 处均为零，则其期望为

\[
c(2pq+pq^2+p^2q)=3cpq.
\]

取 $c=1/3$ 即得

\[
\boxed{\delta^*(X)=\frac13\mathbf1_{\{X\in\{-1,1,2\}\}}.}
\]

它满足

\[
E_p\delta^*=\frac13(3pq)=pq,
\]

且在 $-1,1,2$ 处函数值相等，故由正交判别，确为 UMVUE。

其方差为

\[
\boxed{\operatorname{Var}_p(\delta^*)
=\frac19(3pq)(1-3pq)
=\frac{pq}{3}-p^2q^2.}
\]

 不完备并不意味着“所有目标都没有 UMVUE”。这里同一个不完备模型中，$p$ 没有 UMVUE，而 $p(1-p)$ 有。

\section{第 8 题 正态均值指数函数的 UMVUE 与 CR 下界}

\textbf{题意。} $X_i\overset{\mathrm{i.i.d.}}{\sim}N(\mu,\sigma^2)$，$\mu\in\mathbb R$ 未知，$\sigma^2>0$ 已知，$t\ne0$ 固定。求 $e^{t\mu}$ 的 UMVUE，比较其方差与 CR 下界，并求二者比值的极限。

\subsection{(a) 构造并证明 UMVUE}

由于

\[
\bar X\sim N\left(\mu,\frac{\sigma^2}{n}\right),
\]

正态矩母函数给出

\[
E_\mu e^{t\bar X}
=\exp\left(t\mu+\frac{t^2\sigma^2}{2n}\right).
\]

因此

\[
\widehat g
=\exp\left(t\bar X-\frac{t^2\sigma^2}{2n}\right)
\]

满足 $E_\mu\widehat g=e^{t\mu}$，是无偏估计。

再证明其最优性。联合密度可写为

\[
f_\mu(x)
=(2\pi\sigma^2)^{-n/2}
\exp\left(-\frac{\sum_i x_i^2}{2\sigma^2}\right)
\exp\left\{\frac{\mu}{\sigma^2}\sum_i x_i
-\frac{n\mu^2}{2\sigma^2}\right\}.
\]

由因子分解定理，$\sum_iX_i$ 充分。它构成自然参数 $\eta=\mu/\sigma^2\in\mathbb R$ 的一参数指数族，自然参数空间含非空开区间，因此由指数族完备性定理，$\sum_iX_i$ 完备。$\bar X$ 是它的一一变换，亦完备且充分。

$\widehat g$ 是完备充分统计量 $\bar X$ 的无偏函数，且正态指数矩有限，故由 Lehmann--Scheffé 定理，

\[
\boxed{\operatorname{UMVUE}(e^{t\mu})
=\exp\left(t\bar X-\frac{t^2\sigma^2}{2n}\right).}
\]

\subsection{(b) 精确方差严格大于 CR 下界}

记 $a=t^2\sigma^2/n>0$。先计算二阶矩：

\[
\begin{aligned}
E_\mu\widehat g^{\,2}
&=e^{-t^2\sigma^2/n}E_\mu e^{2t\bar X}\\
&=e^{-t^2\sigma^2/n}
\exp\left(2t\mu+\frac{(2t)^2\sigma^2}{2n}\right)\\
&=\exp(2t\mu+a).
\end{aligned}
\]

所以

\[
\boxed{\operatorname{Var}_\mu(\widehat g)
=e^{2t\mu}(e^a-1).}
\]

另一方面，一份观测的得分与信息为

\[
S_\mu(X)=\frac{X-\mu}{\sigma^2},
\qquad I_1(\mu)=\frac1{\sigma^2}.
\]

样本信息为 $I_n(\mu)=n/\sigma^2$。对 $g(\mu)=e^{t\mu}$，有 $g'(\mu)=te^{t\mu}$，故无偏估计的 CR 下界为

\[
\mathrm{CRLB}
=\frac{[g'(\mu)]^2}{I_n(\mu)}
=\frac{t^2\sigma^2}{n}e^{2t\mu}
=ae^{2t\mu}.
\]

当 $a>0$ 时，由

\[
e^a-1=\int_0^a e^u\,du>\int_0^a1\,du=a,
\]

可得

\[
\boxed{\operatorname{Var}_\mu(\widehat g)>\mathrm{CRLB}.}
\]

\subsection{(c) 方差与下界之比趋于 1}

由前面的公式，

\[
\frac{\operatorname{Var}_\mu(\widehat g)}{\mathrm{CRLB}}
=\frac{e^a-1}{a}.
\]

固定 $t,\sigma^2$，当 $n\to\infty$ 时 $a\to0$，且

\[
\frac{e^a-1}{a}
=\int_0^1 e^{av}\,dv\to1.
\]

因此

\[
\boxed{
\frac{\operatorname{Var}_\mu(\widehat g)}{\mathrm{CRLB}}\longrightarrow1.}
\]

 UMVUE 的意思是在所有无偏估计中方差最小，不意味着必须达到 CR 下界。本题中有限样本下界不可达，但 UMVUE 的方差与下界在一阶渐近意义下相同。

\section{第 9 题 U 统计量的条件投影及协方差}

\textbf{题意与记号。} $X_1,\ldots,X_n$ 独立同分布，$h$ 为对称的 $m$ 元核，$m\le n$。假定
\[
 E h(X_1,\ldots,X_m)^2<\infty.
\]
定义
\[
\begin{gathered}
 U_n=\binom nm^{-1}\sum_{1\le i_1<\cdots<i_m\le n}h(X_{i_1},\ldots,X_{i_m}),\\
 \vartheta=E h(X_1,\ldots,X_m).
\end{gathered}
\]

\[
h_k(x_1,\ldots,x_k)
=E[h(x_1,\ldots,x_k,X_{k+1},\ldots,X_m)],
\qquad k=1,\ldots,m,
\]

其中 $h_m=h$。记

\[
\widetilde h_k=h_k-\vartheta,\qquad
\widetilde h=h-\vartheta,\qquad
\zeta_k=\operatorname{Var}(h_k(X_1,\ldots,X_k)).
\]

 $h_k$ 保留前 $k$ 个观测，把剩下的 $m-k$ 个观测平均掉。作为随机变量，

\[
h_k(X_1,\ldots,X_k)
=E[h(X_1,\ldots,X_m)\mid X_1,\ldots,X_k].
\tag{9.1}
\]

二阶可积性保证 $h$ 一阶可积，故可以使用 Fubini 定理与条件期望；由条件 Jensen 不等式，

\[
E[h_k(X_1,\ldots,X_k)^2]\le E[h(X_1,\ldots,X_m)^2]<\infty.
\]

对任意具体 $x_1,\ldots,x_k$ 的积分公式在积分存在处使用；条件期望恒等式均按几乎处处理解。仅有题目的二阶矩条件，不需要也不能额外假定每个样本空间点上的截面积分都有限。

\subsection{(a) 逐次积分恒等式}

\textbf{证明。} 对 $1\le k\le m-1$，利用独立同分布性和 Fubini 定理，

\[
\begin{aligned}
&E[h_{k+1}(x_1,\ldots,x_k,X_{k+1})]\\
&=\int h_{k+1}(x_1,\ldots,x_k,u_{k+1})\,dP(u_{k+1})\\
&=\int\cdots\int
h(x_1,\ldots,x_k,u_{k+1},\ldots,u_m)
\prod_{j=k+1}^m dP(u_j)\\
&=h_k(x_1,\ldots,x_k).
\end{aligned}
\]

所以

\[
\boxed{h_k(x_1,\ldots,x_k)
=E[h_{k+1}(x_1,\ldots,x_k,X_{k+1})].}
\]

随机变量形式就是塔式法则：

\[
E[h_{k+1}(X_1,\ldots,X_{k+1})\mid X_1,\ldots,X_k]
=h_k(X_1,\ldots,X_k).
\]

\subsection{(b) 中心化 U 统计量等于中心化核的平均}

每一个 $m$ 元子样本都与 $(X_1,\ldots,X_m)$ 同分布，所以每项核的期望均为 $\vartheta$。求和中共有 $\binom nm$ 项，因此

\[
EU_n=\binom nm^{-1}\binom nm\,\vartheta=\vartheta.
\]

于是

\[
\begin{aligned}
U_n-EU_n
&=\binom nm^{-1}\sum_{i_1<\cdots<i_m}h(X_{i_1},\ldots,X_{i_m})
-\vartheta\\
&=\binom nm^{-1}\sum_{i_1<\cdots<i_m}
[h(X_{i_1},\ldots,X_{i_m})-\vartheta]\\
&=\boxed{\binom nm^{-1}\sum_{i_1<\cdots<i_m}
\widetilde h(X_{i_1},\ldots,X_{i_m})}.
\end{aligned}
\]

\subsection{(c) 重叠下标数决定两个核的协方差}

\textbf{证明。} 先考虑 $1\le k<m$。因为核对称且样本独立同分布，可以把两组核中的共同观测记为向量 $Z$，各自剩余观测记为 $A,B$。其中 $Z$ 有 $k$ 项，$A,B$ 各有 $m-k$ 项，三组相互独立。

两项中心化核可写成

\[
\widetilde h(Z,A),\qquad \widetilde h(Z,B).
\]

给定 $Z$ 后，$A$ 与 $B$ 仍相互独立，因此两个核条件独立。由条件期望与 (9.1)，

\[
E[\widetilde h(Z,A)\mid Z]
=h_k(Z)-\vartheta,
\]

对另一个核同样成立。故

\[
\begin{aligned}
E[\widetilde h(Z,A)\widetilde h(Z,B)]
&=E\{E[\widetilde h(Z,A)\widetilde h(Z,B)\mid Z]\}\\
&=E\{E[\widetilde h(Z,A)\mid Z]\,
E[\widetilde h(Z,B)\mid Z]\}\\
&=E[(h_k(Z)-\vartheta)^2]\\
&=\operatorname{Var}(h_k(Z))\\
&=\zeta_k.
\end{aligned}
\]

倒数第二步用了 $E h_k(Z)=\vartheta$，这也是塔式法则的直接结果。

若 $k=m$，两组下标完全相同，核对称使两项核相同，此时期望直接等于

\[
E\widetilde h^2=\operatorname{Var}(h)=\zeta_m.
\]

因此对所有题设的 $k$，

\[
\boxed{
E[\widetilde h(X_{i_1},\ldots,X_{i_m})
\widetilde h(X_{j_1},\ldots,X_{j_m})]=\zeta_k.}
\]

补充：若没有共同下标，即 $k=0$，两项核独立且已中心化，协方差为零。

\subsection{(d) 投影方差随保留的变量个数单调不减}

记

\[
H_k=h_k(X_1,\ldots,X_k),\qquad
\mathcal F_k=\sigma(X_1,\ldots,X_k).
\]

由 (a)，$E(H_{k+1}\mid\mathcal F_k)=H_k$。利用全方差公式，

\[
\begin{aligned}
\zeta_{k+1}
&=\operatorname{Var}(H_{k+1})\\
&=\operatorname{Var}(E[H_{k+1}\mid\mathcal F_k])
+E[\operatorname{Var}(H_{k+1}\mid\mathcal F_k)]\\
&=\zeta_k+E[\operatorname{Var}(H_{k+1}\mid\mathcal F_k)]\\
&\ge\zeta_k.
\end{aligned}
\]

所以

\[
\boxed{0\le\zeta_1\le\zeta_2\le\cdots\le\zeta_m.}
\]

 保留的观测越多，平均掉的随机性越少，条件平均本身可以有更大的波动。这里不是说“信息越多，估计误差越大”，而是在比较同一个核的不同条件投影的方差。

\section{第 10 题 三个观测之和为负的 U 统计量}

\textbf{题意。} $n\ge3$，$X_i$ 为独立同分布实值随机变量，核为

\[
h(x_1,x_2,x_3)=\mathbf1_{\{x_1+x_2+x_3<0\}}.
\]

沿用第 9 题记号，写出 $U_n$、$h_1,h_2,h_3$ 和 $\zeta_1,\zeta_2,\zeta_3$。

\subsection{第一步 约定分布函数并处理严格不等号}

设

\[
F(u)=P(X_1\le u),\qquad
F(u-)=P(X_1<u).
\]

记两个独立观测之和的分布函数为

\[
F_2(u)=P(X_1+X_2\le u),
\qquad F_2(u-)=P(X_1+X_2<u).
\]

题目没有假设连续分布，所以阈值处可能有正概率，必须保留左极限。令

\[
\vartheta=P(X_1+X_2+X_3<0).
\]

\subsection{第二步 写出 U 统计量}

核显然对称，且 $0\le h\le1$，故二阶矩条件自动满足。所求统计量为

\[
\boxed{
U_n=\binom n3^{-1}
\sum_{1\le i<j<\ell\le n}
\mathbf1_{\{X_i+X_j+X_\ell<0\}}.}
\]

它是所有三元子样本中“总和为负”的比例，且 $EU_n=\vartheta$。

\subsection{第三步 计算条件投影}

固定第一个观测为 $x$，对其余两个取期望：

\[
\begin{aligned}
h_1(x)
&=P(x+X_2+X_3<0)\\
&=P(X_2+X_3<-x)\\
&=\boxed{F_2((-x)-)}\\
&=\int_{\mathbb R}F((-x-y)-)\,dF(y).
\end{aligned}
\]

这里 $F_2((-x)-)$ 表示在点 $-x$ 处的左极限。

固定前两个观测为 $x_1,x_2$，则

\[
h_2(x_1,x_2)
=P(X_3<-x_1-x_2)
=\boxed{F((-x_1-x_2)-)}.
\]

三个观测都固定时，没有剩余随机变量需要平均，所以

\[
\boxed{h_3(x_1,x_2,x_3)
=\mathbf1_{\{x_1+x_2+x_3<0\}}.}
\]

同时，

\[
\vartheta=\int F_2((-x)-)\,dF(x)
=\iint F((-x-y)-)\,dF(x)\,dF(y).
\]

\subsection{第四步 计算三个投影方差}

由于 $E h_k=\vartheta$，

\[
\zeta_k=E h_k^2-\vartheta^2.
\]

逐一代入得

\[
\boxed{
\zeta_1=\int_{\mathbb R}
[F_2((-x)-)]^2\,dF(x)-\vartheta^2,}
\]

\[
\boxed{
\zeta_2=\iint_{\mathbb R^2}
[F((-x-y)-)]^2\,dF(x)\,dF(y)-\vartheta^2.}
\]

最后 $h_3$ 是概率为 $\vartheta$ 的 0--1 随机变量，故

\[
\boxed{\zeta_3=\vartheta(1-\vartheta).}
\]

原题没有指定共同分布 $F$，因此这些分布积分就是一般情形下的完整答案，不能给出统一数值。若另有连续分布假设，才可以将相应左极限写成通常的分布函数值。

\section{第 11 题 均值与倒数均值乘积的 U 统计量}

\textbf{题意。} $n\ge2$，$X_i$ 独立同分布，并满足

\[
E X_1^2<\infty,\qquad E X_1^{-2}<\infty.
\]

记

\[
\mu=E X_1,\qquad \bar\mu=E(X_1^{-1}).
\]

构造 $\mu\bar\mu$ 的无偏 U 统计量，求精确方差及渐近分布。

为避免将题目的 $\bar\mu$ 与样本均值混淆，下文令

\[
\beta=\bar\mu=E(X_1^{-1}),\qquad
\tau=\mu\beta,\qquad
A=EX_1^2,\quad B=EX_1^{-2}.
\]

逆二阶矩有限意味着 $P(X_1=0)=0$；由 Cauchy--Schwarz，$E|X_1|$ 与 $E|X_1^{-1}|$ 均有限。因此下面涉及的期望都有定义。

\subsection{(a) 构造对称二元核并验证无偏性}

两个独立观测满足

\[
E\frac{X_1}{X_2}
=E X_1\cdot E X_2^{-1}
=\mu\beta.
\]

但 $x_1/x_2$ 不是对称核，所以取对称化

\[
\boxed{h(x,y)=\frac12\left(\frac{x}{y}+\frac{y}{x}\right)}.
\]

于是

\[
Eh(X_1,X_2)
=\frac12(\mu\beta+\mu\beta)=\tau.
\]

对应的二阶 U 统计量为

\[
\begin{aligned}
U_n
&=\binom n2^{-1}\sum_{i<j}
\frac12\left(\frac{X_i}{X_j}+\frac{X_j}{X_i}\right)\\
&=\boxed{\frac1{n(n-1)}\sum_{i\ne j}\frac{X_i}{X_j}}.
\end{aligned}
\]

每项的期望均为 $\tau$，共有 $n(n-1)$ 项，所以

\[
\boxed{EU_n=\tau=\mu\bar\mu.}
\]

为了后续渐近分析，再记

\[
\bar X_n=\frac1n\sum_iX_i,\qquad
\bar Y_n=\frac1n\sum_iX_i^{-1}.
\]

由于

\[
\left(\sum_iX_i\right)\left(\sum_jX_j^{-1}\right)
=\sum_{i\ne j}\frac{X_i}{X_j}+n,
\]

得到重要的等价表达式

\[
\boxed{U_n=\frac{n\bar X_n\bar Y_n-1}{n-1}.}
\tag{11.1}
\]

注意 $\bar Y_n$ 是“倒数的平均”，一般不等于 $1/\bar X_n$。

\subsection{(b) 计算一阶投影方差与核方差}

固定第一个变量为 $x$，有

\[
h_1(x)=E[h(x,X_2)]
=\frac12\left(xE X_2^{-1}+\frac{E X_2}{x}\right)
=\boxed{\frac12\left(\beta x+\frac{\mu}{x}\right)}.
\]

其期望为 $\tau$。平方后，

\[
h_1(X_1)^2
=\frac14\left(\beta^2X_1^2+\mu^2X_1^{-2}+2\mu\beta\right).
\]

因此

\[
E h_1(X_1)^2=\frac14(\beta^2A+\mu^2B+2\tau),
\]

\[
\boxed{
\zeta_1
=\frac14(\beta^2A+\mu^2B+2\tau-4\tau^2).}
\tag{11.2}
\]

再计算核的二阶矩。由 $X_1/X_2$ 与 $X_2/X_1$ 的乘积恒为 1，

\[
h(X_1,X_2)^2
=\frac14\left(\frac{X_1^2}{X_2^2}
+\frac{X_2^2}{X_1^2}+2\right).
\]

独立性给出 $E(X_1^2/X_2^2)=AB$，所以

\[
E h(X_1,X_2)^2=\frac12(AB+1)<\infty.
\]

因此

\[
\boxed{\zeta_2=\frac12(AB+1)-\tau^2.}
\tag{11.3}
\]

\subsection{(c) 通过核项重叠计数求精确方差}

记 $H_{ij}=h(X_i,X_j)$，$i<j$，以及 $N=\binom n2$。则

\[
\operatorname{Var}(U_n)
=\frac1{N^2}
\sum_{i<j}\sum_{k<\ell}\operatorname{Cov}(H_{ij},H_{k\ell}).
\]

对固定的无序下标对 $\{i,j\}$，内层的 $\{k,\ell\}$ 分三类：

\begin{enumerate}
\item 完全相同：只有 1 对，协方差为 $\zeta_2$。
\item 恰有 1 个共同下标：从 $i,j$ 中选共同的一个，再从其余 $n-2$ 个中选另一个，共 $2(n-2)$ 对。由第 9(c) 题，协方差为 $\zeta_1$。
\item 没有共同下标：两项核独立，协方差为 0。

\end{enumerate}
外层共有 $N$ 对，故

\[
\begin{aligned}
\operatorname{Var}(U_n)
&=\frac{N[\zeta_2+2(n-2)\zeta_1]}{N^2}\\
&=\boxed{
\frac{4(n-2)}{n(n-1)}\zeta_1
+\frac2{n(n-1)}\zeta_2}.
\end{aligned}
\tag{11.4}
\]

这里双重求和本身就是有序的“核项对”求和，所以 $2(n-2)$ 不需再除以 2。

代入 (11.2)、(11.3)，也可写为

\[
\boxed{
\operatorname{Var}(U_n)=
\frac{(n-2)(\beta^2A+\mu^2B+2\tau-4\tau^2)
+AB+1-2\tau^2}{n(n-1)}.}
\tag{11.5}
\]

当 $n=2$ 时，(11.4) 给出 $\operatorname{Var}(U_2)=\zeta_2$，与 $U_2=h(X_1,X_2)$ 一致。

\subsection{(d) 一般的一阶渐近分布}

令

\[
V_i=\begin{pmatrix}X_i\\X_i^{-1}\end{pmatrix}.
\]

因为两项都有有限二阶矩，多元中心极限定理给出

\[
\sqrt n
\begin{pmatrix}
\bar X_n-\mu\\
\bar Y_n-\beta
\end{pmatrix}
\xrightarrow{d}
\begin{pmatrix}G_1\\G_2\end{pmatrix}
\sim N_2(0,\Sigma),
\]

其中

\[
\Sigma=
\begin{pmatrix}
A-\mu^2 & 1-\tau\\
1-\tau & B-\beta^2
\end{pmatrix}.
\tag{11.6}
\]

非对角项来自

\[
\operatorname{Cov}(X_1,X_1^{-1})
=E(X_1X_1^{-1})-\mu\beta=1-\tau.
\]

写

\[
D_{X,n}=\bar X_n-\mu,\qquad
D_{Y,n}=\bar Y_n-\beta.
\]

直接展开乘积：

\[
\bar X_n\bar Y_n-\tau
=\beta D_{X,n}+\mu D_{Y,n}+D_{X,n}D_{Y,n}.
\tag{11.7}
\]

由 CLT，$D_{X,n},D_{Y,n}=O_P(n^{-1/2})$，故

\[
\sqrt n D_{X,n}D_{Y,n}=O_P(n^{-1/2})=o_P(1).
\]

由 (11.1)，

\[
\begin{aligned}
\sqrt n(U_n-\tau)
&=\frac n{n-1}\sqrt n(\bar X_n\bar Y_n-\tau)
+\frac{\sqrt n}{n-1}(\tau-1)\\
&=\frac n{n-1}
\left[\beta\sqrt n D_{X,n}+\mu\sqrt n D_{Y,n}\right]+o_P(1).
\end{aligned}
\]

应用连续映射与 Slutsky 定理，

\[
\sqrt n(U_n-\tau)\xrightarrow d\beta G_1+\mu G_2.
\]

极限为均值零的正态变量，其方差为

\[
\begin{aligned}
V
&=(\beta,\mu)\Sigma(\beta,\mu)^\top\\
&=\beta^2(A-\mu^2)+\mu^2(B-\beta^2)+2\mu\beta(1-\tau)\\
&=\beta^2A+\mu^2B+2\tau-4\tau^2\\
&=4\zeta_1.
\end{aligned}
\]

因此一般成立

\[
\boxed{
\sqrt n(U_n-\mu\bar\mu)
\xrightarrow d
N\left(0,\,
\bar\mu^2EX_1^2+\mu^2EX_1^{-2}
+2\mu\bar\mu-4\mu^2\bar\mu^2\right).}
\tag{11.8}
\]

这里允许方差为零；$N(0,0)$ 表示集中在零点的退化分布。若 $V>0$，还可写成通常的标准化形式

\[
\boxed{\frac{\sqrt n(U_n-\tau)}{\sqrt V}\xrightarrow dN(0,1).}
\]

\subsection{(e) 一阶退化时的极限}

题设未假定 $\zeta_1>0$，故需区分一阶退化情形。若 $\zeta_1=0$，(11.8) 仍正确，但只告诉我们 $\sqrt n(U_n-\tau)\to_P0$。下面给出此时 $n$ 尺度的极限。

由 $4\zeta_1=0$，均值为零的随机变量

\[
W=\beta(X_1-\mu)+\mu(X_1^{-1}-\beta)
\]

方差为零，故

\[
W=0\quad\text{a.s.}
\]

于是对每个 $n$，

\[
\beta D_{X,n}+\mu D_{Y,n}=0\quad\text{a.s.}
\]

代入 (11.7)，得到精确等式

\[
\bar X_n\bar Y_n-\tau=D_{X,n}D_{Y,n}.
\]

再由 (11.1)，

\[
n(U_n-\tau)
=\frac n{n-1}
\left[(\sqrt n D_{X,n})(\sqrt n D_{Y,n})+\tau-1\right].
\]

使用联合 CLT 与连续映射定理，得到

\[
\boxed{n(U_n-\tau)\xrightarrow dG_1G_2+\tau-1,\qquad \zeta_1=0,}
\tag{11.9}
\]

其中 $(G_1,G_2)$ 的协方差仍由 (11.6) 给出。这个极限一般不是正态；若模型进一步退化，也可能恒为零。

\textbf{具体例子。} 若 $P(X_1=1)=P(X_1=-1)=1/2$，则

\[
\mu=\beta=0,\quad A=B=1,\quad X_1^{-1}=X_1.
\]

所以 $\tau=0,\zeta_1=0$，并且 $G_1=G_2=Z$，$Z\sim N(0,1)$。此时

\[
U_n=\frac{n\bar X_n^2-1}{n-1},
\qquad
\boxed{nU_n\xrightarrow d Z^2-1=\chi_1^2-1.}
\]

这说明题设确实允许一阶退化，补充分支不是多余假设。

\end{document}
